% Lesson 15 — Vocabulary: Words and expressions related to comparing prices or quality to determine the best buys for consumer economy and community resources/services. # Semaine / Période Chapitre Titre Date effective Visa de l'enseignant — Anglais Tle A — Cours signé OnBuch+
% Programme officiel MINESEC. Compiler avec : tectonic EN15-vocabulary-words-and-expressions-related-to-comparing-p.tex
\newcommand{\DOCMATIERE}{Anglais}
\newcommand{\DOCNIVEAU}{Tle A}
\newcommand{\DOCMODULE}{Chapter 2 : Economic Life and Occupations}
\newcommand{\DOCLECON}{EN15}
\newcommand{\DOCTITRE}{Lesson 15 — Vocabulary: Words and expressions related to comparing prices or quality to determine the best buys for consumer economy and community resources/services. \# Semaine / Période Chapitre Titre Date effective Visa de l'enseignant}
% OnBuch+ — préambule « cours signés » ENRICHI (compilé avec tectonic / XeLaTeX)
% Système visuel « Pop » d'OnBuch+ : fond crème (jamais blanc), bordures encre
% épaisses, ombres « dures » décalées, titres Archivo Black en capitales.
% Version MATHÉMATIQUES : graphes (pgfplots/tikz), boîtes pédagogiques
% (définition, propriété, méthode, attention, le savais-tu, exercice résolu,
% exercice, TP, bilan), page de garde et signature OnBuch+ ancrée en bas.
%
% Macros attendues dans main.tex AVANT \input{preamble} :
%   \DOCMATIERE  \DOCNIVEAU  \DOCMODULE  \DOCLECON (numéro)  \DOCTITRE
\documentclass[11pt]{article}

\usepackage{fontspec}
\usepackage[english,french]{babel} % français = langue principale ; l'anglais via \eng{..} / textebox
\usepackage[a4paper,margin=2cm,top=3.4cm,bottom=2.8cm,headheight=1.4cm,footskip=1.3cm]{geometry}
\usepackage{amsmath,amssymb,amsfonts}
\usepackage{mathtools} % \xmapsto, \coloneqq... souvent utilisés par les modèles
\usepackage{cancel} % simplifications barrées (bilans de forces)
% \abs{z} (module, valeur absolue) et \norm{u} (norme), utilisés par les modèles
\providecommand{\abs}{}\let\abs\relax\DeclarePairedDelimiter{\abs}{\lvert}{\rvert}
\providecommand{\norm}{}\let\norm\relax\DeclarePairedDelimiter{\norm}{\lVert}{\rVert}
\usepackage[table]{xcolor} % [table] : fournit aussi \rowcolors (alternance de lignes)
\usepackage{pagecolor}
\usepackage{tikz}
\usetikzlibrary{arrows.meta,positioning,calc,shapes.geometric,shapes.misc,shapes.arrows,decorations.pathmorphing,decorations.pathreplacing,decorations.markings,patterns,shadows,mindmap,trees,fit,backgrounds,matrix,intersections,angles,quotes}
\usepackage{pgfplots}
\usepackage[version=4]{mhchem} % équations nucléaires \ce{^{235}_{92}U}
\usepackage[european,siunitx]{circuitikz} % schémas électriques
\pgfplotsset{compat=1.18}
\usepgfplotslibrary{fillbetween,groupplots} % aires entre courbes (calcul intégral)
\usepackage{enumitem}
\usepackage{fancyhdr}
% [most] charge toutes les bibliothèques tcolorbox (listings, minted, raster,
% documentation...) et gonfle inutilement la mémoire TeX sur un document long
% (~300 boîtes) : on ne charge que ce qui est réellement utilisé.
\usepackage[skins,breakable]{tcolorbox}
\usepackage{graphicx}
\usepackage{colortbl}
\usepackage{booktabs}
\usepackage{tabularx}
\usepackage{diagbox} % case d'en-tête diagonale des tableaux à double entrée
% Colonnes extensibles centrée / gauche / droite (C, L, R), souvent utilisées par les modèles
\newcolumntype{C}{>{\centering\arraybackslash}X}
\newcolumntype{L}{>{\raggedright\arraybackslash}X}
\newcolumntype{R}{>{\raggedleft\arraybackslash}X}
\usepackage{array}
\usepackage{multicol}
\usepackage{multirow}
\usepackage{siunitx}
% parse-numbers=false : accepte aussi \SI{2,5 \times 10^{-3}}{...}
\sisetup{locale=FR,output-decimal-marker={,},group-minimum-digits=4,parse-numbers=false}
\DeclareSIUnit\ohm{\ensuremath{\symup{\Omega}}} % U+2126 absent de la police
\DeclareSIUnit\division{div} % réglages d'oscilloscope (ms/div, V/div)
\DeclareSIUnit\tr{tr} % tours (vitesse de rotation, ex. \SI{1500}{\tr\per\minute})
\DeclareSIUnit\molar{M} % concentration molaire (ex. \SI{3}{\molar}, \SI{1}{\micro\molar})
\DeclareSIUnit\an{an} % année (le modèle confond parfois avec \year, un compteur TeX) — ex. \SI{5}{\kilo\meter\per\an}
\DeclareSIUnit\FCFA{FCFA} % franc CFA (ex. \SI{500}{\FCFA\per\kilo\gram})
\DeclareSIUnit\degreeBrix{\textdegree Brix} % teneur en sucre d'un sirop/jus (réfractométrie)
\DeclareSIUnit\month{mois} % le modèle utilise parfois \month, un compteur TeX, comme unité de durée
\usepackage{qrcode}
% \degree : commande fréquemment utilisée par les modèles pour un angle ou une
% température en dehors de \SI{}{} (non fournie par les paquets chargés).
\providecommand{\degree}{^{\circ}}
% \qed (fin de démonstration), utilisé par les modèles sans amsthm
\providecommand{\qed}{\hfill$\square$}
% \barre : double barre des valeurs interdites dans un tableau de variations (tkz-tab)
\providecommand{\barre}{\ensuremath{\|}}
% environnement proof (amsthm non chargé) utilisé par les modèles
\ifdefined\proof\else\newenvironment{proof}[1][Démonstration]{\par\noindent\textit{#1.}\ }{\qed\par}\fi
\usepackage{lastpage}
\usepackage{hyperref}
\hypersetup{hidelinks}

% ============================================================
% Palette « Pop » OnBuch+ — mêmes hex que la classe P (plus_ui.dart)
% ============================================================
\definecolor{poporange}{HTML}{F59321}
\definecolor{popdark}{HTML}{E07A0C}
\definecolor{popcream}{HTML}{FFF6E8}
\definecolor{popink}{HTML}{1C1714}
\definecolor{poppurple}{HTML}{7A5AE0}
\definecolor{popgreen}{HTML}{2E9E6A}
\definecolor{poppink}{HTML}{E0457B}
\definecolor{popblue}{HTML}{2F7DE1}
\definecolor{popgold}{HTML}{C98A12}
\definecolor{popmuted}{HTML}{8A7F76}
\definecolor{popline}{HTML}{EADFD2}
\definecolor{popgray}{HTML}{8A7F76}
\colorlet{popgrey}{popgray}
% Alias tolérants (le modèle nomme parfois les couleurs différemment)
\colorlet{popwhite}{white}
\colorlet{popblack}{popink}
\colorlet{popred}{poppink}
\colorlet{popyellow}{popgold}
% Teintes claires (fonds de tableaux/encadrés) — le modèle nomme parfois
% les couleurs avec un suffixe L (ex. popblueL) sans les avoir définies.
\colorlet{poporangeL}{poporange!20}
\colorlet{popdarkL}{popdark!20}
\colorlet{poppurpleL}{poppurple!20}
\colorlet{popgreenL}{popgreen!20}
\colorlet{poppinkL}{poppink!20}
\colorlet{popblueL}{popblue!20}
\colorlet{popgoldL}{popgold!20}
\colorlet{popredL}{popred!20}
\colorlet{popgrayL}{popgray!20}
% Teintes claires pour fonds de tableaux / zones
\colorlet{popblueL}{popblue!10!white}
\colorlet{poporangeL}{poporange!14!white}
\colorlet{popgreenL}{popgreen!12!white}
\colorlet{poppurpleL}{poppurple!12!white}
\colorlet{poppinkL}{poppink!12!white}
\colorlet{popgoldL}{popgold!14!white}

\pagecolor{popcream}

% ============================================================
% Typographie
% ============================================================
\defaultfontfeatures{Ligatures=TeX}
\setmainfont{PlusJakartaSans-Medium.ttf}[Path=../../../fonts/,BoldFont=PlusJakartaSans-Bold.ttf]
% Symboles, API et emojis absents de la police principale : repli sur DejaVu Sans (caractères actifs)
\newfontface\symfont{DejaVuSans.ttf}[Path=../../../fonts/]
\makeatletter
\newcommand\symactive[1]{\catcode#1=13 \begingroup\lccode`\~=#1 \lowercase{\endgroup\def~}{{\symfont\char#1}}}
\newcommand\symblank[1]{\catcode#1=13 \begingroup\lccode`\~=#1 \lowercase{\endgroup\def~}{}}
\newcommand\symhyph[1]{\catcode#1=13 \begingroup\lccode`\~=#1 \lowercase{\endgroup\def~}{\mbox{-}}}
\newcount\symc
\newcommand\symrange[2]{\symc=#1 \@whilenum\symc<#2 \do{\expandafter\symactive\expandafter{\the\symc}\advance\symc by 1 }}
\newcommand\symblankrange[2]{\symc=#1 \@whilenum\symc<#2 \do{\expandafter\symblank\expandafter{\the\symc}\advance\symc by 1 }}
\makeatother
\symrange{"0250}{"02FF}\symactive{"2713}\symactive{"2714}\symactive{"2717}\symactive{"2718}\symactive{"2610}\symactive{"2611}\symactive{"2612}
\symhyph{"2011}\symhyph{"2E1A}\symblankrange{"1F300}{"1FAFF}\symblank{"FE0F}
% Maths en sans-serif (Fira Math) avec chiffres et lettres droites de Plus
% Jakarta, pour que les formules se fondent dans le texte.
\usepackage[math-style=ISO,bold-style=ISO]{unicode-math}
\setmathfont{FiraMath-Regular.otf}[Path=../../../fonts/]
\setmathfont{PlusJakartaSans-Medium.ttf}[Path=../../../fonts/,range={up/num,up/{Latin,latin},\mathup}]
\usepackage{icomma} % virgule décimale sans espace en mode math
% Caractères mathématiques Unicode absents des polices de texte (Plus Jakarta,
% Archivo Black) : rendus par leur équivalent mathématique, en texte comme en
% formule — sinon ils s'affichent en carré vide (ex. « Arithmétique dans ℤ »).
\usepackage{newunicodechar}
\newunicodechar{ℝ}{\ensuremath{\mathbb{R}}}
\newunicodechar{ℤ}{\ensuremath{\mathbb{Z}}}
\newunicodechar{ℂ}{\ensuremath{\mathbb{C}}}
\newunicodechar{ℕ}{\ensuremath{\mathbb{N}}}
\newunicodechar{ℚ}{\ensuremath{\mathbb{Q}}}
\newunicodechar{𝔻}{\ensuremath{\mathbb{D}}}
\newunicodechar{α}{\ensuremath{\alpha}}
\newunicodechar{β}{\ensuremath{\beta}}
\newunicodechar{γ}{\ensuremath{\gamma}}
\newunicodechar{δ}{\ensuremath{\delta}}
\newunicodechar{ε}{\ensuremath{\varepsilon}}
\newunicodechar{θ}{\ensuremath{\theta}}
\newunicodechar{λ}{\ensuremath{\lambda}}
\newunicodechar{μ}{\ensuremath{\mu}}
\newunicodechar{σ}{\ensuremath{\sigma}}
\newunicodechar{τ}{\ensuremath{\tau}}
\newunicodechar{φ}{\ensuremath{\varphi}}
\newunicodechar{ψ}{\ensuremath{\psi}}
\newunicodechar{ω}{\ensuremath{\omega}}
\newunicodechar{Σ}{\ensuremath{\Sigma}}
% majuscules produites par \MakeUppercase dans les titres (α-aminés -> Α)
\newunicodechar{Α}{\ensuremath{\alpha}}
\newunicodechar{Β}{\ensuremath{\beta}}
\newunicodechar{Γ}{\ensuremath{\Gamma}}
\newunicodechar{Δ}{\ensuremath{\Delta}}
\newunicodechar{Ε}{\ensuremath{\varepsilon}}
\newunicodechar{Θ}{\ensuremath{\Theta}}
\newunicodechar{Λ}{\ensuremath{\Lambda}}
\newunicodechar{Μ}{\ensuremath{\mu}}
\newunicodechar{Τ}{\ensuremath{\tau}}
\newunicodechar{Φ}{\ensuremath{\Phi}}
\newunicodechar{Ψ}{\ensuremath{\Psi}}
\newunicodechar{Ω}{\ensuremath{\Omega}}
\newunicodechar{↦}{\ensuremath{\mapsto}}
\newunicodechar{∈}{\ensuremath{\in}}
\newunicodechar{∉}{\ensuremath{\notin}}
\newunicodechar{∧}{\ensuremath{\wedge}}
\newunicodechar{⇔}{\ensuremath{\Leftrightarrow}}
\newunicodechar{⟺}{\ensuremath{\Longleftrightarrow}}
\newunicodechar{⇒}{\ensuremath{\Rightarrow}}
\newunicodechar{⟹}{\ensuremath{\Longrightarrow}}
\newunicodechar{∘}{\ensuremath{\circ}}
\newunicodechar{∪}{\ensuremath{\cup}}
\newunicodechar{∩}{\ensuremath{\cap}}
\newunicodechar{≡}{\ensuremath{\equiv}}
\newunicodechar{⊂}{\ensuremath{\subset}}
\newunicodechar{⊥}{\ensuremath{\perp}}
\newunicodechar{∃}{\ensuremath{\exists}}
\newunicodechar{∀}{\ensuremath{\forall}}
\newunicodechar{∛}{\ensuremath{\sqrt[3]{\,}}}
\newunicodechar{∜}{\ensuremath{\sqrt[4]{\,}}}
\newunicodechar{⃗}{\ensuremath{{}^{\to}}}
\newunicodechar{ⁿ}{\ifmmode^{n}\else\textsuperscript{\ensuremath{n}}\fi}
\newunicodechar{ˣ}{\ifmmode^{x}\else\textsuperscript{\ensuremath{x}}\fi}
\newunicodechar{ᵇ}{\ifmmode^{b}\else\textsuperscript{\ensuremath{b}}\fi}
\newunicodechar{ʳ}{\ifmmode^{r}\else\textsuperscript{\ensuremath{r}}\fi}
\newunicodechar{ᵘ}{\ifmmode^{u}\else\textsuperscript{\ensuremath{u}}\fi}
\newunicodechar{ʸ}{\ifmmode^{y}\else\textsuperscript{\ensuremath{y}}\fi}
\newunicodechar{ᵃ}{\ifmmode^{a}\else\textsuperscript{\ensuremath{a}}\fi}
\newunicodechar{ᵉ}{\ifmmode^{e}\else\textsuperscript{\ensuremath{e}}\fi}
\newunicodechar{ᵏ}{\ifmmode^{k}\else\textsuperscript{\ensuremath{k}}\fi}
\newunicodechar{ᵖ}{\ifmmode^{p}\else\textsuperscript{\ensuremath{p}}\fi}
\newunicodechar{ᴺ}{\ifmmode^{N}\else\textsuperscript{\ensuremath{N}}\fi}
\newunicodechar{ᶜ}{\ifmmode^{c}\else\textsuperscript{\ensuremath{c}}\fi}
\newunicodechar{ʰ}{\ifmmode^{h}\else\textsuperscript{\ensuremath{h}}\fi}
\newunicodechar{ᵠ}{\ifmmode^{q}\else\textsuperscript{\ensuremath{q}}\fi}
\newunicodechar{ₙ}{\ifmmode_{n}\else\textsubscript{\ensuremath{n}}\fi}
\newunicodechar{ₐ}{\ifmmode_{a}\else\textsubscript{\ensuremath{a}}\fi}
\newunicodechar{ᵢ}{\ifmmode_{i}\else\textsubscript{\ensuremath{i}}\fi}
\newunicodechar{ᵣ}{\ifmmode_{r}\else\textsubscript{\ensuremath{r}}\fi}
\newunicodechar{ₖ}{\ifmmode_{k}\else\textsubscript{\ensuremath{k}}\fi}
\newunicodechar{ᵦ}{\ifmmode_{\beta}\else\textsubscript{\ensuremath{\beta}}\fi}
\newunicodechar{ₓ}{\ifmmode_{x}\else\textsubscript{\ensuremath{x}}\fi}
\newfontfamily\popdisplay{ArchivoBlack-Regular.ttf}[Path=../../../fonts/]
\newfontfamily\poplogo{SpaceGrotesk-Bold.ttf}[Path=../../../fonts/]
\newfontfamily\popsemi{PlusJakartaSans-SemiBold.ttf}[Path=../../../fonts/]
\newfontfamily\popxbold{PlusJakartaSans-ExtraBold.ttf}[Path=../../../fonts/]
\newfontfamily\popxboldls{PlusJakartaSans-ExtraBold.ttf}[Path=../../../fonts/,LetterSpace=3.0]

\setlist[enumerate]{leftmargin=1.7em,itemsep=5pt,topsep=4pt}
\setlist[itemize]{leftmargin=1.7em,itemsep=4pt,topsep=4pt,label={\color{poporange}\textbullet}}
\setlength{\parindent}{0pt}
\setlength{\parskip}{5pt}
\renewcommand{\arraystretch}{1.25}

% pgfplots : style Pop par défaut
\pgfplotsset{
  every axis/.append style={axis line style={popink,line width=1pt},tick style={popink},
    grid style={popline,line width=0.5pt},label style={font=\small},tick label style={font=\footnotesize}},
  popcurve/.style={poporange,line width=1.8pt,no markers,smooth},
}
% Même style disponible dans un \draw TikZ simple (hors pgfplots)
\tikzset{popcurve/.style={poporange,line width=1.8pt}}
\tikzset{popgrid/.style={popline,very thin}}
\colorlet{popcurve}{poporange} % le nom du style est parfois utilisé comme couleur

% ============================================================
% Pill / squiggle
% ============================================================
\newcommand{\poppill}[3]{%
  \tikz[baseline=-0.55ex]\node[draw=popink,line width=1.2pt,rounded corners=2.6pt,%
    fill=#2,inner xsep=6.5pt,inner ysep=3pt]{%
    \popxboldls\fontsize{7}{7}\selectfont\color{#3}\MakeUppercase{#1}};%
}
\newcommand{\popsquiggle}[1]{%
  \tikz[baseline=-0.4ex]{
    \draw[#1,line width=2.6pt,line cap=round]
      plot [smooth] coordinates {(0,0.09) (0.34,0.01) (0.68,0.21) (1.02,0.01) (1.36,0.10)};
  }%
}
% Mot-clé surligné (marqueur orange sous le texte)
\newcommand{\cle}[1]{{\popxbold\color{popdark}#1}}

% ============================================================
% En-tête / pied
% ============================================================
\pagestyle{fancy}
\fancyhf{}
\renewcommand{\headrulewidth}{0pt}
\renewcommand{\footrulewidth}{0pt}
\lhead{\raisebox{-0.15ex}{\poplogo\fontsize{13}{13}\selectfont\color{popink}OnBuch\color{poporange}+}}
\rhead{\poppill{\DOCMATIERE\ \textbullet\ \DOCNIVEAU\ \textbullet\ Leçon \DOCLECON}{white}{popink}}
\lfoot{\popxbold\fontsize{7.6}{7.6}\selectfont\color{popmuted}\MakeUppercase{Cours signé OnBuch+}\ \textbullet\ {\popsemi\DOCTITRE}}
\rfoot{\tikz[baseline=-0.6ex]\node[rounded corners=8.5pt,fill=popink,minimum height=17pt,minimum width=17pt,inner xsep=5pt,inner sep=0]{\popxbold\fontsize{8}{8}\selectfont\color{popcream}\thepage\,/\,\pageref*{LastPage}};}

% ============================================================
% Titres
% ============================================================
\newcommand{\chaptitre}[2]{%
  \begin{tcolorbox}[enhanced,colback=poporange,colframe=popink,boxrule=2.6pt,arc=11pt,
      drop shadow=popink,left=15pt,right=15pt,top=12pt,bottom=11pt,before skip=3pt,after skip=18pt]
    \poppill{#2}{popink}{popcream}\par\vspace{9pt}
    {\raggedright\hyphenpenalty=10000\exhyphenpenalty=10000\popdisplay\fontsize{22}{24}\selectfont\color{popcream}\MakeUppercase{#1}\par}
  \end{tcolorbox}
}
% Section numérotée : pastille orange avec le numéro + titre Archivo
\newcounter{popsec}
\newcommand{\coursec}[1]{%
  \stepcounter{popsec}\setcounter{popsub}{0}%
  \par\vspace{16pt}\noindent
  \begin{minipage}{\linewidth}
  \tikz[baseline=-0.7ex]\node[draw=popink,line width=1.6pt,fill=poporange,rounded corners=4pt,minimum size=22pt,inner sep=0]{\popdisplay\fontsize{12}{12}\selectfont\color{popcream}\thepopsec};%
  \hspace{8pt}{\popdisplay\fontsize{15}{17}\selectfont\color{popink}\MakeUppercase{#1}}\par
  \vspace{4pt}\noindent{\color{poporange}\rule{36pt}{3.6pt}}
  \end{minipage}\par\nopagebreak\vspace{8pt}
}
\newcounter{popsub}
\newcommand{\courssub}[1]{%
  \stepcounter{popsub}%
  \par\vspace{9pt}\noindent{\popxbold\fontsize{11.5}{13}\selectfont\color{popink}{\color{poporange}\thepopsec.\thepopsub}\hspace{6pt}#1}\par\nopagebreak\vspace{4pt}
}

% ============================================================
% Boîtes pédagogiques — même recette : fond blanc, bordure épaisse colorée,
% ombre dure encre, badge plein. Titre optionnel [#1].
% ============================================================
\tcbset{popbase/.style={breakable,enhanced,colback=white,boxrule=2.3pt,arc=9pt,
  shadow={3pt}{-3pt}{0pt}{popink},left=14pt,right=14pt,top=11pt,bottom=11pt,before skip=11pt,after skip=15pt,
  fontupper=\popsemi\fontsize{10.6}{14.5}\selectfont\color{popink},
  fontlower=\popsemi\fontsize{10.6}{14.5}\selectfont\color{popink}}}
% Coche dessinée (le glyphe ✓ n'existe pas dans les polices du document)
\newcommand{\popcheck}[1]{\tikz[baseline=-0.5ex]\draw[#1,line width=1.6pt,line cap=round,line join=round](-0.13,0.0)--(-0.04,-0.09)--(0.13,0.1);}
\providecommand{\coche}{\popcheck{popgreen}}
\providecommand{\resultat}[1]{\par\smallskip\noindent\fcolorbox{popgreen}{popgreenL}{\parbox{\dimexpr\linewidth-2\fboxsep-2\fboxrule\relax}{\textbf{Résultat.} #1}}\par\smallskip}
\newcommand{\popboxhead}[3]{\poppill{#1}{#2}{white}\ifx\relax#3\relax\else\hspace{6pt}{\popxbold\fontsize{10.6}{12}\selectfont\color{popink}#3}\fi\par\vspace{7pt}}

\newtcolorbox{aretenir}[1][]{popbase,colframe=popgreen,before upper={\popboxhead{À retenir}{popgreen}{#1}}}
% Texte en anglais (document de lecture, dialogue, extrait) : cadre
% d'étude. Un titre optionnel (source, auteur) s'affiche dans l'en-tête.
\newtcolorbox{textebox}[1][]{popbase,colframe=popblue,before upper={\popboxhead{Text}{popblue}{#1}\begin{otherlanguage}{english}},after upper={\end{otherlanguage}}}
% Marques ✓ / ✗ (forme correcte / fautive) souvent utilisées par les modèles
\providecommand{\cross}{\textcolor{poppink}{\ensuremath{\times}}}
\providecommand{\xmark}{\cross}\providecommand{\crossmark}{\cross}\providecommand{\cmark}{\ensuremath{\checkmark}}
% Ligne de réponse / justification à compléter (inventée par les modèles)
\providecommand{\justification}[1]{\par\noindent\textit{Justification~:} \dotfill\par}
% \textipa (package tipa non chargé) : la police ne couvre pas l'API -> simple passage
\providecommand{\textipa}[1]{#1}
% Soulignement (ulem non chargé) : \uline, \ul
\providecommand{\uline}[1]{\underline{#1}}\providecommand{\ul}[1]{\underline{#1}}
% \glossaire{\item ...} inventée par les modèles : liste de glossaire
\providecommand{\glossaire}[1]{\begin{itemize}#1\end{itemize}}
% \alert (beamer) inventée par les modèles : texte mis en évidence
\providecommand{\alert}[1]{\textbf{\textcolor{poppink}{#1}}}
% variantes ASCII d'ordinaux écrites en macro
\providecommand{\iere}{\textsuperscript{re}}\providecommand{\ieme}{\textsuperscript{e}}\providecommand{\ere}{\textsuperscript{re}}\providecommand{\eme}{\textsuperscript{e}}\providecommand{\er}{\textsuperscript{er}}\providecommand{\ie}{\textsuperscript{e}}
% Ordinaux français écrits en macro par les modèles : 1\ière, 3\ième
\newcommand{\ière}{\textsuperscript{re}}\newcommand{\ième}{\textsuperscript{e}}
% \exemple (inventée par les modèles) : étiquette « Exemple : »
\providecommand{\exemple}{\textbf{Exemple~:}\ }
% \square utilisable aussi hors mode math (cases à cocher)
\AtBeginDocument{\let\squareMath\square\renewcommand{\square}{\ensuremath{\squareMath}}}
% Mot ou phrase en anglais à l'intérieur d'un paragraphe français
\newcommand{\eng}[1]{\ifmmode\text{\textit{#1}}\else\begingroup\selectlanguage{english}\itshape #1\endgroup\fi}
\let\esp\eng % alias toléré
% flèches d'intonation inventées par les modèles
\providecommand{\textsearrow}{}\renewcommand{\textsearrow}{\ensuremath{\searrow}}\providecommand{\textnearrow}{}\renewcommand{\textnearrow}{\ensuremath{\nearrow}}
% \fr{..} : mot français (inventé par les modèles)
\providecommand{\fr}[1]{\textit{#1}}
\newtcolorbox{exemplebox}[1][]{popbase,colframe=poporange,before upper={\popboxhead{Exemple}{poporange}{#1}}}
\newtcolorbox{definition}[1][]{popbase,colframe=popblue,before upper={\popboxhead{Définition}{popblue}{#1}}}
\newtcolorbox{propriete}[1][]{popbase,colframe=poppurple,before upper={\popboxhead{Propriété}{poppurple}{#1}}}
\newtcolorbox{methode}[1][]{popbase,colframe=popgold,before upper={\popboxhead{Méthode}{popgold}{#1}}}
\newtcolorbox{attention}[1][]{popbase,colframe=poppink,before upper={\popboxhead{Attention}{poppink}{#1}}}
\newtcolorbox{experience}[1][]{popbase,colframe=popblue,colback=popblueL,before upper={\popboxhead{Expérience}{popblue}{#1}}}
% « Le savais-tu ? » : carte violette pleine (contexte, culture, Cameroun)
\newtcolorbox{savaistu}[1][]{popbase,colback=poppurple,colframe=popink,
  fontupper=\popsemi\fontsize{10.4}{14.2}\selectfont\color{white},
  before upper={\poppill{Le savais-tu\,?}{white}{poppurple}\ifx\relax#1\relax\else\hspace{6pt}{\popxbold\color{white}#1}\fi\par\vspace{7pt}}}
% Exercice résolu : énoncé en haut, \tcblower puis la solution sur fond crème
\newcounter{popexo}\newcounter{popexr}
\newtcolorbox{exoresolu}[1][]{popbase,colframe=poporange,colbacklower=popcream,
  segmentation style={popink,line width=1.2pt,dashed},
  before upper={\stepcounter{popexr}\popboxhead{Exercice résolu \thepopexr}{poporange}{#1}},
  before lower={\poppill{Solution}{popink}{popcream}\par\vspace{6pt}}}
% Exercice (non résolu en place) — la correction est dans la partie Corrigés
\newtcolorbox{exercice}[1][]{popbase,colframe=popink,
  before upper={\stepcounter{popexo}\popboxhead{Exercice \thepopexo}{popink}{#1}}}
% Corrigé
\newtcolorbox{corrige}[1][]{popbase,colframe=popgreen,colback=popgreenL,
  before upper={\popboxhead{Corrigé}{popgreen}{#1}}}
% Objectifs / prérequis / bilan
\newtcolorbox{objectifs}{popbase,colframe=popink,colback=poporangeL,
  before upper={\popboxhead{Objectifs de la leçon}{popink}{}}}
\newtcolorbox{prerequis}{popbase,colframe=popink,colback=popblueL,
  before upper={\popboxhead{Prérequis}{popblue}{}}}
\newtcolorbox{bilan}{popbase,colframe=popink,colback=white,boxrule=2.8pt,
  before upper={\popboxhead{Fiche bilan}{poporange}{L'essentiel à connaître pour l'examen}}}
% Figure encadrée avec légende
\newtcolorbox{popfigure}[1][]{enhanced,breakable=false,colback=white,colframe=popline,boxrule=1.4pt,arc=8pt,
  left=10pt,right=10pt,top=10pt,bottom=8pt,before skip=10pt,after skip=12pt,halign=center,
  lower separated=false,halign lower=center,
  fontlower=\popsemi\fontsize{9}{11.5}\selectfont\color{popmuted},
  #1}
\newcommand{\legende}[1]{\par\vspace{6pt}{\popsemi\fontsize{9}{11.5}\selectfont\color{popmuted}#1}}

% ============================================================
% Signature OnBuch+
% ============================================================
\newcommand{\poplogoinline}[1]{{\poplogo\fontsize{#1}{#1}\selectfont\color{popink}OnBuch\color{poporange}+}}
\newcommand{\signatureonbuch}{%
  \begin{tcolorbox}[enhanced,colback=popink,colframe=popink,boxrule=2.2pt,arc=10pt,
      drop shadow=poporange,left=14pt,right=14pt,top=10pt,bottom=10pt,before skip=0pt,after skip=0pt]
    \tikz[baseline=-0.7ex]\node[circle,fill=popgreen,draw=popcream,line width=1.3pt,minimum size=22pt,inner sep=0]{\popcheck{white}};%
    \hspace{9pt}\begin{minipage}[c]{0.62\linewidth}
      {\popxbold\fontsize{12}{13}\selectfont\color{popcream}Signé\ \textbullet\ {\poplogo OnBuch\color{poporange}+}}\par\vspace{2pt}
      {\raggedright\popsemi\fontsize{8}{10.5}\selectfont\color{popline}\MakeUppercase{Équipe pédagogique OnBuch+}\\\MakeUppercase{Conforme au programme officiel MINESEC}\par}
    \end{minipage}\hfill
    {\poplogo\fontsize{22}{22}\selectfont\color{popcream}OnBuch\color{poporange}+}
  \end{tcolorbox}%
}

% ============================================================
% Page de garde
% #1 titre · #2 matière · #3 niveau · #4 module · #5 n° leçon · #6 durée ·
% #7 description / objectifs (texte ou itemize)
% ============================================================
\newcommand{\pagegarde}[7]{%
  \thispagestyle{empty}
  \newgeometry{a4paper,margin=1.7cm,top=1.6cm,bottom=1.4cm}%
  \begin{tikzpicture}[remember picture,overlay]
    \fill[poporange!18] ([xshift=-3.2cm,yshift=-3.6cm]current page.north east) circle (4.6cm);
    \fill[poppurple!14] ([xshift=2.4cm,yshift=7.6cm]current page.south west) circle (3.8cm);
  \end{tikzpicture}%
  {\poplogo\fontsize{30}{30}\selectfont\color{popink}OnBuch\color{poporange}+}\hfill
  \raisebox{0.6ex}{\poppill{Cours signé \textbullet\ Premium}{popink}{popcream}}\par
  \vspace{1pt}\popsquiggle{poppurple}\par
  \vspace{8pt}
  \begin{tcolorbox}[enhanced,colback=poporange,colframe=popink,boxrule=2.8pt,arc=13pt,
      drop shadow=popink,left=18pt,right=18pt,top=12pt,bottom=13pt,after skip=10pt]
    \poppill{#2 \textbullet\ #3}{popink}{popcream}\hspace{5pt}\poppill{#4}{white}{popink}\par\vspace{8pt}
    {\popdisplay\fontsize{14}{14}\selectfont\color{popink}LEÇON #5}\par\vspace{4pt}
    {\raggedright\hyphenpenalty=10000\exhyphenpenalty=10000\popdisplay\fontsize{25}{27}\selectfont\color{popcream}\MakeUppercase{#1}\par}
    \vspace{8pt}
    {\popxbold\fontsize{9.5}{9.5}\selectfont\color{popink}\MakeUppercase{Durée conseillée : #6}}
  \end{tcolorbox}
  \begin{tcolorbox}[enhanced,colback=white,colframe=popink,boxrule=2.2pt,arc=10pt,
      drop shadow=popink,left=16pt,right=16pt,top=10pt,bottom=10pt,after skip=8pt]
    \poppill{Dans cette leçon}{poppurple}{white}\par\vspace{5pt}
    \popsemi\fontsize{9.3}{12.6}\selectfont\color{popink}#7
  \end{tcolorbox}
  \noindent\begin{minipage}[t]{0.487\linewidth}
    \begin{tcolorbox}[enhanced,colback=popgreen,colframe=popink,boxrule=2.2pt,arc=10pt,
        drop shadow=popink,left=11pt,right=11pt,top=10pt,bottom=10pt,height=3.1cm,valign=center]
      \tcbox[colback=white,colframe=popink,boxrule=1.3pt,arc=5pt,boxsep=0pt,nobeforeafter,
        left=3pt,right=3pt,top=3pt,bottom=3pt,valign=center]{\qrcode[height=1.9cm]{https://play.google.com/store/apps/details?id=com.luvvix.onbuch&hl=fr}}%
      \hspace{8pt}\begin{minipage}[c]{0.5\linewidth}
        {\popdisplay\fontsize{11}{12.5}\selectfont\color{white}\MakeUppercase{Télécharge OnBuch}}\par\vspace{4pt}
        {\raggedright\popsemi\fontsize{8.4}{11}\selectfont\color{white}Gratuit \textbullet\ Google Play\\Tuteur IA, annales, concours, résultats\par}
      \end{minipage}
    \end{tcolorbox}
  \end{minipage}\hfill
  \begin{minipage}[t]{0.487\linewidth}
    \begin{tcolorbox}[enhanced,colback=poppurple,colframe=popink,boxrule=2.2pt,arc=10pt,
        drop shadow=popink,left=11pt,right=11pt,top=10pt,bottom=10pt,height=3.1cm,valign=center]
      \tikz[baseline=-0.7ex]\node[circle,fill=white,draw=popink,line width=1.3pt,minimum size=1.6cm,inner sep=0]{\popdisplay\fontsize{24}{24}\selectfont\color{poppurple}+};%
      \hspace{8pt}\begin{minipage}[c]{0.6\linewidth}
        {\popdisplay\fontsize{11}{12.5}\selectfont\color{white}\MakeUppercase{Passe à OnBuch+}}\par\vspace{4pt}
        {\raggedright\popsemi\fontsize{8.4}{11}\selectfont\color{white}Tous les cours signés, Tuteur IA illimité, fiches PDF — onglet OnBuch+ de l'app\par}
      \end{minipage}
    \end{tcolorbox}
  \end{minipage}
  \vspace{8pt}
  \popcontenu
  \vfill
  \signatureonbuch
  \restoregeometry
  \newpage
}

% Rangée « Ce cours contient » (page de garde)
\newcommand{\poptile}[3]{% #1 couleur #2 gros texte #3 légende
  \begin{tcolorbox}[enhanced,colback=#1,colframe=popink,boxrule=2pt,arc=9pt,shadow={2.5pt}{-2.5pt}{0pt}{popink},
      left=6pt,right=6pt,top=9pt,bottom=9pt,halign=center,valign=center,height=1.75cm,width=\linewidth,nobeforeafter]
    {\popdisplay\fontsize{12.5}{13}\selectfont\color{white}\MakeUppercase{#2}}\par\vspace{3pt}
    {\centering\popsemi\fontsize{7.8}{9.5}\selectfont\color{white}#3\par}
  \end{tcolorbox}}
\newcommand{\popcontenu}{%
  \par\noindent{\popxboldls\fontsize{7.5}{7.5}\selectfont\color{popmuted}CE COURS SIGNÉ CONTIENT}\par\vspace{7pt}
  \noindent\begin{minipage}[t]{0.235\linewidth}\poptile{popblue}{Cours}{complet, illustré, pas à pas}\end{minipage}\hfill
  \begin{minipage}[t]{0.235\linewidth}\poptile{poporange}{Méthodes}{et exercices résolus}\end{minipage}\hfill
  \begin{minipage}[t]{0.235\linewidth}\poptile{poppink}{Exercices}{type examen + corrigés}\end{minipage}\hfill
  \begin{minipage}[t]{0.235\linewidth}\poptile{popgreen}{Fiche bilan}{l'essentiel à retenir}\end{minipage}\par}

% Sommaire
\newtcolorbox{sommaire}{enhanced,colback=white,colframe=popink,boxrule=2.2pt,arc=10pt,
  drop shadow=popink,left=18pt,right=18pt,top=13pt,bottom=13pt,before skip=6pt,after skip=20pt,
  before upper={\poppill{Sommaire}{poppurple}{white}\par\vspace{9pt}\popsemi\fontsize{10.6}{14.5}\selectfont\color{popink}}}

% Signature de fin de cours — poussée tout en bas de la dernière page
\newcommand{\findecours}{%
  \par\vfill\nopagebreak
  \begin{center}\popsquiggle{poporange}\end{center}\vspace{4pt}
  \signatureonbuch
}
\AtBeginDocument{\def\llbracket{\mathopen{[\mkern-3.5mu[}}\def\rrbracket{\mathclose{]\mkern-3.5mu]}}} % intervalles d'entiers (glyphe ⟦ absent de la police)
% Glyphes absents de Fira Math : redéfinis à partir de glyphes disponibles
\AtBeginDocument{%
  \def\setminus{\mathbin{\backslash}}\let\smallsetminus\setminus
  \def\vdots{\mathord{\vbox{\baselineskip3.2pt\lineskiplimit0pt\kern4pt\hbox{.}\hbox{.}\hbox{.}}}}%
  \def\ddots{\mathinner{\mkern1mu\raise6.5pt\hbox{.}\mkern2mu\raise3.6pt\hbox{.}\mkern2mu\raise0.7pt\hbox{.}\mkern1mu}}%
  \def\triangle{\mathord{\scalebox{0.9}{$\symup{\Delta}$}}}%
  \def\nabla{\mathord{\raisebox{\depth}{\scalebox{1}[-1]{$\symup{\Delta}$}}}}%
  \def\checkmark{\popcheck{popdark}}%
  \def\star{\mathbin{\ast}}\def\bigstar{\mathord{\ast}}%
  \def\diamond{\mathbin{\scalebox{0.6}{$\Diamond$}}}%
  \def\bigcup{\mathop{\mathchoice{\raisebox{-0.3em}{\scalebox{1.7}{$\cup$}}}{\scalebox{1.25}{$\cup$}}{\cup}{\cup}}\displaylimits}%
  \def\bigcap{\mathop{\mathchoice{\raisebox{-0.3em}{\scalebox{1.7}{$\cap$}}}{\scalebox{1.25}{$\cap$}}{\cap}{\cap}}\displaylimits}%
  \def\lesseqqgtr{\mathrel{\lessgtr}}\def\bigoplus{\mathop{\oplus}\displaylimits}%
}
\providecommand{\ssi}{si et seulement si} % abréviation fréquente
\AtBeginDocument{\providecommand{\wideparen}{\overparen}} % arc de cercle

%% FIGFIX-BEGIN v5 — correctifs globaux des figures (docs/onbuch-generation/tools/figfix)
\usepackage{adjustbox}\usepackage{etoolbox}\tcbuselibrary{raster}
% toute figure TikZ/pgfplots plus large que la ligne est réduite à la largeur disponible
\BeforeBeginEnvironment{tikzpicture}{\begin{adjustbox}{max width=\linewidth}}
\AfterEndEnvironment{tikzpicture}{\end{adjustbox}}
% légendes pgfplots lisibles par-dessus les courbes
\pgfplotsset{every axis/.append style={legend style={fill=white,fill opacity=0.92,text opacity=1,draw=black!15,inner sep=3pt}}}
% étiquettes d'arêtes (syntaxe edge["..."]) avec fond blanc
\tikzset{every edge quotes/.append style={fill=white,inner sep=1.2pt}}
%% FIGFIX-END
\begin{document}

\pagegarde{Lesson 15 — Vocabulary: Words and expressions related to comparing prices or quality to determine the best buys for consumer economy and community resources/services. \# Semaine / Période Chapitre Titre Date effective Visa de l'enseignant}{Anglais}{Tle A}{Chapter 2 : Economic Life and Occupations}{EN15}{2 h}{This vocabulary lesson equips Terminale A students with the words, collocations and expressions used to compare prices and quality, identify the best buys, and talk about consumer economy and community services. Through thematic lexical fields, false friends, word families and contextual practice, students learn to use this lexis accurately both orally and in writing.
\vspace{4pt}
\begin{multicols}{2}\small
\begin{itemize}
\item À la fin de cette leçon, je sais utiliser le vocabulaire des prix (price, cost, fee, charge, fare) sans les confondre
\item À la fin de cette leçon, je sais comparer la qualité et les prix avec les bons adjectifs (cheap, affordable, overpriced, value for money)
\item À la fin de cette leçon, je sais employer les collocations du commerce (to bargain, to get a discount, value for money, to shop around)
\item À la fin de cette leçon, je sais éviter les faux amis et calques du français (economic/economical, to demand/to ask for, commerce)
\item À la fin de cette leçon, je sais former et utiliser les mots de la même famille (consume, consumer, consumption, consumerism)
\item À la fin de cette leçon, je sais parler des ressources et services communautaires (public services, amenities, facilities)
\end{itemize}\end{multicols}}

\begin{sommaire}
\begin{multicols}{2}
\begin{enumerate}
\item Talking about prices and money
\item Comparing quality and finding the best buy
\item Shopping, bargaining and consumer behaviour
\item Community resources and public services
\item Word families, false friends and pronunciation
\item Using the lexis in context
\item Activité d'intégration
\item Méthodes et exercices résolus
\item Exercices
\item Corrigés des exercices
\item Fiche bilan
\end{enumerate}
\end{multicols}
\end{sommaire}

\begin{objectifs}
\begin{itemize}
\item À la fin de cette leçon, je sais utiliser le vocabulaire des prix (price, cost, fee, charge, fare) sans les confondre
\item À la fin de cette leçon, je sais comparer la qualité et les prix avec les bons adjectifs (cheap, affordable, overpriced, value for money)
\item À la fin de cette leçon, je sais employer les collocations du commerce (to bargain, to get a discount, value for money, to shop around)
\item À la fin de cette leçon, je sais éviter les faux amis et calques du français (economic/economical, to demand/to ask for, commerce)
\item À la fin de cette leçon, je sais former et utiliser les mots de la même famille (consume, consumer, consumption, consumerism)
\item À la fin de cette leçon, je sais parler des ressources et services communautaires (public services, amenities, facilities)
\item À la fin de cette leçon, je sais réemployer ce lexique dans des phrases, dialogues et courts textes au Bac
\end{itemize}
\end{objectifs}

\begin{prerequis}
\begin{itemize}
\item Les comparatifs et superlatifs (cheaper than, the best) vus en classe de Première
\item Le vocabulaire de base de l'argent et des achats (money, to buy, to sell, market, shop)
\item La formation des noms et adjectifs par suffixes (-er, -tion, -al, -less)
\item L'expression de l'opinion simple (I think, in my opinion)
\end{itemize}
\end{prerequis}

\newpage

\coursec{Talking about prices and money}

\courssub{Situation d'accroche : At Mokolo Market}

\begin{textebox}[At Mokolo Market, Yaoundé]
At Mokolo Market in Yaoundé, Amina wants to buy a bag of rice. One seller offers it at 12,000 FCFA, another at 10,500 FCFA, but the cheaper bag contains broken grains. Which one is the \cle{best buy}? Every day, Cameroonian consumers compare prices and quality before spending their money. In this lesson, you will learn the English words and expressions you need to compare, bargain and choose wisely — a key skill for consumer economy and for managing community resources.
\end{textebox}

\textbf{Problématique.} Comment dire, en anglais correct et naturel, qu'un prix est élevé, qu'un service est facturé, qu'un billet de car coûte plus cher qu'avant ? Le français utilise souvent le mot unique \emph{prix} ou \emph{coût} ; l'anglais, lui, possède une famille de mots spécialisés — \eng{price}, \eng{cost}, \eng{fee}, \eng{charge}, \eng{fare}, \eng{rate} — qu'il faut savoir choisir avec précision. C'est tout l'objet de cette première section.

\begin{definition}[Best buy]
A \cle{best buy} is the product that offers the best quality for the lowest price. En français : \emph{le meilleur achat, le meilleur rapport qualité-prix}. \eng{This bag of rice is the best buy: good quality at a fair price.} « Ce sac de riz est le meilleur achat : bonne qualité à un prix raisonnable. »
\end{definition}

\courssub{Observation : un texte pour découvrir le lexique}

\begin{textebox}[Travelling on a budget — a student's diary]
Last Saturday, I travelled from Bafoussam to Douala to visit my uncle. The bus \cle{fare} has gone up again: it now costs 3,500 FCFA, but the agency charges no booking \cle{fee}, so it is still the best buy for students like me. When I arrived, my uncle complained about the \cle{cost} of living in Douala. “Everything is expensive here,” he said. “The \cle{rent} for my small flat is 60,000 FCFA a month, and the electricity \cle{bill} keeps rising. Even the water \cle{rate} has increased.” On Sunday, we went to a restaurant. The waiter brought the bill, and my uncle left a small \cle{tip} because the service was good. On Monday morning, I went to the bank to change some money. The exchange \cle{rate} was favourable, and the bank charged only a small commission. Before leaving, I asked for a \cle{receipt}. Back in Bafoussam, I calculated my expenses: I had spent 15,000 FCFA in three days. “I must \cle{save} more money,” I told myself, “and stop wasting it on things I don't need. Next time, I will compare prices more carefully before I buy anything.”

\textbf{Glossaire.} fare (n.) : prix du transport ; fee (n.) : honoraires, frais ; rent (n.) : loyer ; bill (n.) : facture, addition ; tip (n.) : pourboire ; rate (n.) : tarif, taux ; receipt (n.) : reçu ; to save (v.) : économiser ; to waste (v.) : gaspiller.
\end{textebox}

\textbf{Questions de compréhension.}
\begin{enumerate}
\item How much is the bus fare from Bafoussam to Douala?
\item Why does the narrator's uncle find life in Douala expensive?
\item Why did the uncle leave a tip at the restaurant?
\item What two things did the narrator do at the bank?
\item What resolution does the narrator make at the end?
\end{enumerate}

\courssub{Le champ lexical du prix : price, cost, fee, charge, fare, rate}

Observons d'abord les phrases suivantes, tirées du texte ou de la vie quotidienne :

\begin{exemplebox}[Observer avant de comprendre]
\eng{The price of this bag is 5,000 FCFA.} « Le prix de ce sac est de 5 000 FCFA. »\\
\eng{The cost of living in Douala is high.} « Le coût de la vie à Douala est élevé. »\\
\eng{School fees have increased this year.} « Les frais de scolarité ont augmenté cette année. »\\
\eng{The agency charges no booking fee.} « L'agence ne facture aucun frais de réservation. »\\
\eng{The bus fare from Bafoussam to Douala has gone up.} « Le prix du car de Bafoussam à Douala a augmenté. »\\
\eng{The exchange rate is favourable today.} « Le taux de change est favorable aujourd'hui. »
\end{exemplebox}

\begin{propriete}[Choisir le bon mot : price, cost, fee, charge, fare, rate, rent, bill, fine, tip]
Chaque mot correspond à un type de dépense précis :
\begin{itemize}
\item \cle{price} : le prix d'un article, d'un produit précis. \eng{What is the price of this phone?}
\item \cle{cost} : le coût global, souvent abstrait (cost of living, cost of a project). \eng{The cost of the new bridge was enormous.}
\item \cle{fee} : les honoraires, les frais payés pour un service professionnel ou officiel : \eng{school fees} (frais de scolarité), \eng{lawyer's fees} (honoraires d'avocat), \eng{entrance fee} (droits d'entrée).
\item \cle{charge} : les frais facturés pour un service : \eng{a service charge}, \eng{delivery charges} (frais de livraison).
\item \cle{fare} : le prix d'un transport (bus, taxi, train, avion) : \eng{bus fare}, \eng{taxi fare}.
\item \cle{rate} : un tarif, un taux : \eng{exchange rate} (taux de change), \eng{interest rate} (taux d'intérêt), \eng{hourly rate} (tarif horaire).
\item \cle{rent} : le loyer. \eng{The rent for this flat is 60,000 FCFA a month.}
\item \cle{bill} : la facture (électricité, eau) ou l'addition au restaurant. \eng{Can I have the bill, please?}
\item \cle{fine} : l'amende. \eng{He paid a fine for parking in the wrong place.}
\item \cle{tip} : le pourboire. \eng{She left a tip for the waiter.}
\end{itemize}
\end{propriete}

\begin{attention}[Faux amis et calques du français]
\begin{itemize}
\item \eng{charge} ne signifie pas « charge » au sens de fardeau : c'est un \emph{frais facturé}. Un fardeau se dit \eng{a burden} ou \eng{a load}.
\item \eng{fare} = prix du transport ; ne pas confondre avec \eng{fair} (juste, équitable), qui se prononce exactement pareil : « fèr ».
\item \eng{rent} = loyer que l'on paie ; « louer » (une voiture) se dit \eng{to rent a car} ou \eng{to hire a car}.
\item \eng{bill} = facture ou addition ; un billet de banque se dit \eng{a banknote} (GB) ou \eng{a bill} (US uniquement).
\end{itemize}
\end{attention}

\courssub{Les verbes de l'argent : to cost, to charge, to pay, to spend, to afford, to save, to waste, to owe}

\begin{propriete}[Emploi et constructions des verbes]
\begin{itemize}
\item \cle{to cost} (cost, cost) : le sujet est la chose. \eng{This bag costs 5,000 FCFA.} « Ce sac coûte 5 000 FCFA. »
\item \cle{to charge} : le vendeur facture. Structure : \eng{to charge somebody an amount for something}. \eng{They charged me 500 FCFA for delivery.}
\item \cle{to pay} (paid, paid) : \eng{to pay an amount for something}. \eng{I paid 3,500 FCFA for the ticket.}
\item \cle{to spend} (spent, spent) : \eng{to spend money on something}. \eng{She spent 15,000 FCFA on food.}
\item \cle{to afford} : pouvoir se permettre (financièrement), presque toujours avec \eng{can / can't}. \eng{I can't afford a new phone.} « Je ne peux pas me permettre d'acheter un nouveau téléphone. »
\item \cle{to save} : économiser. \eng{He saves 10,000 FCFA every month.}
\item \cle{to waste} : gaspiller. \eng{Don't waste your money on useless things.}
\item \cle{to owe} : devoir (de l'argent à quelqu'un). \eng{I owe my brother 2,000 FCFA.}
\end{itemize}
\end{propriete}

\begin{exemplebox}[Exemples traduits]
\eng{School fees have increased.} « Les frais de scolarité ont augmenté. »\\
\eng{How much did you pay for this watch?} « Combien as-tu payé cette montre ? »\\
\eng{We can't afford to waste water.} « Nous ne pouvons pas nous permettre de gaspiller l'eau. »\\
\eng{The mechanic charged me 8,000 FCFA for the repair.} « Le mécanicien m'a facturé 8 000 FCFA pour la réparation. »
\end{exemplebox}

\courssub{Expressions utiles : parler de la hauteur d'un prix}

\begin{propriete}[Expressions figées à mémoriser]
\begin{itemize}
\item \eng{at a high price} / \eng{at a low price} : à un prix élevé / bas. \eng{They sell rice at a low price.}
\item \eng{free of charge} : gratuit. \eng{The advice is free of charge.} « Le conseil est gratuit. »
\item \eng{at half price} : à moitié prix. \eng{I bought these shoes at half price.}
\item \eng{prices go up} / \eng{prices come down} : les prix augmentent / baissent. \eng{Fuel prices have gone up again.}
\item \eng{value for money} : le rapport qualité-prix. \eng{This restaurant offers good value for money.}
\end{itemize}
\end{propriete}

\begin{attention}[Ne dites jamais “the price is expensive”]
Erreur classique des francophones, calquée sur « le prix est cher ». En anglais, un \cle{prix} est \eng{high} ou \eng{low} ; un \cle{produit} est \eng{expensive} ou \eng{cheap}.\\
\eng{The price is high.} « Le prix est élevé. »\\
\eng{The bag is expensive.} « Le sac est cher. »\\
Interdit : \emph{The price is expensive.}
\end{attention}

\begin{center}
\begin{tabularx}{\linewidth}{|X|X|X|}
\hline
\rowcolor{popblueL} Français & English & Remarque \\
\hline
le prix d'un article & the price & \eng{the price of bread} \\
\hline
le coût de la vie & the cost of living & sens global, abstrait \\
\hline
les frais de scolarité & school fees & toujours au pluriel \\
\hline
les honoraires d'avocat & lawyer's fees & service professionnel \\
\hline
les frais de livraison & delivery charges & service facturé \\
\hline
le prix du car / du taxi & the bus / taxi fare & transports uniquement \\
\hline
le taux de change & the exchange rate & tarif, taux \\
\hline
le loyer & the rent & \eng{to pay the rent} \\
\hline
la facture d'électricité & the electricity bill & aussi : addition au restaurant \\
\hline
l'amende & the fine & \eng{to pay a fine} \\
\hline
le pourboire & the tip & \eng{to leave a tip} \\
\hline
gratuit & free of charge & expression figée \\
\hline
le prix est élevé & the price is high & jamais \emph{expensive} avec \eng{price} \\
\hline
\end{tabularx}
\end{center}

\courssub{Carte mentale : organiser le lexique}

\begin{popfigure}
\begin{tikzpicture}[mindmap, grow cyclic, every node/.style={concept, font=\small},
root concept/.append style={concept color=popblue, font=\small\bfseries},
level 1/.append style={level distance=3.2cm, sibling angle=90},
level 2/.append style={level distance=2.2cm, sibling angle=45, font=\scriptsize}]
\node[root concept]{MONEY \& PRICES}
  child[concept color=poporange]{ node {nouns}
    child { node {price, cost} }
    child { node {fee, charge} }
    child { node {fare, rate} }
    child { node {rent, bill, fine, tip} }
  }
  child[concept color=popgreen]{ node {verbs}
    child { node {cost, charge} }
    child { node {pay, spend} }
    child { node {afford, save} }
    child { node {waste, owe} }
  }
  child[concept color=poppurple]{ node {adjectives}
    child { node {expensive, cheap} }
    child { node {high, low} }
    child { node {free, fair} }
  }
  child[concept color=poppink]{ node {expressions}
    child { node {free of charge} }
    child { node {at half price} }
    child { node {go up, come down} }
    child { node {value for money} }
  };
\end{tikzpicture}
\legende{Carte mentale du vocabulaire de l'argent et des prix : quatre familles à mémoriser.}
\end{popfigure}

\courssub{Prononciation et orthographe : les pièges}

\begin{propriete}[Mots difficiles à prononcer et à orthographier]
\begin{itemize}
\item \eng{price} se prononce « praïce » : le \emph{i} est une diphtongue, comme dans \eng{my}, \eng{five}.
\item \eng{charge} se prononce « tchârdj » : \emph{ch} = « tch » et \emph{ge} final = « dj ».
\item \eng{receipt} se prononce « rissite » : le \emph{p} est muet, comme dans \eng{psychology} ou \eng{pneumonia}.
\item \eng{fare} et \eng{fair} se prononcent tous deux « fèr » : attention au sens.
\item \eng{owe} se prononce « ô » : le \emph{w} est muet.
\end{itemize}
\end{propriete}

\begin{savaistu}[Le P muet de “receipt”]
En anglais, \eng{receipt} (reçu) s'écrit avec un \emph{p} qui ne se prononce pas. Ce \emph{p} vient du latin \emph{recepta} : les savants de la Renaissance l'ont réintroduit dans l'orthographe anglaise pour rappeler l'origine latine du mot, alors que le français \emph{reçu} l'a perdu. C'est un piège classique à l'écrit du Baccalauréat : beaucoup de candidats écrivent \emph{receit} ou \emph{reciept}. Retenez : \textbf{re-cei-pt}, avec \emph{ei} après le \emph{c} (comme dans \eng{receive}, \eng{deceive}) et un \emph{p} silencieux.
\end{savaistu}

\courssub{Exercice résolu}

\begin{exoresolu}[Choose the right word]
Complétez chaque phrase avec le mot qui convient : \eng{price}, \eng{cost}, \eng{fee}, \eng{fare}, \eng{rate}, \eng{bill}, \eng{fine}, \eng{tip}.
\begin{enumerate}
\item The \dots\ of living in Yaoundé is rising every year.
\item The taxi \dots\ from the airport to the town centre is 5,000 FCFA.
\item My parents pay my school \dots\ every term.
\item The waiter brought the \dots\ and we left a \dots\ on the table.
\item He had to pay a \dots\ because he drove too fast.
\item What is the exchange \dots\ for the euro today?
\item The \dots\ of this dictionary is 8,000 FCFA.
\end{enumerate}
\tcblower
\textbf{Solution détaillée.}
\begin{enumerate}
\item \textbf{cost} — il s'agit du coût global de la vie (\eng{the cost of living}), expression figée.
\item \textbf{fare} — prix d'un transport (taxi) : on dit toujours \eng{taxi fare}, \eng{bus fare}.
\item \textbf{fees} — frais de scolarité : \eng{school fees}, toujours au pluriel dans ce sens.
\item \textbf{bill} ; \textbf{tip} — au restaurant, on demande \eng{the bill} (l'addition) et on laisse \eng{a tip} (un pourboire).
\item \textbf{fine} — une amende pour excès de vitesse : \eng{to pay a fine}.
\item \textbf{rate} — taux de change : \eng{exchange rate}.
\item \textbf{price} — le prix d'un article précis : \eng{the price of this dictionary}.
\end{enumerate}
\end{exoresolu}

\begin{exercice}[À vous de jouer]
Traduisez en anglais :
\begin{enumerate}
\item Je ne peux pas me permettre d'acheter une nouvelle moto.
\item Le prix du car a augmenté, mais c'est encore le meilleur achat pour les étudiants.
\item Le mécanicien m'a facturé 10 000 FCFA pour la réparation.
\item Ne gaspille pas ton argent ; économise pour tes frais de scolarité.
\item Ce sac est trop cher, mais l'autre est à moitié prix.
\end{enumerate}
\end{exercice}

\begin{corrige}[Exercice « À vous de jouer »]
\begin{enumerate}
\item \eng{I can't afford to buy a new motorbike.}
\item \eng{The bus fare has gone up, but it is still the best buy for students.}
\item \eng{The mechanic charged me 10,000 FCFA for the repair.}
\item \eng{Don't waste your money; save for your school fees.}
\item \eng{This bag is too expensive, but the other one is at half price.}
\end{enumerate}
\end{corrige}

\begin{aretenir}[L'essentiel de la section]
\begin{itemize}
\item Un article a un \eng{price} ; la vie a un \eng{cost} ; un professionnel facture des \eng{fees} ; un service entraîne des \eng{charges} ; un transport a un \eng{fare} ; une banque applique un \eng{rate}.
\item Un prix est \eng{high} ou \eng{low} ; un produit est \eng{expensive} ou \eng{cheap}.
\item \eng{to afford} s'emploie avec \eng{can} : \eng{I can't afford it.}
\item \eng{free of charge} = gratuit ; \eng{at half price} = à moitié prix ; \eng{prices go up / come down} = les prix montent / baissent.
\item \eng{receipt} s'écrit avec un \emph{p} muet et se prononce « rissite ».
\end{itemize}
\end{aretenir}

\coursec{Comparing quality and finding the best buy}

Dans la section précédente, nous avons appris à \cle{parler des prix et de l'argent} : demander un prix, dire combien coûte un article, parler de monnaie et de budget. Mais connaître le prix ne suffit pas pour bien acheter. Un consommateur avisé doit aussi \cle{juger la qualité} d'un produit et \cle{comparer} plusieurs articles avant de décider lequel est \eng{the best buy} — le meilleur achat. C'est précisément l'objet de cette section : le vocabulaire et les structures qui permettent de comparer la qualité et les prix pour faire le bon choix.

\courssub{Observing the language: a dialogue at the market}

Avant toute règle, observons le vocabulaire en action. Lis le mini-dialogue suivant, qui se déroule au marché Mokolo à Yaoundé : deux amies, Sandrine et Mirabelle, comparent deux sacs à main.

\begin{textebox}[At Mokolo Market — “Which bag is the best buy?”]
\textbf{Sandrine:} Look at these two handbags, Mirabelle. The leather one costs 15,000 FCFA and the plastic one is only 5,000 FCFA.\par
\textbf{Mirabelle:} The plastic one is much cheaper, but look at the stitching — it is already coming apart. It looks cheap, honestly.\par
\textbf{Sandrine:} You are right. The leather bag is more expensive, but it is genuine leather, and it feels very durable. It will last for years.\par
\textbf{Mirabelle:} Let's weigh up the pros and cons. The plastic bag is affordable, but it is poor quality — it will probably break in a month. The leather bag costs three times more, but it is reliable and well made.\par
\textbf{Sandrine:} So which one is better value for money?\par
\textbf{Mirabelle:} Definitely the leather one. It is not as cheap as the plastic one, but the quality is far better. At 15,000 FCFA, it is a real bargain — genuine leather usually costs twice that price in town.\par
\textbf{Sandrine:} Wait — is it really genuine? Some traders here sell fake leather and pass it off as genuine.\par
\textbf{Mirabelle:} Good point. Let me check the label and smell it. Yes, it smells like real leather, and the brand label is original, not counterfeit.\par
\textbf{Sandrine:} Then it is worth the money. But before we buy, let's shop around and check prices in the other stalls. Maybe we can find an even better deal.\par
\textbf{Mirabelle:} Good idea. A wise shopper always compares before paying!
\end{textebox}

\textbf{Glossaire :}

\begin{center}
\begin{tabularx}{\linewidth}{|l|l|X|}
\hline
\rowcolor{popblueL} \textbf{Mot anglais} & \textbf{Nature} & \textbf{Traduction française} \\
\hline
stitching & nom & couture (d'un vêtement, d'un sac) \\
\hline
to come apart & verbe & se découdre, se défaire \\
\hline
genuine & adjectif & authentique, véritable \\
\hline
durable & adjectif & durable, résistant \\
\hline
reliable & adjectif & fiable \\
\hline
stall & nom & étal, stand (au marché) \\
\hline
counterfeit & adjectif/nom & contrefait / contrefaçon \\
\hline
wise & adjectif & sage, avisé \\
\hline
\end{tabularx}
\end{center}

\textbf{Questions de compréhension :}

\begin{enumerate}
\item How much does each handbag cost?
\item Why does Mirabelle reject the plastic bag?
\item What does Sandrine fear about the leather bag, and how does Mirabelle check it?
\item Which bag is “the best buy” according to the two friends, and why?
\item What do they decide to do before paying?
\end{enumerate}

\courssub{Adjectives of quality}

À partir du dialogue, classons d'abord les \cle{adjectifs de qualité}. Ils permettent de dire si un produit est bien fait, authentique, neuf ou d'occasion.

\begin{definition}[Quality adjectives — les adjectifs de qualité]
\begin{itemize}
\item \eng{high-quality} : de haute qualité. \eng{This is a high-quality product.} « C'est un produit de haute qualité. »
\item \eng{durable} : durable, qui dure longtemps. \eng{Leather shoes are durable.} « Les chaussures en cuir sont durables. »
\item \eng{reliable} : fiable, sur lequel on peut compter. \eng{Toyota cars are reliable.} « Les voitures Toyota sont fiables. »
\item \eng{genuine} : authentique, véritable (contraire : \eng{fake}, \eng{counterfeit}). \eng{Is this note genuine or fake?} « Ce billet est-il authentique ou faux ? »
\item \eng{second-hand} : d'occasion. \eng{She bought a second-hand laptop in Douala.} « Elle a acheté un ordinateur portable d'occasion à Douala. »
\item \eng{brand new} : tout neuf, flambant neuf. \eng{He rides a brand new motorbike.} « Il conduit une moto toute neuve. »
\item \eng{good value} : de bon rapport qualité-prix. \eng{This market offers good value.} « Ce marché offre un bon rapport qualité-prix. »
\end{itemize}
\end{definition}

\begin{attention}[Fake, counterfeit, false — ne confonds pas]
\eng{Fake} est le terme courant pour « faux » (produit, billet, document) : \eng{a fake watch} « une fausse montre ». \eng{Counterfeit} est plus soutenu et désigne surtout la contrefaçon de marques et de billets : \eng{counterfeit money} « de la fausse monnaie ». Attention : on ne dit jamais \eng{a false watch} pour une montre contrefaite ; \eng{false} s'emploie pour les idées (\eng{a false idea}, \eng{false teeth} « un dentier »).
\end{attention}

\courssub{Adjectives of price}

Passons maintenant aux \cle{adjectifs de prix}, qui évaluent le coût d'un article.

\begin{definition}[Price adjectives — les adjectifs de prix]
\begin{itemize}
\item \eng{cheap} : bon marché, pas cher (parfois péjoratif : de mauvaise qualité).
\item \eng{affordable} : abordable (sens positif : à la portée de ta bourse). \eng{The school fees are affordable.} « Les frais de scolarité sont abordables. »
\item \eng{inexpensive} : peu coûteux (sens positif, plus formel que \eng{cheap}).
\item \eng{reasonable} : raisonnable (prix juste). \eng{The price is reasonable.} « Le prix est raisonnable. »
\item \eng{expensive} : cher, coûteux.
\item \eng{overpriced} : trop cher, surfait (le prix ne correspond pas à la valeur réelle). \eng{These imported biscuits are overpriced.} « Ces biscuits importés sont trop chers. »
\item \eng{priceless} : inestimable, qui n'a pas de prix (sens très positif). \eng{Her friendship is priceless.} « Son amitié n'a pas de prix. »
\item \eng{worthless} : sans valeur, qui ne vaut rien. \eng{The fake phone stopped working after two days — it is worthless.} « Le faux téléphone a cessé de fonctionner au bout de deux jours — il ne vaut rien. »
\end{itemize}
\end{definition}

\begin{attention}[« Cheap » : un mot à double visage]
En français, « pas cher » est neutre ou positif. En anglais, \eng{cheap} peut être \cle{péjoratif} : il suggère souvent « de mauvaise qualité, bas de gamme ». \eng{These shoes are cheap — they fell apart in a week.} « Ces chaussures sont de mauvaise qualité — elles se sont abîmées en une semaine. » Pour un sens positif, préfère \eng{affordable}, \eng{inexpensive} ou \eng{good value}. De même, attention au faux ami : \eng{priceless} ne veut pas dire « sans valeur » (c'est \eng{worthless}) mais « inestimable » !
\end{attention}

\courssub{Key expressions: judging a purchase}

Voici les \cle{expressions idiomatiques} essentielles pour donner ton jugement sur un achat.

\begin{definition}[Value for money]
\eng{Value for money} signifie \emph{the quality you get compared with the price you pay} — la qualité obtenue par rapport au prix payé, c'est-à-dire le \cle{rapport qualité-prix}. \eng{This restaurant offers excellent value for money.} « Ce restaurant offre un excellent rapport qualité-prix. » On dit aussi \eng{good value} : \eng{The meal was good value.} « Le repas valait son prix. »
\end{definition}

\begin{exemplebox}[Expressions clés en contexte]
\begin{itemize}
\item \eng{a bargain} : une bonne affaire. \eng{This phone is a real bargain at 25,000 FCFA.} « Ce téléphone est une vraie bonne affaire à 25 000 FCFA. »
\item \eng{a good deal} : une bonne affaire, un bon arrangement. \eng{We got a good deal on the second-hand fridge.} « Nous avons fait une bonne affaire avec le réfrigérateur d'occasion. »
\item \eng{a rip-off} : une arnaque (familier). \eng{That shop is a rip-off!} « Ce magasin, c'est de l'arnaque ! »
\item \eng{worth the money} : qui vaut son prix. \eng{The bag is expensive, but it is worth the money.} « Le sac est cher, mais il vaut son prix. »
\item \eng{to be worth it} : en valoir la peine. \eng{Paying more for genuine parts is worth it.} « Payer plus cher pour des pièces authentiques en vaut la peine. »
\end{itemize}
\end{exemplebox}

\begin{popfigure}
\begin{center}
\begin{tikzpicture}
\node[draw=popgreen, fill=popgreenL, rounded corners, thick, minimum width=6.5cm, minimum height=1cm, align=center, font=\bfseries] (pos) at (0,0) {POSITIVE judgement};
\node[draw=popgreen, fill=white, rounded corners, align=center, minimum width=6.5cm] (poslist) at (0,-2.4) {a bargain\\[2pt] good value (for money)\\[2pt] a good deal\\[2pt] worth the money / worth it\\[2pt] affordable, reasonable};
\node[draw=popink, fill=poppinkL, rounded corners, thick, minimum width=6.5cm, minimum height=1cm, align=center, font=\bfseries] (neg) at (8.5,0) {NEGATIVE judgement};
\node[draw=popink, fill=white, rounded corners, align=center, minimum width=6.5cm] (neglist) at (8.5,-2.4) {a rip-off\\[2pt] overpriced\\[2pt] worthless\\[2pt] poor quality\\[2pt] fake, counterfeit};
\end{tikzpicture}
\end{center}
\legende{Deux colonnes de jugement : à gauche, les expressions positives pour un bon achat ; à droite, les expressions négatives pour un mauvais achat.}
\end{popfigure}

\courssub{Verbs of comparing}

Pour trouver le meilleur achat, il faut \cle{comparer} activement. Voici les verbes utiles.

\begin{definition}[Verbs of comparison — les verbes de comparaison]
\begin{itemize}
\item \eng{to compare} : comparer. \eng{Compare the two phones before you choose.} « Compare les deux téléphones avant de choisir. » Construction : \eng{to compare A with/to B}.
\item \eng{to shop around} : comparer les prix dans plusieurs magasins avant d'acheter. \eng{We shopped around before buying the television.} « Nous avons comparé les prix dans plusieurs magasins avant d'acheter la télévision. »
\item \eng{to check prices} : vérifier les prix. \eng{Always check prices online first.} « Vérifie toujours les prix en ligne d'abord. »
\item \eng{to weigh up the pros and cons} : peser le pour et le contre. \eng{Let us weigh up the pros and cons of each option.} « Pesons le pour et le contre de chaque option. »
\end{itemize}
\end{definition}

\begin{attention}[Compare... with, pas « compare... and »]
En français on dit « comparer A et B » ; en anglais, la préposition est obligatoire : \eng{to compare A \textbf{with} B} (ou \eng{to}). Ne dis jamais \eng{compare the two prices and choose} pour « compare les deux prix » ; dis \eng{compare the two prices}. Autre piège : \eng{to shop around} est inséparable — on ne peut pas dire \eng{to shop around the market} au sens de « comparer » ; on dira \eng{to shop around in the market} ou \eng{to shop around for a better price}.
\end{attention}

\courssub{Comparison structures}

Pour formuler la comparaison, l'anglais utilise des structures précises que tu connais déjà ; rappelons-les appliquées au vocabulaire de cette leçon.

\begin{propriete}[Les structures de comparaison appliquées aux achats]
\begin{itemize}
\item \textbf{Comparatif de supériorité} : adjectif court + \eng{-er than} ; \eng{more} + adjectif long + \eng{than}. \eng{The plastic bag is cheaper than the leather one.} « Le sac en plastique est moins cher que celui en cuir. » \eng{This shop is more expensive than the market.} « Ce magasin est plus cher que le marché. »
\item \textbf{Comparatif avec « value »} : \eng{better value than}. \eng{The leather bag is better value than the plastic one.} « Le sac en cuir a un meilleur rapport qualité-prix que celui en plastique. »
\item \textbf{Superlatif} : \eng{the most} + adjectif long ; adjectif court + \eng{-est}. \eng{Mokolo Market is the most affordable place to buy vegetables.} « Le marché Mokolo est l'endroit le plus abordable pour acheter des légumes. »
\item \textbf{Comparatif d'infériorité} : \eng{not as ... as}. \eng{The second-hand phone is not as expensive as the brand new one.} « Le téléphone d'occasion n'est pas aussi cher que le neuf. » \eng{This fabric is not as durable as that one.} « Ce tissu n'est pas aussi résistant que celui-là. »
\end{itemize}
\end{propriete}

\begin{center}
\begin{tabularx}{\linewidth}{|X|X|X|}
\hline
\rowcolor{popblueL} \textbf{Français} & \textbf{English} & \textbf{Remarque} \\
\hline
moins cher que & cheaper than / less expensive than & \eng{less expensive} est plus formel \\
\hline
un meilleur rapport qualité-prix que & better value (for money) than & \eng{value} est invariable ici \\
\hline
le plus abordable & the most affordable / the cheapest & superlatif \\
\hline
pas aussi cher que & not as expensive as & \eng{as ... as} avec négation \\
\hline
une bonne affaire & a bargain / a good deal & \eng{bargain} : prononcer « bar-guine » \\
\hline
une arnaque & a rip-off & familier, à l'oral \\
\hline
en valoir la peine & to be worth it & jamais \eng{to worth it} \\
\hline
\end{tabularx}
\end{center}

\begin{methode}[How to compare two products in writing]
Pour rédiger un court paragraphe comparant deux produits, suis ces trois étapes :
\begin{enumerate}
\item \textbf{State the price of each} : indique le prix de chaque article. \eng{Bag A costs 5,000 FCFA; bag B costs 15,000 FCFA.}
\item \textbf{Compare quality with adjectives} : compare la qualité avec les adjectifs et structures de cette leçon. \eng{Bag B is more expensive, but it is genuine leather, durable and reliable, whereas bag A looks cheap and poorly made.}
\item \textbf{Conclude} : conclus avec la formule \eng{the best buy is ... because ...}. \eng{The best buy is bag B because it offers better value for money.}
\end{enumerate}
\end{methode}

\begin{exoresolu}[Choosing the best buy]
Énoncé — Complete the sentences with a word or expression from the lesson: \emph{bargain, overpriced, worth it, shop around, second-hand, not as ... as}.\par
1. This radio costs 40,000 FCFA here but only 20,000 FCFA downtown — it is completely \dots\ in this shop.\par
2. My father bought a \dots\ car; it is old but very reliable.\par
3. Before buying your textbooks, you should \dots\ to find the lowest price.\par
4. The market is \dots\ expensive \dots\ the supermarket.\par
5. At 10,000 FCFA for genuine leather shoes, that is a real \dots !
\tcblower
Solution détaillée :\par
1. \eng{overpriced} — le prix (40 000 FCFA) est le double du prix réel : l'article est « trop cher ».\par
2. \eng{second-hand} — « d'occasion » ; la voiture est vieille mais fiable.\par
3. \eng{shop around} — « comparer les prix dans plusieurs magasins » avant d'acheter.\par
4. \eng{not as ... as} — comparatif d'infériorité : \eng{The market is not as expensive as the supermarket.} « Le marché n'est pas aussi cher que le supermarché. »\par
5. \eng{bargain} — des chaussures en cuir véritable à 10 000 FCFA, c'est « une vraie bonne affaire ».
\end{exoresolu}

\begin{exercice}[Your turn — compare and conclude]
Two phones are on sale in Buea. Phone A: brand new, 80,000 FCFA, reliable brand, two-year guarantee. Phone B: second-hand, 30,000 FCFA, no guarantee, battery problems. Write a five-sentence paragraph comparing the two phones, using at least three expressions from this lesson (\emph{value for money, worth it, not as ... as, a bargain, overpriced, durable}). End with: \eng{The best buy is ... because ...}.
\end{exercice}

\begin{corrige}[Exercice « Your turn »]
Modèle de production possible : \eng{Phone A costs 80,000 FCFA, while phone B costs only 30,000 FCFA. Phone B is much cheaper, but it is second-hand and has battery problems, so it may not be durable. Phone A is more expensive, but it is brand new, reliable and has a two-year guarantee. Phone B is not as reliable as phone A, and repairing it could cost a lot. The best buy is phone A because it offers better value for money and is worth the extra cost.} « Le téléphone A coûte 80 000 FCFA, tandis que le B ne coûte que 30 000 FCFA. Le B est bien moins cher, mais c'est une occasion avec des problèmes de batterie, donc il risque de ne pas durer. Le A est plus cher, mais il est neuf, fiable et garanti deux ans. Le B n'est pas aussi fiable que le A, et le réparer pourrait coûter cher. Le meilleur achat est le téléphone A car il offre un meilleur rapport qualité-prix et vaut la dépense supplémentaire. »
\end{corrige}

\begin{experience}[Pair work — at the stall]
Par groupes de deux : l'un joue le vendeur au marché, l'autre le client. Le client compare deux articles (par exemple deux sacs de riz, deux paires de chaussures) en utilisant au moins quatre expressions de la leçon : \eng{How much is this one? Is it genuine? This one is cheaper, but is it durable? Let me weigh up the pros and cons. It is worth the money.} Puis échangez les rôles.
\end{experience}

\begin{savaistu}[“Best buy” awards in Britain]
Au Royaume-Uni, les magazines de consommateurs comme \emph{Which?} testent chaque année des centaines de produits — téléphones, voitures, machines à laver — et décernent des \eng{best buy} awards, des récompenses de « meilleur achat », pour guider les acheteurs vers le meilleur rapport qualité-prix. L'expression \eng{best buy} est ainsi devenue un label de confiance : quand un produit porte la mention \eng{Which? Best Buy}, les consommateurs britanniques savent qu'il a été testé et approuvé par des experts indépendants.
\end{savaistu}

\begin{aretenir}[L'essentiel de la section]
\begin{itemize}
\item Qualité : \eng{high-quality, durable, reliable, genuine} (authentique) vs \eng{fake, counterfeit, second-hand, brand new}.
\item Prix : \eng{cheap} (parfois péjoratif !), \eng{affordable, reasonable, expensive, overpriced} ; attention aux faux amis \eng{priceless} (inestimable) et \eng{worthless} (sans valeur).
\item Jugement : \eng{a bargain, a good deal, value for money, worth it} vs \eng{a rip-off, overpriced}.
\item Comparer : \eng{to compare A with B, to shop around, to check prices, to weigh up the pros and cons}.
\item Structures : \eng{cheaper than, better value than, the most affordable, not as expensive as}.
\item Conclusion d'un paragraphe de comparaison : \eng{The best buy is ... because ...}.
\end{itemize}
\end{aretenir}

\coursec{Shopping, bargaining and consumer behaviour}

In the previous section, we learnt how to compare prices and quality in order to find the \cle{best buy}. But finding the best buy is only the beginning of the story. Once you have chosen your product, you still have to \emph{buy} it: you must speak to the seller, perhaps \cle{bargain} over the price, pay, keep your \cle{receipt}, and sometimes \cle{complain} or ask for a \cle{refund}. In this section, we will learn all the words and expressions you need to shop confidently, to negotiate prices like a customer at Mokolo market in Yaoundé, and to defend your rights as a consumer.

\courssub{Observing the language: a dialogue at the market}

\begin{textebox}[At the market — a model dialogue]
\textbf{Customer:} Good morning, madam. How much is this bag?\\
\textbf{Vendor:} It costs 8,000 CFA francs, my dear. Very good quality, genuine leather.\\
\textbf{Customer:} That's a bit steep! Can you give me a discount?\\
\textbf{Vendor:} No discount, madam. This is the best price in the whole market.\\
\textbf{Customer:} Come on! I saw the same bag at the next stall for 6,500 CFA francs.\\
\textbf{Vendor:} Ah, but theirs is not genuine leather. OK, because you are my first customer today, 7,000 CFA francs — final price.\\
\textbf{Customer:} That's a good deal. I'll take it. Can I pay by mobile money?\\
\textbf{Vendor:} Of course. Here is your receipt.\\
\textbf{Customer:} Thank you. And if the zip breaks, can I bring it back?\\
\textbf{Vendor:} Yes, you have a one-month guarantee. Just come back with the receipt.
\end{textebox}

\begin{definition}[Glossaire du dialogue]
\begin{itemize}
\item \eng{steep} (adj., familier) : (prix) excessif, exagéré — « That's a bit steep! » = « C'est un peu cher ! »
\item \eng{stall} (n.) : étal, stand (au marché)
\item \eng{genuine leather} : cuir véritable
\item \eng{a good deal} : une bonne affaire
\item \eng{final price} : dernier prix, prix ferme
\item \eng{guarantee} (n.) : garantie
\item \eng{CFA francs} : francs CFA (monnaie utilisée au Cameroun et dans la zone CEMAC)
\end{itemize}
\end{definition}

Observez bien : le client ne dit jamais « I want » ni « Give me ». Il utilise des formules polies (\eng{How much is...?}, \eng{Can you give me...?}, \eng{Can I pay...?}). En anglais, la politesse passe par les questions et les modaux (\eng{can}, \eng{could}), pas seulement par \eng{please}.

\courssub{The people of buying and selling}

\begin{definition}[Customer, consumer, seller... qui est qui ?]
\begin{itemize}
\item a \cle{customer} : a person who buys something from a particular shop or business — un \emph{client} (d'un magasin précis).
\item a \cle{consumer} : anyone who buys and uses goods or services — un \emph{consommateur} (terme général et économique).
\item a \cle{seller} / a \cle{vendor} : a person who sells something — un \emph{vendeur}. \eng{Vendor} est fréquent pour les petits commerçants : \eng{a street vendor} (un vendeur ambulant).
\item a \cle{retailer} : a business that sells goods directly to the public — un \emph{détaillant}.
\item a \cle{wholesaler} : a business that sells goods in large quantities to retailers — un \emph{grossiste}.
\item a \cle{shopkeeper} : a person who owns or manages a small shop — un \emph{commerçant}, un \emph{boutiquier}.
\item a \cle{hawker} : a person who sells goods in the street, moving from place to place — un \emph{colporteur}, un \emph{vendeur à la sauvette} (très visible à Douala et Yaoundé : vendeurs de mouchoirs, de ceintures dans les embouteillages).
\end{itemize}
\end{definition}

\begin{attention}[Customer ou consumer ?]
Ne confondez pas \eng{customer} et \eng{consumer}. Le \eng{customer} est le client d'\emph{un} magasin ou d'\emph{une} entreprise précise : \eng{The customer complained to the shopkeeper.} (Le client s'est plaint au commerçant.) Le \eng{consumer} est une notion économique générale : \eng{Consumers are protected by law.} (Les consommateurs sont protégés par la loi.) On parle de \eng{consumer rights} (droits du consommateur), jamais de « customer rights » dans ce sens abstrait. De même, \eng{customer service} (le service client d'une entreprise) ne se dit pas « consumer service ».
\end{attention}

\courssub{Key actions: bargaining, paying, complaining}

Voici les verbes essentiels du parcours d'achat. Notez leur construction (avec ou sans préposition).

\begin{definition}[Les verbes de l'achat et de la négociation]
\begin{itemize}
\item \cle{to bargain} / \cle{to haggle} : to argue about the price before buying — \emph{marchander}. Construction : \eng{to bargain \textbf{over} something}, \eng{to haggle \textbf{with} somebody \textbf{over/about} the price}.
\item \cle{to negotiate} : to discuss in order to reach an agreement — \emph{négocier} (plus formel que \eng{bargain}).
\item \cle{to ask for a discount} : demander une remise, un rabais.
\item \cle{to reduce} / \cle{to lower the price} : baisser le prix. \eng{The vendor lowered the price from 8,000 to 7,000 CFA francs.}
\item \cle{to refund} : rembourser. Nom : \eng{a refund} (un remboursement).
\item \cle{to exchange} : échanger (un article contre un autre). \eng{Can I exchange this shirt for a larger size?}
\item \cle{to complain (about something)} : se plaindre (de quelque chose). Attention à la préposition \eng{about} !
\end{itemize}
\end{definition}

\begin{exemplebox}[Exemples traduits]
\eng{The shopkeeper refused to refund my money.} = « Le commerçant a refusé de me rembourser. »\par
\eng{This radio has a one-year warranty.} = « Cette radio a une garantie d'un an. »\par
\eng{She haggled with the vendor over the price of tomatoes.} = « Elle a marchandé avec la vendeuse sur le prix des tomates. »\par
\eng{If the shoes are too small, you can exchange them within a week.} = « Si les chaussures sont trop petites, vous pouvez les échanger dans un délai d'une semaine. »\par
\eng{Customers complained about the poor quality of the product.} = « Les clients se sont plaints de la mauvaise qualité du produit. »
\end{exemplebox}

\begin{attention}[Pièges de construction]
\begin{itemize}
\item On dit \eng{to refund \textbf{somebody}} ou \eng{to refund \textbf{the money}}, jamais « to refund \emph{to} somebody » : \eng{They refunded me.} (Ils m'ont remboursé.)
\item \eng{to complain} exige la préposition \eng{about} : \eng{He complained about the service}, et non « he complained the service ».
\item \eng{to bargain} est suivi de \eng{over} : \eng{to bargain over the price}. Ne traduisez pas mot à mot « marchander le prix » par « to bargain the price ».
\item Faux ami : \eng{a discount} = une remise (sur le prix), pas « un escompte » au sens bancaire dans le langage courant ; et \eng{actually} ne veut pas dire « actuellement » (cf. section 5).
\end{itemize}
\end{attention}

\courssub{Places to shop}

\begin{center}
\begin{tabularx}{\linewidth}{|l|X|X|}\hline
\rowcolor{popblueL} \textbf{English} & \textbf{Definition} & \textbf{Français} \\ \hline
a market & an open place where many vendors sell food and goods & un marché \\ \hline
a stall & a small open stand in a market & un étal, un stand \\ \hline
a shop (US: a store) & a small building where goods are sold & un magasin, une boutique \\ \hline
a supermarket & a large self-service shop for food and household goods & un supermarché \\ \hline
a department store & a very large shop divided into departments (clothes, electronics...) & un grand magasin \\ \hline
an online shop & a shop on the internet (Jumia, for example) & une boutique en ligne \\ \hline
\end{tabularx}
\end{center}

\eng{I bought these plantains at a stall in Mvog-Mbi market.} = « J'ai acheté ces plantains à un étal du marché Mvog-Mbi. »\par
\eng{More and more Cameroonians order phones from online shops.} = « De plus en plus de Camerounais commandent des téléphones dans des boutiques en ligne. »

\courssub{Essential collocations and consumer rights}

En anglais, certains verbes vont naturellement avec certains noms : ce sont les \cle{collocations}. Les apprendre par blocs vous fera gagner en naturel.

\begin{center}
\begin{tabularx}{\linewidth}{|X|X|}\hline
\rowcolor{popblueL} \textbf{Collocation} & \textbf{Traduction} \\ \hline
to make a purchase & faire un achat \\ \hline
to place an order & passer une commande \\ \hline
to pay in cash & payer en espèces / en liquide \\ \hline
to pay by mobile money & payer par mobile money (Orange Money, MTN MoMo) \\ \hline
to get / to keep a receipt & obtenir / garder un reçu \\ \hline
to complain about a product & se plaindre d'un produit \\ \hline
to ask for a refund & demander un remboursement \\ \hline
to return a faulty product & retourner un produit défectueux \\ \hline
\end{tabularx}
\end{center}

\begin{definition}[Consumer rights — les droits du consommateur]
\begin{itemize}
\item \cle{consumer rights} : the legal rights of people who buy goods and services — les droits du consommateur.
\item a \cle{warranty} / a \cle{guarantee} : a written promise to repair or replace a product if it breaks within a period of time — une garantie. \eng{The phone comes with a six-month warranty.}
\item a \cle{refund policy} : the rules of a shop about giving money back — la politique de remboursement.
\item \cle{customer service} : the department that helps customers — le service client.
\end{itemize}
\eng{Always keep your receipt: without it, the shop may refuse a refund.} = « Gardez toujours votre reçu : sans lui, le magasin peut refuser le remboursement. »
\end{definition}

\courssub{The buying journey: from shopping around to complaining}

\begin{popfigure}
\begin{tikzpicture}[node distance=0.5cm, every node/.style={font=\small}]
\node[draw=popblue, fill=popblueL, rounded corners, align=center, minimum width=3.4cm] (a) {Shop around\\ (faire le tour des prix)};
\node[draw=popgreen, fill=popgreenL, rounded corners, align=center, minimum width=3.4cm, right=of a] (b) {Compare prices\\ (comparer les prix)};
\node[draw=poporange, fill=poporangeL, rounded corners, align=center, minimum width=3.4cm, right=of b] (c) {Bargain\\ (marchander)};
\node[draw=poppurple, fill=poppurpleL, rounded corners, align=center, minimum width=3.4cm, below=0.7cm of c] (d) {Pay\\ (payer)};
\node[draw=popgold, fill=popgoldL, rounded corners, align=center, minimum width=3.4cm, left=of d] (e) {Keep the receipt\\ (garder le reçu)};
\node[draw=poppink, fill=poppinkL, rounded corners, align=center, minimum width=3.4cm, left=of e] (f) {Complain / Ask for\\ a refund (se plaindre)};
\draw[-{Stealth}, thick, popblue] (a) -- (b);
\draw[-{Stealth}, thick, popgreen] (b) -- (c);
\draw[-{Stealth}, thick, poporange] (c) -- (d);
\draw[-{Stealth}, thick, poppurple] (d) -- (e);
\draw[-{Stealth}, thick, popgold] (e) -- (f);
\end{tikzpicture}
\legende{The buying journey: six steps of a smart consumer.}
\end{popfigure}

\begin{methode}[How to succeed in a market role-play]
Pour un jeu de rôle au marché, structurez toujours votre dialogue en trois temps :
\begin{enumerate}
\item \textbf{Opening} (ouverture polie) : \eng{Good morning! How much do you sell these mangoes?} ou \eng{How much is this dress?}
\item \textbf{Counter-proposal} (contre-proposition) : \eng{That's too much! What about 2,000 CFA francs?} ou \eng{Can you lower the price?} Le vendeur répond : \eng{That's my final price.} ou \eng{OK, because it's you...}
\item \textbf{Conclusion} : \eng{Deal! I'll take it.} (Marché conclu !) ou \eng{No, thanks. I'll think about it.} (Je vais réfléchir.) Terminez par le paiement : \eng{Can I pay by mobile money? Could I have a receipt, please?}
\end{enumerate}
\end{methode}

\begin{exoresolu}[Complétez le dialogue]
Énoncé — Complete the dialogue with: \emph{discount, receipt, steep, refund, warranty, bargain}.\par
A: How much is this fan?\par
B: 25,000 CFA francs.\par
A: That's a bit (1) \ldots ! Can you give me a (2) \ldots ?\par
B: If you (3) \ldots well, I can make it 22,000 CFA francs.\par
A: Deal! Does it have a (4) \ldots ?\par
B: Yes, one year. And if it doesn't work, bring it back with the (5) \ldots and we will (6) \ldots your money.
\tcblower
Solution — (1) \textbf{steep} : le client trouve le prix exagéré, c'est l'adjectif familier pour « trop cher ». (2) \textbf{discount} : on demande une remise avec la collocation \eng{to give a discount}. (3) \textbf{bargain} : le vendeur invite à marchander (\eng{if you bargain well} = « si tu marchandes bien »). (4) \textbf{warranty} : la garantie (ici d'un an). (5) \textbf{receipt} : le reçu, indispensable pour tout retour. (6) \textbf{refund} : rembourser — notez la construction \eng{refund your money} sans préposition.
\end{exoresolu}

\begin{exercice}[Your turn to bargain]
Write a short dialogue (8 to 10 lines) between a customer and a vendor at Marché Central in Douala. The customer wants to buy a school bag. Use at least five of these expressions: \emph{How much is...?, a bit steep, discount, final price, a good deal, pay in cash, receipt, warranty}.
\end{exercice}

\begin{corrige}[Exercice — modèle de production]
\textbf{A:} Good afternoon! How much is this school bag?\par
\textbf{B:} 12,000 CFA francs. Very strong, my friend.\par
\textbf{A:} That's a bit steep! Can you give me a discount?\par
\textbf{B:} How much do you offer?\par
\textbf{A:} What about 9,000 CFA francs?\par
\textbf{B:} No, no. 10,500 CFA francs — final price.\par
\textbf{A:} OK, that's a good deal. I'll pay in cash. Does it have a warranty?\par
\textbf{B:} Three months. Here is your receipt. Keep it well!\par
Ce modèle contient les huit expressions demandées, une ouverture polie, une contre-proposition (\eng{What about 9,000 CFA francs?}) et une conclusion claire (\eng{that's a good deal}).
\end{corrige}

\begin{experience}[Class activity: Mokolo Market in the classroom]
Work in pairs. Student A is a vendor at Mokolo market selling second-hand shoes; Student B is a customer with only 5,000 CFA francs. The vendor starts at 8,000 CFA francs. Bargain for two minutes, then conclude with \eng{Deal!} or \eng{I'll think about it.} Then swap roles. The class votes for the best bargainer.
\end{experience}

\begin{aretenir}[What you must remember]
\begin{itemize}
\item \eng{customer} = client d'un magasin ; \eng{consumer} = consommateur (terme économique) ; \eng{retailer} = détaillant ; \eng{wholesaler} = grossiste ; \eng{hawker} = vendeur ambulant.
\item Actions clés : \eng{to bargain / to haggle (over the price)}, \eng{to ask for a discount}, \eng{to lower the price}, \eng{to refund somebody}, \eng{to exchange a product}, \eng{to complain about something}.
\item Collocations à connaître par cœur : \eng{to make a purchase}, \eng{to place an order}, \eng{to pay in cash / by mobile money}, \eng{to keep the receipt}.
\item Droits : \eng{consumer rights}, \eng{warranty/guarantee}, \eng{refund policy}, \eng{customer service}.
\item Au marché : ouverture polie, contre-proposition (\eng{What about...?}), conclusion (\eng{Deal!}).
\end{itemize}
\end{aretenir}

\coursec{Community resources and public services}

So far, we have compared prices in shops and markets and learned how to bargain for the \emph{best buy}. But a wise consumer does not only buy \emph{products}: every family also depends on \emph{services} that no single person can buy alone — clean water, electricity, health care, schools, roads. These are \cle{community resources} and \cle{public services}. In this section, you will learn the exact English words to name them, describe their quality, and discuss how they are managed — a theme that appears very often in Bac texts about development and citizenship.

\courssub{Observing the vocabulary in context}

\begin{textebox}[A community meeting in Bafoussam]
At last Saturday's meeting at the community hall, the residents of Ndiangdam discussed their neighbourhood's public services. “Our council should provide clean running water,” said Mrs Kamga. “The borehole near the market broke down three months ago, and nobody has repaired it. Meanwhile, we pay our water bills every month!”

Mr Fon, a retired teacher, agreed: “The health centre lacks basic equipment — there is not even a microscope. And look at the library: it is so run-down that students cannot study there during the rainy season.”

The councillor promised to allocate more resources to utilities. “Public spending must benefit everyone,” he said. “We will improve the electricity supply, because the constant power cuts are killing small businesses. A welder cannot work during a blackout.”

A young nurse added that the new subsidized health insurance scheme was affordable for most families, and that the primary school remained free and accessible to all children. “Our amenities are not perfect,” she concluded, “but if we maintain them properly, they give better value for money than expensive private services.”

\textbf{Glossary}: borehole (n.) = forage, puits foré ; councillor (n.) = conseiller municipal ; scheme (n.) = dispositif, programme ; welder (n.) = soudeur ; value for money = rapport qualité-prix.

\textbf{Comprehension questions}:
\begin{enumerate}
\item What three problems do the residents mention?
\item What does the councillor promise to improve, and why?
\item According to the nurse, what advantages do public services have over private ones?
\end{enumerate}
\end{textebox}

\courssub{The lexical field: naming community resources}

\begin{definition}[Amenities]
\cle{Amenities} (always plural in this sense) are \emph{facilities that make a place pleasant and convenient to live in} — commodités, équipements collectifs : water supply, schools, parks, libraries, health centres.

\eng{The new housing estate has excellent amenities: a school, a market and a health centre.} « Le nouveau lotissement dispose d'excellentes commodités : une école, un marché et un centre de santé. »
\end{definition}

\begin{definition}[Key nouns]
\begin{itemize}
\item \cle{public services} : services publics (health, education, transport).
\item \cle{facilities} : équipements, installations (sports facilities, health facilities).
\item \cle{infrastructure} (uncountable) : infrastructure (roads, bridges, networks).
\item \cle{utilities} : services d'utilité publique — eau, électricité, gaz ; \eng{utility bills} = les factures d'eau et d'électricité.
\item \cle{resources} : ressources (money, staff, equipment available to a community).
\end{itemize}
\end{definition}

\begin{definition}[Institutions]
\begin{itemize}
\item \cle{town hall} / \cle{council} : mairie / conseil municipal ; \eng{the local council} = la municipalité.
\item \cle{health centre} : centre de santé ; \cle{public library} : bibliothèque municipale.
\item \cle{community hall} : salle communautaire (des fêtes).
\item \cle{water board} : régie des eaux ; \cle{electricity company} : compagnie d'électricité.
\end{itemize}
\end{definition}

\begin{popfigure}
\begin{tikzpicture}[mindmap, grow cyclic, every node/.style={concept, font=\small},
root concept/.append style={concept color=popblue, font=\small\bfseries},
level 1/.append style={level distance=3.2cm, sibling angle=90},
level 2/.append style={level distance=2.2cm, sibling angle=45}]
\node[root concept]{COMMUNITY RESOURCES}
  child[concept color=popgreen] { node {health} child { node {health centre} } child { node {nurses} } }
  child[concept color=poporange] { node {education} child { node {public school} } child { node {library} } }
  child[concept color=poppurple] { node {utilities} child { node {water} } child { node {electricity} } }
  child[concept color=poppink] { node {leisure} child { node {community hall} } child { node {playground} } };
\end{tikzpicture}
\legende{Carte mentale du champ lexical : les quatre grandes familles de ressources communautaires.}
\end{popfigure}

\courssub{Describing quality: adjectives}

\begin{propriete}[Adjectives for public services]
\begin{itemize}
\item \cle{free} : gratuit — \eng{Primary education is free.} « L'enseignement primaire est gratuit. »
\item \cle{subsidized} : subventionné (prononcé « sub-si-daïzd ») — \eng{subsidized housing} « des logements subventionnés ».
\item \cle{affordable} : abordable (financièrement) — \eng{affordable health care} « des soins abordables ».
\item \cle{accessible} : accessible (physiquement) — \eng{The library is accessible to wheelchair users.} « La bibliothèque est accessible aux personnes en fauteuil roulant. »
\item \cle{available} : disponible — \eng{Running water is not available in every quarter.} « L'eau courante n'est pas disponible dans tous les quartiers. »
\item \cle{efficient} : efficace ; \cle{poor-quality} : de mauvaise qualité ; \cle{run-down} : délabré, dégradé — \eng{a run-down clinic} « un dispensaire délabré ».
\end{itemize}
\end{propriete}

\begin{attention}[Affordable et accessible : deux idées différentes]
Ne confondez pas : \eng{affordable} concerne le \emph{prix} (on peut payer), \eng{accessible} concerne l'\emph{accès} (on peut y aller, l'utiliser). \eng{The service is affordable but not accessible from our village.} « Le service est abordable mais pas accessible depuis notre village. »
\end{attention}

\courssub{Verbs and collocations: managing resources}

\begin{propriete}[Key verbs]
\begin{itemize}
\item \cle{to provide} : fournir (un service) — \eng{The council provides clean water.} Attention : \eng{to provide somebody \emph{with} something}.
\item \cle{to supply} : approvisionner — \eng{The company supplies electricity to the whole town.} « La compagnie approvisionne toute la ville en électricité. »
\item \cle{to maintain} : entretenir — \eng{to maintain roads and boreholes} « entretenir les routes et les forages ».
\item \cle{to improve} : améliorer ; \cle{to lack} : manquer de (sans préposition !) — \eng{The school lacks textbooks.} « L'école manque de manuels. »
\item \cle{to manage resources} : gérer les ressources ; \cle{to allocate resources} : allouer, affecter des ressources.
\end{itemize}
\end{propriete}

\begin{exemplebox}[Essential collocations]
\begin{itemize}
\item \eng{to pay a water bill} « payer une facture d'eau » ; \eng{to pay electricity bills} « payer les factures d'électricité ».
\item \eng{a power cut / a blackout} « une coupure de courant » ; \eng{There was a blackout last night.} « Il y a eu une coupure de courant hier soir. »
\item \eng{running water} « l'eau courante » (au robinet) ; \eng{clean running water} « de l'eau courante potable ».
\item \eng{to allocate resources to health} « affecter des ressources à la santé ».
\item \eng{public spending} « les dépenses publiques » ; \eng{to cut public spending} « réduire les dépenses publiques ».
\item \eng{Our council should provide clean running water.} « Notre commune devrait fournir de l'eau courante potable. »
\item \eng{The health centre lacks basic equipment.} « Le centre de santé manque d'équipement de base. »
\end{itemize}
\end{exemplebox}

\begin{attention}[Faux amis et pièges]
\begin{itemize}
\item \eng{facilities} signifie « équipements, infrastructures », \textbf{pas} « facilités » au sens de facilité. \eng{Sports facilities} = « des installations sportives ». Pour « avec facilité, aisément », dites \eng{with great facility} (singulier, sens abstrait) : \eng{She learns languages with great facility.} « Elle apprend les langues avec une grande aisance. »
\item \eng{to lack} ne prend \textbf{jamais} de préposition : écrivez \eng{The centre lacks staff.}, pas \emph{The centre lacks of staff}.
\item \eng{running water} = eau courante (au robinet) ; ne traduisez pas mot à mot « eau qui court ».
\item \eng{infrastructure} est généralement indénombrable en anglais : \eng{The region's infrastructure is poor.} (pas \emph{infrastructures are}).
\item \eng{scheme} = programme, dispositif (ex. \eng{insurance scheme}) ; ne pas confondre avec \eng{schema} (schéma technique) ni avec le sens péjoratif « combine, arnaque » (informel).
\end{itemize}
\end{attention}

\begin{center}
\begin{tabularx}{\linewidth}{|X|X|X|}\hline
\rowcolor{popblueL} English & Français & Remarque \\ \hline
amenities & commodités, équipements collectifs & toujours au pluriel \\ \hline
utilities & eau, électricité, gaz & \eng{utility bills} = factures \\ \hline
council / town hall & conseil municipal / mairie & \eng{local council} \\ \hline
running water & eau courante & collocation figée \\ \hline
power cut / blackout & coupure de courant & \eng{blackout} = coupure totale \\ \hline
subsidized & subventionné & ≠ \eng{free} (gratuit) \\ \hline
run-down & délabré & pour les bâtiments \\ \hline
to allocate resources & allouer des ressources & \eng{allocate ... to ...} \\ \hline
public spending & dépenses publiques & indénombrable \\ \hline
\end{tabularx}
\end{center}

\courssub{Word families: building your lexicon}

\begin{propriete}[Derivation: verbs, nouns and adjectives]
Apprenez ces familles de mots pour varier votre expression écrite (composition, résumé) et réussir les exercices de dérivation (Bac).
\begin{center}
\begin{tabular}{|l|l|l|l|}\hline
\rowcolor{popblueL} \textbf{Verb} & \textbf{Noun (thing/action)} & \textbf{Noun (person/agent)} & \textbf{Adjective} \\ \hline
provide & provision, provider & provider & provided \\ \hline
maintain & maintenance & maintainer & maintained / maintenance-free \\ \hline
allocate & allocation & allocator & allocated \\ \hline
improve & improvement & — & improved / improving \\ \hline
manage & management & manager & manageable / managerial \\ \hline
consume & consumption & consumer & consumer (adj.) \\ \hline
supply & supply, supplier & supplier & — \\ \hline
access & access, accessibility & — & accessible \\ \hline
afford & affordability & — & affordable \\ \hline
subsidize & subsidy, subsidization & — & subsidized \\ \hline
\end{tabular}
\end{center}
\end{propriete}

\begin{methode}[Comment réussir un exercice de dérivation (Bac)]
\begin{enumerate}
\item Identifiez la catégorie grammaticale manquante (nom, verbe, adjectif, adverbe) grâce à la structure de la phrase (ex. après \eng{the}, \eng{a}, \eng{of} $\rightarrow$ nom ; après \eng{to} $\rightarrow$ verbe ; après \eng{be}, \eng{very} $\rightarrow$ adjectif).
\item Cherchez la racine du mot dans la même famille.
\item Appliquez le suffixe correct (\eng{-tion/-sion}, \eng{-ment}, \eng{-ance/-ence}, \eng{-able/-ible}, \eng{-ity}, \eng{-ness}, \eng{-ly}).
\item Vérifiez l'orthographe (doublement de consonne, chute de \emph{e} muet, \eng{y} $\rightarrow$ \eng{i}).
\item Relisez la phrase complète pour contrôler le sens.
\end{enumerate}
\end{methode}

\courssub{Pronunciation and spelling focus}

\begin{propriete}[Prononciation (repères en lettres françaises) et orthographe]
\begin{itemize}
\item \cle{subsidized} : « sub-si-daïzd » (accent sur 1ère syllabe, \eng{s} = /s/, \eng{i} = /ɪ/).
\item \cle{affordable} : « a-ford-a-beul » (accent sur 2ème syllabe, \eng{a} final = /ə/).
\item \cle{accessible} : « ac-cè-sseu-beul » (accent sur 2ème, double \eng{ss} = /s/, \eng{c} = /s/).
\item \cle{infrastructure} : « in-fra-steuc-teur » (accent sur 1ère, \eng{structure} se entend « steuc-teur »).
\item \cle{bureaucracy} : « bu-ro-cra-sie » (accent sur 2ème, \eng{eu} = /ʊə/).
\item \cle{council} / \cle{counsel} : homophones « kaon-seul » (conseil municipal / avocat).
\item \cle{borehole} : « bor-ole » (deux syllabes, \eng{e} muet final).
\item \cle{programme} (UK spelling) : attention, \eng{program} est l'orthographe US ou informatique ; au Bac, on attend \eng{programme} pour un programme politique, social, TV.
\item \cle{centre} (UK) vs \eng{center} (US) : écrivez \eng{health centre}, \eng{shopping centre}.
\end{itemize}
\end{propriete}

\begin{savaistu}[Utility bills in West Africa]
Au Ghana et au Nigeria, on parle couramment de \eng{utility bills} pour désigner les factures d'eau et d'électricité, et de \eng{utility companies} pour les sociétés qui les fournissent. Ce vocabulaire revient très souvent dans les textes du Bac portant sur le développement et la vie économique : apprenez les paires \eng{water bill / electricity bill}, \eng{water board / electricity company}.
\end{savaistu}

\begin{textebox}[The best buy is a service]
In many Cameroonian communities, the best buy is not a product but a service: a well-maintained borehole gives better value for money than expensive private water vendors. When a village invests in a solar pump, every family saves money every single day. The same is true of a public library: one building, properly managed, offers thousands of books that no family could afford to buy. Wise consumers therefore look beyond the market stall: they ask whether their council maintains the roads, whether the health centre has enough nurses, and whether public spending really improves their daily lives.
\end{textebox}

\begin{exoresolu}[Gap-filling: community services]
Énoncé — Complete the sentences with: \emph{amenities, lacks, running water, power cut, subsidized, allocate, maintain, available}.
\begin{enumerate}
\item The council must \dots\ more resources to the health centre.
\item There was a \dots\ yesterday, so the freezer defrosted.
\item Our quarter has no \dots\ ; we fetch water from the well.
\item The clinic \dots\ basic equipment such as thermometers.
\item The school fees are low because the college is \dots.
\end{enumerate}
\tcblower
Solution détaillée :
\begin{enumerate}
\item \eng{allocate} — collocation \eng{to allocate resources to} « affecter des ressources à ».
\item \eng{power cut} — « une coupure de courant » ; le congélateur a dégelé.
\item \eng{running water} — « eau courante » ; sinon on puise au puits.
\item \eng{lacks} — \eng{to lack} + complément direct, sans préposition : « manque de ».
\item \eng{subsidized} — « subventionné » : l'État paie une partie, donc les frais sont bas.
\end{enumerate}
\end{exoresolu}

\begin{exercice}[Your turn]
Translate into English:
\begin{enumerate}
\item Notre commune devrait fournir de l'eau courante potable à tous les quartiers.
\item Le centre de santé manque de médicaments et le personnel n'est pas disponible le week-end.
\item Les coupures de courant empêchent les petites entreprises de travailler efficacement.
\end{enumerate}
Then write three sentences about the amenities in your own neighbourhood, using \emph{affordable}, \emph{accessible} and \emph{run-down}.
\end{exercice}

\begin{corrige}[Exercice : Your turn]
\begin{enumerate}
\item \eng{Our council should provide clean running water to all quarters.}
\item \eng{The health centre lacks medicines, and the staff is not available at weekends.}
\item \eng{Power cuts prevent small businesses from working efficiently.}
\end{enumerate}
Free production: accept any correct sentence, e.g. \eng{The community hall is affordable but run-down.} / \eng{The public library is accessible to all students.}
\end{corrige}

\begin{aretenir}[Key points]
\begin{itemize}
\item \eng{Amenities} = commodités ; \eng{utilities} = eau et électricité ; \eng{utility bills} = factures.
\item \eng{Affordable} (prix) ≠ \eng{accessible} (accès) ; \eng{subsidized} ≠ \eng{free}.
\item \eng{To lack} + complément direct, jamais \emph{of}.
\item Collocations à connaître par cœur : \eng{running water}, \eng{power cut}, \eng{to pay a water bill}, \eng{to allocate resources}, \eng{public spending}.
\item Dérivations utiles : \eng{maintain/maintenance}, \eng{provide/provision}, \eng{allocate/allocation}, \eng{improve/improvement}, \eng{manage/management}, \eng{access/accessibility}, \eng{afford/affordability}.
\item Orthographe UK : \eng{programme}, \eng{centre}, \eng{council}, \eng{subsidized} (z autorisé en UK, s aussi possible \eng{subsidised}).
\end{itemize}
\end{aretenir}

\coursec{Word families, false friends and pronunciation}

In the previous section, we explored the vocabulary of \emph{community resources and public services} — words like \eng{welfare}, \eng{facilities} and \eng{public transport}. You may have noticed that many of these words are built from the same roots: \eng{consume} gives \eng{consumer}, \eng{consumption}, \eng{consumerism}... This is exactly what we are going to study now. English, like French, builds whole \emph{families of words} from one root. Mastering these families will multiply your vocabulary effortlessly — but beware: some words look French and are traps. That is the world of \emph{false friends}.

\courssub{Observing word families in a text}

\begin{textebox}[An article from “The Cameroon Consumer” — original text]
Yaoundé — As the \cle{economy} slows down, Cameroonian \cle{consumers} are changing their habits. According to a local \cle{economist}, household \cle{consumption} has fallen in recent months, and families now compare prices carefully before buying. “\cle{Consumerism} pushed people to buy more than they needed,” she explains. “Today, shoppers look for \cle{economical} products — a stove that uses less gas, a phone that lasts longer.”

At Mokolo market, trader Mrs Etoundi agrees: “Customers bargain hard. They want \cle{value} for money. A cheap but \cle{valueless} product is a bad buy; a slightly pricey but \cle{valuable} one is often the best choice.” She adds that good after-sales service is \cle{invaluable}: “When a customer trusts you, he comes back.”

Economists warn that the current \cle{economic} situation — rising prices and falling incomes — will continue to shape buying behaviour. Meanwhile, consumer associations advise families to read guarantees carefully and to avoid goods of no real value, however attractive their packaging may be.
\end{textebox}

\textbf{Glossaire}

\begin{center}
\begin{tabularx}{\linewidth}{|l|l|X|}
\hline
\rowcolor{popblueL} English word & Nature & Français \\
\hline
consumer & noun & consommateur \\
\hline
consumption & noun & consommation \\
\hline
consumerism & noun & consumérisme (société de consommation) \\
\hline
economical & adjective & économe, qui fait économiser \\
\hline
valueless & adjective & sans valeur \\
\hline
invaluable & adjective & inestimable, très précieux \\
\hline
pricey & adjective (informal) & cher, coûteux (familier) \\
\hline
\end{tabularx}
\end{center}

\textbf{Comprehension questions}
\begin{enumerate}
\item According to the economist, how have consumers' habits changed?
\item Why does Mrs Etoundi say a cheap product can be a bad buy?
\item What does she consider “invaluable”, and why?
\end{enumerate}

\courssub{How word families are built: the suffix system}

Observe the text again: \eng{consume} (verb) → \eng{consumer} (the person) → \eng{consumption} (the action) → \eng{consumerism} (the social trend). English builds these words with \emph{suffixes}.

\begin{propriete}[Forming word families with suffixes]
\begin{itemize}
\item \textbf{-er / -or} forms the \emph{agent} (the person or thing that does the action): \eng{consume} → \eng{consumer} (consommateur), \eng{buy} → \eng{buyer}, \eng{sell} → \eng{seller}, \eng{trade} → \eng{trader}.
\item \textbf{-tion / -sion} forms the \emph{noun of action}: \eng{consume} → \eng{consumption} (consommation), \eng{produce} → \eng{production}, \eng{decide} → \eng{decision}.
\item \textbf{-ism} forms the name of a \emph{system, doctrine or trend}: \eng{consumer} → \eng{consumerism} (consumérisme), \eng{capital} → \eng{capitalism}.
\item \textbf{-ist} forms the \emph{specialist}: \eng{economy} → \eng{economist} (économiste).
\item \textbf{-ic / -ical} form \emph{adjectives}: \eng{economy} → \eng{economic} / \eng{economical}.
\item \textbf{-less} means \emph{without}: \eng{value} → \eng{valueless} (sans valeur); \textbf{-able} means \emph{that can be / worthy of}: \eng{value} → \eng{valuable} (de grande valeur).
\item \textbf{-ics} forms the name of a \emph{science or discipline}: \eng{economy} → \eng{economics} (les sciences économiques).
\end{itemize}
Note the spelling changes: a final silent \emph{-e} disappears before a suffix beginning with a vowel (\eng{consume} → \eng{consumer}, \eng{consumption}).
\end{propriete}

\begin{exemplebox}[The family of “consume” in sentences]
\eng{Consumption increases during festive seasons.} « La consommation augmente pendant les fêtes. »

\eng{Consumers should always ask for a receipt.} « Les consommateurs devraient toujours demander un reçu. »

\eng{Consumerism encourages people to buy things they do not need.} « Le consumérisme pousse les gens à acheter des choses dont ils n'ont pas besoin. »

\eng{We must reduce our consumption of imported goods.} « Nous devons réduire notre consommation de produits importés. »
\end{exemplebox}

\courssub{The family of ECONOMY: a mind map}

The root \eng{economy} (pronounced « i-ko-no-mi ») gives \emph{five} very common words. Study this tree carefully:

\begin{popfigure}
\begin{center}
\begin{tikzpicture}[
  grow=right,
  level 1/.style={sibling distance=2.1cm, level distance=4.6cm},
  every node/.style={draw, rounded corners, align=center, font=\small, inner sep=4pt},
  edge from parent/.style={draw=popmuted, thick}
]
\node[fill=poporangeL, font=\small\bfseries, align=center]{ECONOMY\\ « i-ko-no-mi »\\ l'économie}
  child { node[fill=popblueL, align=center]{economics\\ la science économique\\ \emph{Economics is my favourite subject.}} }
  child { node[fill=popgreenL, align=center]{economist\\ l'économiste\\ \emph{The economist gave a lecture.}} }
  child { node[fill=poppinkL, align=center]{economical\\ économe, économique (qui épargne)\\ \emph{an economical car}} }
  child { node[fill=poppurpleL, align=center]{economic\\ économique (relatif à l'économie)\\ \emph{economic growth}} };
\end{tikzpicture}
\end{center}
\legende{L'arbre lexical de la famille « ECONOMY » : une racine, quatre dérivés.}
\end{popfigure}

\begin{attention}[Economic or economical? Never mix them up!]
This is one of the most frequent errors of francophone learners, because French uses \emph{économique} for both ideas.
\begin{itemize}
\item \textbf{economic} = related to the economy of a country or region: \eng{an economic crisis} (une crise économique), \eng{economic growth} (la croissance économique), \eng{economic policy}.
\item \textbf{economical} = that saves money or resources: \eng{an economical heater} (un chauffage économique), \eng{an economical car} (une voiture qui consomme peu).
\end{itemize}
\eng{This stove is very economical.} « Ce réchaud est très économe. » — Never say \eng{an economic stove}!
\end{attention}

\courssub{The family of VALUE and the family of PRICE}

\begin{center}
\begin{tabularx}{\linewidth}{|l|l|X|}
\hline
\rowcolor{popblueL} Word & Nature & Meaning / Example \\
\hline
value & noun/verb & la valeur / évaluer — \eng{good value for money} (un bon rapport qualité-prix) \\
\hline
valuable & adjective & de grande valeur, précieux — \eng{valuable advice} \\
\hline
valueless & adjective & sans valeur — \eng{a valueless coin} \\
\hline
invaluable & adjective & inestimable, extrêmement précieux — \eng{invaluable help} \\
\hline
price & noun & le prix — \eng{Prices are rising.} \\
\hline
priceless & adjective & inestimable, sans prix — \eng{a priceless work of art} \\
\hline
pricey & adjective (informal) & cher (familier) — \eng{This restaurant is a bit pricey.} \\
\hline
\end{tabularx}
\end{center}

\begin{savaistu}[“Invaluable” — a classic Bac trap!]
Logically, \eng{invaluable} should mean “without value”... but it means exactly the opposite: \emph{inestimable, extremely precious}! \eng{Your help was invaluable.} « Ton aide a été inestimable. » In reading comprehension at the Baccalauréat, many candidates translate it as \emph{sans valeur} and misunderstand the whole passage. Remember: \eng{valueless} = sans valeur; \eng{invaluable} = très précieux. The same surprise exists with \eng{priceless}, which also means \emph{inestimable}, not “free”!
\end{savaistu}

\courssub{False friends: words that look French but are not}

\begin{definition}[What is a false friend?]
A \emph{false friend} (\eng{faux ami}) is an English word that resembles a French word but has a different meaning. In the field of money and shopping, they are numerous and dangerous.
\end{definition}

\begin{center}
\begin{tabularx}{\linewidth}{|X|X|X|}
\hline
\rowcolor{popblueL} English word & Real meaning & French look-alike (trap) \\
\hline
to demand & exiger (fortement) & ≠ demander poliment = \eng{to ask for} \\
\hline
commerce & le commerce (activité, secteur) & ≠ un magasin = \eng{a shop / a store} \\
\hline
a magazine & un magazine, une revue & ≠ un magasin = \eng{a shop / a store} \\
\hline
actually & en fait, en réalité & ≠ actuellement = \eng{currently, at present} \\
\hline
to achieve & accomplir, réaliser & ≠ acheter = \eng{to buy} \\
\hline
a coin & une pièce de monnaie & ≠ un coin = \eng{a corner} \\
\hline
\end{tabularx}
\end{center}

\begin{exemplebox}[False friends in context]
\eng{The workers are demanding higher wages.} « Les travailleurs exigent des salaires plus élevés. » (politely: \eng{They are asking for higher wages.})

\eng{Actually, this product is cheaper at the supermarket.} « En fait, ce produit est moins cher au supermarché. » (\eng{currently}: \eng{I am currently comparing prices.} « Je compare les prix actuellement. »)

\eng{I bought a magazine at the shop near the market.} « J'ai acheté un magazine dans la boutique près du marché. »
\end{exemplebox}

\begin{attention}[Calques du français à éviter]
\begin{itemize}
\item Do not say \eng{the price is costly}: a \emph{product} is expensive or cheap, but a \emph{price} is \textbf{high} or \textbf{low}. \eng{The price of fuel is high.} « Le prix du carburant est élevé. »
\item Do not say \eng{to buy cheaply}: say \eng{to buy cheap} or \eng{to buy at a low price}. \eng{She bought her phone cheap.} « Elle a acheté son téléphone à bas prix. »
\item \eng{economic} ≠ \eng{economical} (see above).
\end{itemize}
\end{attention}

\begin{methode}[How to avoid false friends]
\begin{enumerate}
\item When an English word looks like a French word, be suspicious: resemblance is a warning, not a guarantee.
\item Before using it, mentally check it with \emph{one example sentence learnt by heart}: e.g. \eng{Actually, I disagree.} = « En fait, je ne suis pas d'accord. »
\item If the sentence does not work with your intended meaning, look for the real equivalent (\eng{currently}, \eng{to ask for}, \eng{a shop}).
\item Keep a personal “false friends notebook” and add two examples per word.
\end{enumerate}
\end{methode}

\courssub{Pronunciation and spelling of key words}

English spelling does not always show pronunciation. Learn these key words with their pronunciation in French letters:

\begin{center}
\begin{tabularx}{\linewidth}{|l|l|X|}
\hline
\rowcolor{popblueL} Word & Pronunciation & Remark \\
\hline
consumer & « kon-sio-meur » & stress on « sio » (2nd syllable) \\
\hline
consumption & « kon-somp-cheune » & stress on « somp » (2nd syllable) \\
\hline
economy & « i-ko-no-mi » & stress on « ko » (2nd syllable) \\
\hline
economic & « i-ko-no-mik » & stress on « no » (3rd syllable) \\
\hline
economical & « i-ko-no-mi-kel » & stress on « no » (3rd syllable) \\
\hline
economist & « i-ko-no-mist » & stress on « ko » (2nd syllable) \\
\hline
bargain & « ba-guine » & stress on 1st syllable « ba »; final \emph{-ain} sounds like « ine » \\
\hline
guarantee & « ga-ran-ti » & silent \emph{u}; stress on the last syllable \\
\hline
valuable & « va-liou-beul » & only three syllables \\
\hline
cheap & « tchipe » & beware: \eng{sheep} « chip » (mouton)! \\
\hline
\end{tabularx}
\end{center}

Note how the \emph{stress moves} inside the family: e-\textbf{co}-no-my but e-co-\textbf{no}-mic. This stress shift is typical of English word families and is often tested in oral exams.

\courssub{Practice: solved exercise}

\begin{exoresolu}[Word families and false friends]
\textbf{A. Complete each sentence with the correct word from the family in brackets.}

1. \eng{Household \dots\ has decreased this year.} (consume)
2. \eng{An \dots\ from the university explained the crisis on television.} (economy)
3. \eng{This gas cooker is very \dots : it uses little gas.} (economy)
4. \eng{The government has announced new \dots\ measures.} (economy)
5. \eng{Her advice was \dots\ — it saved our business.} (value)

\textbf{B. Correct the mistake in each sentence.}

6. \eng{Actually, prices are rising in Douala.} (meaning: at present)
7. \eng{The price of bread is very costly.}
8. \eng{I bought some fruit at the magazine.}
\tcblower
\textbf{Solution détaillée}

\textbf{A.} 1. \eng{consumption} (nom d'action, suffixe -tion). 2. \eng{economist} (le spécialiste, suffixe -ist). 3. \eng{economical} (qui fait économiser — attention, pas \eng{economic}!). 4. \eng{economic} (relatif à la politique de l'économie du pays). 5. \eng{invaluable} (inestimable — \eng{valuable} serait possible mais plus faible ; jamais \eng{valueless}, qui veut dire « sans valeur »).

\textbf{B.} 6. \eng{Currently, prices are rising in Douala.} — \eng{actually} signifie « en fait », pas « actuellement ». 7. \eng{The price of bread is very high.} — un prix est \eng{high} ou \eng{low}, jamais \eng{costly}. 8. \eng{I bought some fruit at the shop.} — \eng{a magazine} est une revue ; un magasin se dit \eng{a shop} ou \eng{a store}.
\end{exoresolu}

\begin{exercice}[Your turn]
Translate into English, paying attention to word families and false friends:
\begin{enumerate}
\item « Les consommateurs exigent des prix plus bas. »
\item « Cette voiture est très économique (elle consomme peu). »
\item « En fait, ce magasin est fermé actuellement. »
\item « Son expérience est inestimable pour notre association de consommateurs. »
\end{enumerate}
\end{exercice}

\begin{corrige}[Exercice « Your turn »]
1. \eng{Consumers are demanding lower prices.} — \eng{to demand} convient car il s'agit d'une exigence forte.
2. \eng{This car is very economical.} — jamais \eng{economic} pour un objet qui fait économiser.
3. \eng{Actually, this shop is currently closed.} — les deux faux amis dans une seule phrase : \eng{actually} = en fait ; \eng{currently} = actuellement ; \eng{shop} = magasin.
4. \eng{Her experience is invaluable to our consumers' association.} — \eng{invaluable} = inestimable, et non « sans valeur ».
\end{corrige}

\begin{aretenir}[Key points of this section]
\begin{itemize}
\item One root, many words: \eng{consume} → \eng{consumer} → \eng{consumption} → \eng{consumerism}; \eng{economy} → \eng{economic} / \eng{economical} / \eng{economist} / \eng{economics}.
\item \eng{economic} = de l'économie du pays ; \eng{economical} = qui fait économiser.
\item \eng{invaluable} and \eng{priceless} mean \emph{inestimable}; \eng{valueless} means \emph{sans valeur}.
\item Beware of false friends: \eng{to demand} (exiger), \eng{actually} (en fait), \eng{a magazine} (une revue).
\item A price is \eng{high} or \eng{low}, never \eng{costly}.
\end{itemize}
\end{aretenir}

\coursec{Using the lexis in context}

Dans la section précédente, nous avons travaillé sur les familles de mots (\eng{cheap – cheaper – cheapness}, \eng{value – valuable – valueless}), les faux amis (\eng{actually} ne veut pas dire « actuellement ») et la prononciation des mots-clés de la leçon. Il est maintenant temps de \cle{réemployer tout ce lexique en contexte} : lire un texte authentique, repérer les mots cibles, compléter des phrases, négocier au marché et rédiger un court paragraphe d'opinion. C'est exactement ce que l'on vous demandera au Baccalauréat : non pas réciter des listes, mais \emph{utiliser} le vocabulaire avec précision.

\courssub{Reading: “Smart Shopping in Yaoundé”}

Lisez le texte suivant à voix haute, puis repérez les \textbf{10 mots cibles} de la leçon qui y apparaissent (ils sont en gras dans la version corrigée du professeur).

\begin{textebox}[Smart Shopping in Yaoundé — an original article]
Before the new school year, Mrs Ngono shops around in three markets in Yaoundé: Mokolo, Mvog-Mbi and the Central Market. She is a careful consumer. She compares prices from stall to stall, checks the quality of each item and bargains politely with the vendors. “How much is this bag, madam?” she asks. “Eight thousand francs,” replies the vendor. “That is overpriced! My neighbour bought the same one for five thousand. Can you give me a discount?” After a few minutes of friendly negotiation, they agree on a fair price.

Mrs Ngono never buys the first thing she sees. Last year, she bought a cheap school bag at a roadside shop. It looked attractive, but the zip broke after two weeks. It was poor value for money. This year, she prefers affordable but durable bags: strong zips, solid handles and waterproof material. They cost a little more, but they are reliable and will last the whole school year. In the end, she buys two good-quality bags for her children: a real bargain and true value for money.

Mrs Ngono also thinks about her community. She buys from local traders because their shops create jobs, and she uses public services wisely: she pays her water and electricity bills on time and takes the shared taxi instead of a private car to save money. Her advice to young shoppers? “Compare before you pay. A low price is not always a good deal, and a high price is not always a guarantee of quality. The best buy is the product that gives you the most for your money.”
\end{textebox}

\textbf{Glossaire des mots difficiles}

\begin{center}
\begin{tabularx}{\linewidth}{|l|l|X|}
\hline
\rowcolor{popblueL} Mot anglais & Nature & Traduction française \\
\hline
to shop around & verbe & comparer les prix dans plusieurs magasins \\
\hline
stall & nom & étal, stand (au marché) \\
\hline
to bargain & verbe & marchander, négocier le prix \\
\hline
vendor & nom & vendeur, vendeuse \\
\hline
overpriced & adjectif & trop cher, surfait \\
\hline
discount & nom & remise, réduction \\
\hline
durable & adjectif & durable, résistant \\
\hline
reliable & adjectif & fiable, de confiance \\
\hline
a bargain & nom & une bonne affaire \\
\hline
value for money & expression & rapport qualité-prix \\
\hline
bill & nom & facture (eau, électricité) \\
\hline
a good deal & expression & une bonne affaire, un bon compromis \\
\hline
\end{tabularx}
\end{center}

\begin{exercice}[Comprehension questions]
Répondez en phrases complètes, en anglais.
\begin{enumerate}
\item Why does Mrs Ngono visit three different markets?
\item What happened to the cheap bag she bought last year?
\item What qualities does she look for in a school bag this year?
\item How does Mrs Ngono help her community when she shops?
\item In your own words, what is “the best buy” according to Mrs Ngono?
\end{enumerate}
\end{exercice}

\begin{corrige}[Comprehension questions]
\begin{enumerate}
\item \eng{She visits three markets to shop around and compare prices before buying.} (Elle visite trois marchés pour comparer les prix avant d'acheter.)
\item \eng{The zip broke after two weeks; it was poor value for money.} (La fermeture a cassé au bout de deux semaines ; mauvais rapport qualité-prix.)
\item \eng{She looks for durable and reliable bags with strong zips, solid handles and waterproof material.} (Elle cherche des sacs durables et fiables.)
\item \eng{She buys from local traders, pays her bills on time and uses shared taxis.} (Elle achète chez les commerçants locaux, paie ses factures à temps et prend les taxis collectifs.)
\item \eng{The best buy is the product that gives you the most for your money.} (Le meilleur achat est le produit qui en donne le plus pour votre argent.)
\end{enumerate}
\end{corrige}

\courssub{Spotting the target lexis in the text}

\begin{exoresolu}[Find the 10 target words]
Énoncé : relevez dans le texte 10 mots ou expressions de la leçon (prix, qualité, marchandage, communauté) et classez-les.
\tcblower
Solution : voici les 10 mots cibles, classés par champ lexical.
\begin{center}
\begin{tabular}{|l|l|}
\hline
\rowcolor{popblueL} Champ lexical & Mots du texte \\
\hline
Prices & \textbf{compares prices}, \textbf{overpriced}, \textbf{discount} \\
\hline
Quality & \textbf{durable}, \textbf{reliable}, \textbf{value for money} \\
\hline
Shopping & \textbf{shops around}, \textbf{bargains}, \textbf{a bargain} \\
\hline
Community & \textbf{public services} \\
\hline
\end{tabular}
\end{center}
Remarquez la différence entre \eng{a bargain} (nom : une bonne affaire) et \eng{to bargain} (verbe : marchander). Même famille, deux natures différentes : \emph{She bargains with the vendor and finds a bargain.}
\end{exoresolu}

\begin{popfigure}
\begin{center}
\begin{tikzpicture}
\node[draw=popdark, fill=popblueL, rounded corners, font=\bfseries, minimum width=2.6cm] (t) at (0,0) {BEST BUY};
\node[draw=popblue, fill=popblueL!60, rounded corners, align=center, minimum width=2.9cm] (p) at (-4.6,2.2) {\textbf{Prices}\\ compare prices\\ overpriced\\ discount};
\node[draw=popgreen, fill=popgreenL, rounded corners, align=center, minimum width=2.9cm] (q) at (-1.55,2.2) {\textbf{Quality}\\ durable\\ reliable\\ value for money};
\node[draw=poporange, fill=poporangeL, rounded corners, align=center, minimum width=2.9cm] (s) at (1.55,2.2) {\textbf{Shopping}\\ shop around\\ bargain\\ a good deal};
\node[draw=poppurple, fill=poppurpleL, rounded corners, align=center, minimum width=2.9cm] (c) at (4.6,2.2) {\textbf{Community}\\ local traders\\ public services\\ bills};
\draw[-{Stealth}, popblue, thick] (p) -- (t);
\draw[-{Stealth}, popgreen, thick] (q) -- (t);
\draw[-{Stealth}, poporange, thick] (s) -- (t);
\draw[-{Stealth}, poppurple, thick] (c) -- (t);
\end{tikzpicture}
\end{center}
\legende{Carte de synthèse : quatre champs lexicaux, trois mots-clés chacun, pour mémoriser le vocabulaire du « best buy ».}
\end{popfigure}

\courssub{Gap-fill: activating the vocabulary}

\begin{exercice}[Complete the sentences]
Complétez avec : \eng{affordable}, \eng{overpriced}, \eng{bargain}, \eng{reliable}, \eng{value for money}, \eng{discount}, \eng{shop around}, \eng{durable}.
\begin{enumerate}
\item This restaurant is \dots\dots\dots : a simple meal costs 10,000 francs!
\item Before you buy a phone, \dots\dots\dots in two or three shops.
\item The vendor gave me a 10 \% \dots\dots\dots because I bought three shirts.
\item This radio is cheap and \dots\dots\dots ; it has worked for ten years.
\item Leather shoes are more \dots\dots\dots than plastic ones.
\item At 500 francs, these mangoes are a real \dots\dots\dots .
\item The bus is the most \dots\dots\dots way to travel from Buea to Douala.
\item This computer is expensive, but it offers excellent \dots\dots\dots .
\end{enumerate}
\end{exercice}

\begin{corrige}[Complete the sentences]
1. \eng{overpriced} (trop cher) — 2. \eng{shop around} (comparez les prix) — 3. \eng{discount} (remise) — 4. \eng{reliable} (fiable) — 5. \eng{durable} (durables) — 6. \eng{bargain} (bonne affaire) — 7. \eng{affordable} (abordable) — 8. \eng{value for money} (rapport qualité-prix).
\end{corrige}

\courssub{Building personal sentences: comparing real products}

Observez d'abord ces exemples, puis produisez les vôtres.

\begin{exemplebox}[Comparing two real products]
\eng{After comparing both shops, I found that the second one offered better value for money.} « Après avoir comparé les deux magasins, j'ai trouvé que le second offrait un meilleur rapport qualité-prix. »

\eng{Local rice is cheaper than imported rice, and it is just as tasty.} « Le riz local est moins cher que le riz importé, et il est tout aussi savoureux. »

\eng{This smartphone is more expensive, but its battery is more reliable.} « Ce smartphone est plus cher, mais sa batterie est plus fiable. »

\eng{A shared taxi is more affordable than a private car, so it is the best buy for students.} « Un taxi collectif est plus abordable qu'une voiture privée, c'est donc le meilleur choix pour les élèves. »
\end{exemplebox}

\begin{propriete}[Les structures de comparaison à réemployer]
\begin{itemize}
\item \cle{cheaper than} / \cle{more expensive than} : \eng{Rice is cheaper than spaghetti.}
\item \cle{more affordable than} : \eng{The bus is more affordable than the taxi.}
\item \cle{as ... as} (égalité) : \eng{Local rice is as good as imported rice.}
\item \cle{the best / the worst} + nom : \eng{It is the best buy in the market.}
\item Conclusion type : \eng{X offers (better) value for money because...}
\end{itemize}
Attention à l'ordre des mots : l'adjectif précède le nom en anglais (\eng{a durable bag}, jamais \eng{a bag durable} dans une description simple).
\end{propriete}

\begin{exercice}[Your turn]
Écrivez trois phrases personnelles comparant : (a) deux marques de riz, (b) deux téléphones, (c) deux moyens de transport à Yaoundé ou Douala. Utilisez au moins trois structures de la boîte ci-dessus.
\end{exercice}

\courssub{Guided mini-dialogue: at the market}

Complétez ce dialogue avec votre partenaire, puis jouez-le. Objectif : négocier et conclure sur le \eng{best buy}.

\begin{experience}[Role-play: negotiating at Mokolo market]
\begin{itemize}
\item \textbf{Customer:} Good morning! How much is this school bag?
\item \textbf{Vendor:} It is 7,000 francs, madam.
\item \textbf{Customer:} 7,000? That is \dots\dots\dots (overpriced)! I saw the same bag for 5,000 at another stall. Can you give me a \dots\dots\dots (discount)?
\item \textbf{Vendor:} For you, madam, 6,000. Look at the quality: it is very \dots\dots\dots (durable).
\item \textbf{Customer:} Hmm. Is it \dots\dots\dots (reliable)? The zip of my last bag broke quickly.
\item \textbf{Vendor:} This one has a strong zip and solid handles. It is real \dots\dots\dots (value for money).
\item \textbf{Customer:} All right. At 5,500, it is a \dots\dots\dots (bargain). I will take it. It is the \dots\dots\dots (best buy) of my day!
\end{itemize}
Variante : inversez les rôles et vendez un téléphone ou un sac de riz.
\end{experience}

\courssub{Writing: “The best buy in my community”}

\begin{methode}[How to write a comparison paragraph (5–6 lines)]
\begin{enumerate}
\item \textbf{Introduce the two products} : \eng{In my community, people often hesitate between X and Y.}
\item \textbf{Compare the price} : \eng{X is cheaper than Y / more affordable than Y.}
\item \textbf{Compare the quality} : \eng{However, Y is more durable and more reliable.}
\item \textbf{Conclude} : \eng{The best buy is ... because it offers real value for money.}
\end{enumerate}
Reliez vos idées avec \eng{however}, \eng{because}, \eng{in the end}, \eng{that is why}.
\end{methode}

\begin{exemplebox}[Model paragraph]
\eng{In my community, many students hesitate between a bicycle and a bendskin to go to school. A bicycle is expensive at first, but it is durable and costs nothing to use. A bendskin is faster, but the daily fares add up quickly. In the end, the best buy is the bicycle, because it offers real value for money and keeps students healthy.} « Dans ma communauté, beaucoup d'élèves hésitent entre un vélo et un bendskin pour aller à l'école. Le vélo coûte cher au départ, mais il est durable et ne coûte rien à l'usage. Le bendskin est plus rapide, mais les courses quotidiennes s'accumulent vite. En fin de compte, le meilleur achat est le vélo, car il offre un vrai rapport qualité-prix et garde les élèves en bonne santé. »
\end{exemplebox}

\begin{exercice}[Your paragraph]
Rédigez votre propre paragraphe \eng{The best buy in my community} (5 à 6 lignes) en suivant la méthode en quatre étapes. Choisissez deux produits ou services réels (riz local / riz importé, taxi / bus, marché / supermarché, eau de source / eau en sachet...).
\end{exercice}

\begin{attention}[Au Bac : bannissez “good” et “bad” !]
Les correcteurs pénalisent les copies qui répètent \eng{good} et \eng{bad}. Remplacez-les par le lexique précis de la leçon :

\begin{center}
\begin{tabularx}{\linewidth}{|X|X|}
\hline
\rowcolor{popblueL} Au lieu de... & Écrivez plutôt... \\
\hline
a good price & an affordable price, a fair price, a bargain \\
\hline
a bad price & an overpriced item, a rip-off \\
\hline
good quality & durable, reliable, good value for money \\
\hline
bad quality & poor quality, fragile, not worth the price \\
\hline
to look in many shops & to shop around, to compare prices \\
\hline
to discuss the price & to bargain, to negotiate, to ask for a discount \\
\hline
\end{tabularx}
\end{center}

Autre piège : ne traduisez pas « rapport qualité-prix » par \eng{quality-price report} ! L'expression correcte est \cle{value for money}. De même, « faire les magasins » ne se dit pas \eng{to make the shops} mais \eng{to shop around}.
\end{attention}

\begin{aretenir}[L'essentiel de la section]
\begin{itemize}
\item Lire un texte en repérant les mots cibles permet de mémoriser le lexique en contexte.
\item Les structures de comparaison (\eng{cheaper than}, \eng{as ... as}, \eng{the best buy}) s'emploient avec le vocabulaire précis : \eng{affordable}, \eng{durable}, \eng{reliable}, \eng{overpriced}, \eng{a bargain}, \eng{value for money}.
\item Un paragraphe de comparaison suit quatre étapes : introduction, prix, qualité, conclusion.
\item Au marché comme à l'écrit, concluez toujours en justifiant votre choix : \eng{... because it offers real value for money.}
\end{itemize}
\end{aretenir}

\newpage

\coursec{Activité d'intégration}

\coursec{Lesson 15 — Vocabulary: Comparing Prices and Quality for Best Buys}

\courssub{Objectifs de la leçon}

À l'issue de cette leçon, l'élève sera capable de :

\begin{itemize}
\item \cle{maîtriser} le vocabulaire thématique des prix, de la qualité, de la négociation et des services publics (\eng{cheaper than}, \eng{value for money}, \eng{durable}, \eng{a rip-off}, \eng{a bargain}, \eng{consumer rights}, \eng{standard bureau}...) ;
\item \cle{structurer} une comparaison à l'oral et à l'écrit en utilisant correctement les comparatifs et superlatifs (réguliers et irréguliers) ainsi que les connecteurs logiques (\eng{however}, \eng{although}, \eng{whereas}) ;
\item \cle{réemployer} le lexique dans une tâche de communication authentique : analyser des offres commerciales, rédiger un rapport de recommandation (\eng{best buy}) et défendre son choix en débat ;
\item \cle{identifier} et \cle{éviter} les faux amis fréquents et les calques du français (\eng{actually} vs \eng{currently}, \eng{economic} vs \eng{economical}, \eng{expensive} vs \eng{expensif}, \eng{payment delay} vs \eng{credit period}) ;
\item \cle{connaître} les principales ressources communautaires de protection du consommateur au Cameroun (MINCOMMERCE, ANOR, CAPHYA).
\end{itemize}

\courssub{Talking about Prices and Money}

\begin{definition}[Vocabulaire de base : prix et argent]
\begin{itemize}
\item \eng{Price} /praɪs/ (n.) : le prix. \eng{Cost} /kɒst/ (n.) : le coût (dépense pour produire/acheter). \eng{Value} /ˈvæljuː/ (n.) : la valeur.
\item \eng{Cheap} /tʃiːp/ (adj.) : bon marché, peu cher. \eng{Expensive} /ɪkˈspensɪv/ (adj.) : cher.
\item \eng{Afford} /əˈfɔːd/ (v.) : avoir les moyens de s'acheter. \textit{Structure :} \eng{can/cannot afford to + V} / \eng{can/cannot afford + N}.
\item \eng{Charge} /tʃɑːdʒ/ (v.) : faire payer, demander comme prix. \eng{The shop charges 12,000 FCFA.}
\item \eng{Pay} /peɪ/ (v., irr. \eng{paid}, \eng{paid}) : payer. \eng{Pay for} (payer pour), \eng{pay in cash} (payer en espèces), \eng{pay by card} (payer par carte).
\item \eng{Discount} /ˈdɪskaʊnt/ (n.) : remise, réduction. \eng{A 5\% discount}. \eng{Give a discount on}.
\item \eng{Bargain} /ˈbɑːɡɪn/ (n.) : une bonne affaire. (v.) : marchander. \eng{It's a real bargain!}
\item \eng{Rip-off} /ˈrɪp ɒf/ (n., familier) : une arnaque, un vol (prix excessif pour qualité médiocre).
\item \eng{Value for money} /ˈvæljuː fɔː ˈmʌni/ (n. phrase) : rapport qualité-prix.
\item \eng{Receipt} /rɪˈsiːt/ (n.) : reçu (le \emph{p} ne se prononce pas). \eng{Invoice} /ˈɪnvɔɪs/ (n.) : facture.
\item \eng{Credit period} /ˈkredɪt ˈpɪəriəd/ (n.) : délai de paiement (attention : \eng{payment delay} = retard de paiement).
\item \eng{Delivery} /dɪˈlɪvəri/ (n.) : livraison. \eng{Free delivery}.
\end{itemize}
\end{definition}

\begin{exemplebox}[Collocations utiles]
\begin{itemize}
\item \eng{To compare prices} — \eng{to shop around} (faire plusieurs magasins).
\item \eng{To fit one's budget} — \eng{to be within one's budget} (rentrer dans le budget).
\item \eng{To be overpriced} / \eng{to be reasonably priced} (surévalué / à prix raisonnable).
\item \eng{To get a refund} / \eng{to ask for a refund} (se faire rembourser / demander un remboursement).
\item \eng{To make a complaint} / \eng{to lodge a complaint} (déposer une réclamation).
\end{itemize}
\end{exemplebox}

\courssub{Comparing Quality and Finding the Best Buy}

\begin{propriete}[Comparatifs et superlatifs : rappels et pièges]
\textbf{Règle de formation :}
\begin{center}
\begin{tabular}{|l|l|l|l|}\hline
\rowcolor{popblueL} Type & Adjectif & Comparatif & Superlatif \\ \hline
Court (1 syllabe) & \eng{cheap} & \eng{cheaper than} & \eng{the cheapest} \\
Court (1 syllabe, e muet) & \eng{nice} & \eng{nicer than} & \eng{the nicest} \\
Court (termine par CVC) & \eng{big} & \eng{bigger than} & \eng{the biggest} \\
Long ($\ge$ 2 syllabes) & \eng{expensive} & \eng{more expensive than} & \eng{the most expensive} \\
Long (termine par -y) & \eng{reliable} & \eng{more reliable than} & \eng{the most reliable} \\
Irréguliers & \eng{good / bad} & \eng{better than / worse than} & \eng{the best / the worst} \\
Irréguliers & \eng{far} & \eng{farther/further than} & \eng{the farthest/furthest} \\
Irréguliers & \eng{little (quantité)} & \eng{less than} & \eng{the least} \\
\hline
\end{tabular}
\end{center}

\textbf{Structures d'égalité et de différence :}
\begin{itemize}
\item \eng{As + adjectif + as} (aussi ... que) : \eng{Seller A is as reliable as Seller C.}
\item \eng{Not as/so + adjectif + as} (moins ... que) : \eng{Seller B is not as durable as Seller A.}
\item \eng{The same price as} / \eng{different from} (UK) / \eng{different than} (US).
\item \eng{Unlike Seller B, Seller C offers a guarantee.}
\item \eng{Whereas / While} (alors que) pour contraster deux sujets : \eng{Seller B is cheap, whereas Seller C is expensive.}
\end{itemize}
\end{propriete}

\begin{definition}[Lexique de la qualité et du jugement]
\begin{itemize}
\item \eng{Quality} /ˈkwɒləti/ (n.) : qualité. Adj. \eng{high-quality} / \eng{poor-quality} / \eng{top-quality}.
\item \eng{Durable} /ˈdjʊərəbəl/ (adj.) : durable, résistant (qui dure longtemps). Pour denrées : \eng{long shelf life} (longue durée de conservation).
\item \eng{Reliable} /rɪˈlaɪəbl/ (adj.) : fiable, sur qui/quoi on peut compter.
\item \eng{Guarantee} /ˌɡærənˈtiː/ (n.) / \eng{Guarantee} (v.) : garantie / garantir. \eng{Warranty} (n.) : garantie (souvent commerciale, écrite). \eng{Under guarantee} (sous garantie).
\item \eng{Standard} /ˈstændəd/ (n.) : norme, standard. \eng{Up to standard} (conforme aux normes).
\item \eng{Defective} /dɪˈfektɪv/ (adj.) : défectueux. \eng{Faulty} /ˈfɔːlti/ (adj.) : défectueux, qui a un défaut.
\item \eng{Worth} /wɜːθ/ (adj., jamais \eng{to be worthing}) : valoir. \eng{It is worth the price.} / \eng{It is worth buying.} / \eng{It is worth it.}
\item \eng{Best buy} (n.) : le meilleur achat (meilleur rapport qualité-prix).
\end{itemize}
\end{definition}

\courssub{Shopping, Bargaining and Consumer Behaviour}

\begin{textebox}[Dialogue : At the Market in Buea]
\textbf{Customer:} Good morning. How much is the bag of local rice?
\textbf{Seller:} It is 15,000 FCFA for 25 kg, madam. Top quality, long grain.
\textbf{Customer:} That's a bit steep. Madam Nguefack down the road charges 14,500 FCFA.
\textbf{Seller:} Ah, but her rice is broken grain. Mine is whole and stores well. I can give you a 2\% discount if you take three bags.
\textbf{Customer:} No discount on a single bag? It's not really value for money if I have to sort out stones.
\textbf{Seller:} No stones here, I guarantee it. Cash or mobile money?
\textbf{Customer:} Mobile money. Let me check the grains first... Okay, it looks clean. I'll take one bag.
\textbf{Seller:} Thank you, madam. Here is your receipt. Keep it for the guarantee.
\end{textebox}

\begin{attention}[Faux amis et calques à bannir]
\begin{itemize}
\item \eng{Actually} \textbf{$\neq$} « actuellement » (=\eng{currently} / \eng{at present}). \eng{Actually} = « en fait, en réalité ». \eng{Actually, the rice is local.}
\item \eng{Economic} \textbf{$\neq$} « économique » (au sens « qui fait économiser »). \eng{Economic} = relatif à l'économie (macro). \eng{Economical} = qui permet d'économiser (peu coûteux à l'usage). \eng{An economical car.}
\item \eng{Expensive} \textbf{$\neq$} \eng{expensif} (n'existe pas).
\item \eng{Payment delay} \textbf{$\neq$} « délai de paiement » (accordé). \eng{Payment delay} = « retard de paiement ». Dire \eng{a credit period} ou \eng{deferred payment}.
\item \eng{Eventually} \textbf{$\neq$} « éventuellement ». \eng{Eventually} = « finalement, à la fin ». \eng{Eventually} = \eng{possibly} / \eng{potentially}.
\item \eng{Quality} (n.) $\to$ Adj. \eng{Qualitative} (rare, technique) vs \eng{Quality} (adj. attributif) : \eng{quality rice}. Pas \eng{qualitative rice}.
\item \eng{Cheaper that} $\to$ \textbf{Interdit}. Toujours \eng{cheaper \textbf{than}}.
\item \eng{The most cheap} $\to$ \textbf{Interdit}. \eng{The cheapest}.
\item \eng{A poor quality rice} $\to$ \textbf{Interdit}. \eng{Rice} est indénombrable. Dire \eng{poor quality rice} ou \eng{rice of poor quality}.
\end{itemize}
\end{attention}

\begin{exemplebox}[Prononciation : paires minimales et pièges]
\begin{center}
\begin{tabular}{|l|l|l|}\hline
\rowcolor{popblueL} Mot & Prononciation (approx. fr.) & Attention \\ \hline
\eng{Cheap} & « tchiip » (voyelle longue /i:/) & $\neq$ \eng{Chip} (« tchip », voyelle courte /ɪ/) \\ \hline
\eng{Price} & « praïss » (diphtongue /aɪ/) & $\neq$ \eng{Prize} (« praïz », /z/ final) \\ \hline
\eng{Receipt} & « risse » (\emph{p} muet, /rɪˈsiːt/) & $\neq$ \eng{Recipe} (« ré-sip », /ˈresəpi/) \\ \hline
\eng{Guarantee} & « ga-ran-ti » (accent sur \emph{ti}) & $\neq$ \eng{Guard} (« gard ») \\ \hline
\eng{Discount} (n.) & « dis-kaount » (accent sur \emph{dis}) & Verbe \eng{to discount} : accent sur \emph{count} \\ \hline
\end{tabular}
\end{center}
\end{exemplebox}

\courssub{Community Resources and Public Services}

\begin{textebox}[Consumer Protection in Cameroon — Key Institutions]
In Cameroon, several public bodies ensure fair trade and protect consumers.
\textbf{MINCOMMERCE} (Ministry of Trade) sets price ceilings on essential goods (rice, fuel, cement) and monitors market practices.
\textbf{ANOR} (Agence des Normes et de la Qualité) certifies products with the \eng{N-mark} (norme camerounaise). It checks standards for food safety, construction materials, and electrical appliances.
\textbf{CAPHYA} (Comité d'Action pour la Protection des Hommes et de l'Environnement) and consumer associations like \textbf{ACOC} (Association Camerounaise des Organisations de Consommateurs) educate buyers on their rights: right to safety, right to information, right to choose, right to be heard, right to redress.
\textbf{The Consumer's Charter:} Every consumer has the right to a receipt, to a product free from defects, and to a refund or replacement if the product is not \eng{fit for purpose}.
\end{textebox}

\begin{aretenir}[Vocabulaire des services publics et droits du consommateur]
\begin{itemize}
\item \eng{Consumer rights} / \eng{Consumer protection} (droits du consommateur / protection).
\item \eng{Price ceiling} / \eng{price cap} (plafond des prix).
\item \eng{Price gouging} (hausse abusive des prix, spéculation).
\item \eng{Standard bureau} / \eng{Standards agency} (organisme de normalisation).
\item \eng{Compliance} / \eng{to comply with standards} (conformité / respecter les normes).
\item \eng{Fit for purpose} (adapté à l'usage prévu — terme juridique clé).
\item \eng{Redress} /rɪˈdres/ (n.) : réparation, indemnisation. \eng{To seek redress}.
\item \eng{Receipt} (reçu) : preuve d'achat obligatoire pour toute réclamation.
\item \eng{Market regulation} / \eng{to regulate markets} (régulation des marchés / réguler).
\end{itemize}
\end{aretenir}

\courssub{Word Families, False Friends and Pronunciation}

\begin{center}
\begin{tabularx}{\linewidth}{|X|X|X|X|}\hline
\rowcolor{popblueL} Verbe & Nom (chose/abstract) & Nom (agent) & Adjectif / Participe \\ \hline
\eng{Consume} & \eng{Consumption} & \eng{Consumer} & \eng{Consumable} \\ \hline
\eng{Produce} & \eng{Production} / \eng{Product} & \eng{Producer} & \eng{Productive} \\ \hline
\eng{Sell} & \eng{Sale} / \eng{Selling} & \eng{Seller} / \eng{Vendor} & \eng{Saleable} / \eng{Marketable} \\ \hline
\eng{Buy} / \eng{Purchase} & \eng{Purchase} / \eng{Buy} & \eng{Buyer} / \eng{Purchaser} & \eng{Purchasable} \\ \hline
\eng{Advertise} & \eng{Advertisement} (UK) / \eng{Ad} & \eng{Advertiser} & \eng{Advertised} \\ \hline
\eng{Guarantee} & \eng{Guarantee} / \eng{Warranty} & \eng{Guarantor} & \eng{Guaranteed} \\ \hline
\eng{Regulate} & \eng{Regulation} & \eng{Regulator} & \eng{Regulatory} \\ \hline
\eng{Complain} & \eng{Complaint} & \eng{Complainant} & \eng{Complaining} \\ \hline
\end{tabularx}
\end{center}

\begin{attention}[Orthographe et pièges morphologiques]
\begin{itemize}
\item \eng{Advertisement} (UK, 4 syllabes /ədˈvɜːtɪsmənt/) vs \eng{Advertizing} (US). Nom : \eng{Advert} (familier UK) / \eng{Ad}.
\item \eng{Complain} (v.) vs \eng{Complaint} (n., prononcé /kəmˈpleɪnt/).
\item \eng{Choice} (n.) vs \eng{Choose} (v., irr. \eng{chose}, \eng{chosen}). Ne pas écrire \eng{choise}.
\item \eng{Lose} (v., perds) vs \eng{Loose} (adj., large). \eng{Don't lose your receipt.} / \eng{Loose change} (monnaie).
\item \eng{Price} / \eng{Prize} / \eng{Praise} (louange).
\item \eng{Customer} (client) vs \eng{Customs} (douane). \eng{Customs officer} = douanier.
\end{itemize}
\end{attention}

\courssub{Using the Lexis in Context: Intégration — The School Canteen Tender}

\courssub{Situation}

Le club d'anglais du lycée de Bafoussam a été chargé par la coopérative des enseignants de choisir le fournisseur de sacs de riz de 25 kg pour la cantine scolaire. Trois vendeurs du marché central ont soumis leurs offres. Le club doit comparer les prix et la qualité, rédiger un rapport et recommander le meilleur achat (\eng{the best buy}).

\begin{textebox}[The Three Sellers' Offers — Central Market, Bafoussam]
\textbf{Seller A — Madam Nguefack's Store.} Price: 14,500 FCFA per 25 kg bag of local rice. Quality: the rice is clean and well dried, but the grains are small and broken. Delivery is free. Madam Nguefack offers a 5\% discount for orders of more than twenty bags. Payment in cash only.

\textbf{Seller B — Cheaper Price Shop.} Price: 12,000 FCFA per 25 kg bag of imported rice. This is the cheapest offer on the market. However, the rice sometimes contains small stones, and several customers have complained about the poor quality. No delivery, no discount, no guarantee.

\textbf{Seller C — Quality First Ltd.} Price: 17,000 FCFA per 25 kg bag of high-quality local rice. The grains are long, whole and have a long shelf life. The rice is guaranteed for six months. Free delivery to the school and a ten-day credit period are offered. It is the most expensive offer, but the shop has an excellent reputation.
\end{textebox}

\begin{center}
\begin{tabularx}{\linewidth}{|X|X|X|X|}\hline
\rowcolor{popblueL} Criterion & Seller A & Seller B & Seller C \\ \hline
Price per bag & 14,500 FCFA & 12,000 FCFA & 17,000 FCFA \\ \hline
Quality & Medium (broken grains) & Poor (stones, impurities) & High (long, whole grains) \\ \hline
Discount / Terms & 5\% over 20 bags & None & 10-day credit period \\ \hline
Delivery & Free & None & Free \\ \hline
Guarantee & None & None & 6 months \\ \hline
\end{tabularx}
\end{center}

\courssub{Consignes}

Travail en groupes de quatre élèves. Chaque groupe suit les étapes dans l'ordre.

\begin{enumerate}
\item \textbf{Compréhension des documents.} Lisez attentivement les trois offres. Soulignez dans chaque offre : le prix, la qualité, les services (livraison, garantie, remise, conditions de paiement). Vérifiez dans le glossaire de la leçon le sens des mots difficiles (\eng{discount}, \eng{guarantee}, \eng{long shelf life}, \eng{credit period}, \eng{fit for purpose}).
\item \textbf{Comparaison orale préparatoire.} Dans votre groupe, comparez les trois offres à voix haute en utilisant au moins cinq structures de comparaison : \eng{Seller B is cheaper than Seller A}, \eng{the most expensive}, \eng{better value for money}, \eng{not as durable as}, \eng{unlike Seller B...}. Chaque membre doit produire au moins deux phrases.
\item \textbf{Choix du best buy.} Mettez-vous d'accord sur le vendeur à recommander. Attention : le moins cher n'est pas toujours le meilleur achat ! Justifiez votre choix par le rapport qualité-prix (\eng{value for money}) et les services inclus.
\item \textbf{Rédaction du rapport.} Rédigez ensemble un rapport de comparaison de 6 à 8 lignes en anglais. Structure attendue : présentation des trois offres (1-2 lignes), comparaison des prix (2 lignes), comparaison de la qualité et des services (2 lignes), conclusion avec le \eng{best buy} et sa justification (1-2 lignes).
\item \textbf{Présentation orale.} Deux membres du groupe présentent le rapport à la classe en deux minutes maximum. Parlez clairement, sans lire mot à mot, en articulant les nombres (14,500 = “fourteen thousand five hundred”).
\item \textbf{Débat final.} Les groupes qui ont choisi des vendeurs différents s'opposent. Chaque groupe défend son choix avec au moins trois expressions de la leçon : \eng{a bargain}, \eng{a rip-off}, \eng{worth the price}, \eng{you get what you pay for}, \eng{fit for purpose}, \eng{false economy}. La classe vote ensuite pour le meilleur argumentaire.
\end{enumerate}

\begin{attention}[Pièges à éviter pendant l'activité]
\begin{itemize}
\item Ne dites pas \eng{the most cheap} : le comparatif de \eng{cheap} est \eng{cheaper}, le superlatif \eng{the cheapest}.
\item \eng{Actually} signifie « en fait, en réalité », jamais « actuellement » (qui se dit \eng{currently} ou \eng{at present}).
\item On dit \eng{cheaper than}, jamais \eng{cheaper that} ni \eng{more cheap than}.
\item \eng{Economic} = relatif à l'économie du pays ; \eng{economical} = qui permet d'économiser. Le riz du vendeur B est \eng{economical} (à l'achat), pas \eng{economic}.
\item \eng{Payment delay} = « retard de paiement ». Pour « délai de paiement accordé », dites \eng{a credit period} ou \eng{deferred payment}.
\item \eng{Rice} est indénombrable : pas \eng{a rice}, pas \eng{rices}. Dire \eng{a bag of rice}, \eng{some rice}, \eng{rice of poor quality}.
\item Prononcez \eng{cheap} « tchiip » (voyelle longue) et ne le confondez pas avec \eng{chip} (« tchip », voyelle courte).
\item \eng{Receipt} se prononce « risse » (le \emph{p} est muet).
\end{itemize}
\end{attention}

\courssub{Ressources à mobiliser}

\begin{aretenir}[Boîte à outils de la comparaison — Synthèse]
\begin{itemize}
\item \textbf{Comparatifs et superlatifs :} \eng{cheap / cheaper / the cheapest} ; \eng{expensive / more expensive / the most expensive} ; \eng{good / better / the best} ; \eng{bad / worse / the worst} ; \eng{far / further / the furthest}.
\item \textbf{Égalité et différence :} \eng{as cheap as}, \eng{not as expensive as}, \eng{the same price as}, \eng{unlike Seller B...}, \eng{whereas / while}.
\item \textbf{Lexique des prix :} \eng{a bargain} (une bonne affaire), \eng{a rip-off} (une arnaque), \eng{a discount} (une remise), \eng{to afford} (avoir les moyens de), \eng{to charge} (demander comme prix), \eng{value for money} (rapport qualité-prix), \eng{price ceiling} (plafond de prix).
\item \textbf{Lexique de la qualité :} \eng{durable} / \eng{long shelf life}, \eng{reliable}, \eng{high-quality / poor quality}, \eng{guaranteed}, \eng{worth the price}, \eng{fit for purpose}, \eng{defective}.
\item \textbf{Expressions du débat :} \eng{In our opinion...}, \eng{We believe that...}, \eng{You get what you pay for}, \eng{It is worth paying more because...}, \eng{False economy} (fausse économie : acheter moins cher mais remplacer plus souvent).
\end{itemize}
\end{aretenir}

\begin{methode}[Comment construire le rapport de comparaison]
\begin{enumerate}
\item \textbf{Introduction (1 phrase) :} Présentez l'objet et le nombre d'offres. \eng{We compared three offers for 25 kg bags of rice from the Central Market in Bafoussam.}
\item \textbf{Comparaison des prix (2 lignes) :} Classez du moins cher au plus cher avec superlatifs et comparatifs. \eng{Seller B is the cheapest at 12,000 FCFA, whereas Seller C is the most expensive at 17,000 FCFA. Seller A is cheaper than Seller C but more expensive than Seller B.}
\item \textbf{Comparaison qualité / services (2 lignes) :} Nuancez avec \eng{however}, \eng{although}, \eng{moreover}. \eng{However, price is not the only criterion. Seller B's rice is of poor quality and contains stones. Although Seller C is more expensive, it offers high-quality rice, a six-month guarantee, free delivery and a credit period.}
\item \textbf{Conclusion (1-2 lignes) :} Recommandation claire + justification \eng{value for money}. \eng{In conclusion, we recommend Seller C as the best buy because it offers the best value for money: the rice is fit for purpose and guaranteed.}
\end{enumerate}
\end{methode}

\courssub{Proposition de production et corrigé commenté}

\begin{exoresolu}[Rapport de comparaison modèle]
Rédiger un rapport de 6 à 8 lignes comparant les trois offres et concluant sur le meilleur achat.
\tcblower
\textbf{Production modèle :}

\eng{We compared three offers for 25 kg bags of rice from the Central Market in Bafoussam. Seller B is the cheapest: he charges only 12,000 FCFA per bag, whereas Seller C is the most expensive at 17,000 FCFA. However, the cheapest offer is not the best buy. Seller B's rice is of poor quality: it sometimes contains stones and there is no guarantee. Seller A's rice is cheaper than Seller C's and the 5\% discount is attractive, but the grains are small and broken. Although Seller C is more expensive, its rice is high-quality, has a long shelf life and is guaranteed for six months, with free delivery and a credit period. In conclusion, we recommend Seller C as the best buy because it offers the best value for money: you get what you pay for.}

\textbf{Commentaire des choix de langue :}
\begin{itemize}
\item \eng{the cheapest} / \eng{the most expensive} : superlatifs correctement formés (adjectif court $\to$ \emph{-est}, adjectif long $\to$ \eng{the most}).
\item \eng{he charges only 12,000 FCFA} : emploi du verbe \eng{to charge} (« demander comme prix »), plus naturel que le calque « il coûte ».
\item \eng{cheaper than Seller C's} : comparatif suivi de \emph{than} et du génitif \eng{Seller C's} pour éviter la répétition de \eng{rice}.
\item \eng{However}, \eng{Although}, \eng{In conclusion} : connecteurs qui structurent le rapport en trois temps logiques (prix, qualité/services, recommandation).
\item \eng{value for money} et \eng{you get what you pay for} : réemploi du lexique idiomatique de la leçon pour justifier le choix.
\item \eng{of poor quality} / \eng{high-quality} : tournures correctes ; éviter le calque \eng{a poor quality rice} (indénombrable).
\item \eng{long shelf life} : terme technique précis pour « durable en stockage » (denrées), préférable à \eng{durable} seul.
\item \eng{credit period} : terme exact pour « délai de paiement accordé » (évite le faux ami \eng{payment delay}).
\end{itemize}
\end{exoresolu}

\begin{experience}[Le débat « Best Buy » : mise en situation]
Après les présentations, l'enseignant désigne deux groupes ayant fait des choix différents (ex. Groupe 1 : Seller C ; Groupe 2 : Seller A ou B). Chaque groupe dispose d'une minute pour défendre son vendeur.

\textbf{Exemple d'argumentaire Groupe 1 (Seller C) :} \eng{Choosing Seller B is a false economy. The rice contains stones, it is not fit for purpose. Seller C is worth the price: the guarantee protects us for six months. You get what you pay for.}

\textbf{Exemple d'argumentaire Groupe 2 (Seller A) :} \eng{The school budget is tight. We cannot afford Seller C. Seller A offers the best compromise: free delivery, a discount over 20 bags, and acceptable quality. It is a real bargain for a community canteen.}

La classe vote pour l'argumentaire le plus convaincant, en comptant les expressions de la leçon correctement employées (\eng{value for money}, \eng{fit for purpose}, \eng{false economy}, \eng{bargain}, \eng{rip-off}, \eng{worth the price}).
\end{experience}

\courssub{Bilan de la leçon}

Cette leçon a permis de mobiliser l'ensemble des savoir-faire du thème « Economic Life and Occupations » dans une situation de communication authentique et contextualisée (appel d'offres pour une cantine scolaire au Cameroun). Les élèves ont :

\begin{itemize}
\item lu et exploité des documents commerciaux réalistes en anglais (compréhension de l'écrit / \eng{reading}) ;
\item comparé des prix et des qualités à l'aide des comparatifs, superlatifs et du lexique thématique (\eng{bargain}, \eng{rip-off}, \eng{discount}, \eng{value for money}, \eng{long shelf life}, \eng{credit period}) ;
\item produit un rapport écrit structuré (\eng{writing}) et une présentation orale argumentée (\eng{speaking}) ;
\item défendu un point de vue dans un débat contradictoire, compétence essentielle de l'expression orale en Terminale ;
\item découvert les institutions camerounaises de protection du consommateur (MINCOMMERCE, ANOR, associations) et le vocabulaire associé (\eng{price ceiling}, \eng{standards}, \eng{redress}, \eng{fit for purpose}) ;
\item intégré une leçon d'économie domestique et civique : le meilleur achat n'est pas toujours le moins cher (\eng{cheapest}), mais celui qui offre le meilleur \cle{rapport qualité-prix} (\eng{value for money}) et respecte les normes (\eng{fit for purpose}).
\end{itemize}

\begin{savaistu}[You get what you pay for — Culture consommateur]
L'expression \eng{you get what you pay for} (« on en a pour son argent ») est un proverbe très courant dans les pays anglophones pour dire qu'un prix trop bas cache souvent une mauvaise qualité. Au Cameroun anglophone, on entend aussi au marché : \eng{Cheap things no good, good things no cheap!} — une version en \eng{Pidgin English} qui montre que la sagesse du consommateur traverse les langues et les cultures. Au Royaume-Uni, l'organisme \eng{Which?} teste les produits pour aider les consommateurs à identifier les \eng{Best Buys}. Au Cameroun, l'ANOR joue ce rôle de certification qualité via la \eng{N-mark}.
\end{savaistu}

\newpage

\coursec{Méthodes et exercices résolus}

\coursec{Lesson 15 — Vocabulary: Words and expressions related to comparing prices or quality to determine the best buys for consumer economy and community resources/services}

\courssub{Savoir-faire 1 : Identifier et employer le vocabulaire des prix et de l'argent}

\begin{methode}[Parler des prix et de l'argent en anglais]
Pour décrire un prix, comparer des coûts et parler d'argent sans calquer le français, suis ces étapes :
\begin{enumerate}
\item \textbf{Choisir le bon verbe.} Une chose \emph{costs} (coûte), une personne \emph{pays} (paie) ou \emph{spends} (dépense) : \eng{This bag costs 5,000 francs.} « Ce sac coûte 5 000 francs. » / \eng{I paid 5,000 francs for this bag.} « J'ai payé ce sac 5 000 francs. » / \eng{I spent 5,000 francs on this bag.} « J'ai dépensé 5 000 francs pour ce sac. »
\item \textbf{Utiliser les bons adjectifs.} Un prix est \emph{high} ou \emph{low} (jamais « expensive price ») ; un objet est \emph{cheap}, \emph{affordable}, \emph{expensive} ou \emph{overpriced} : \eng{The price of fuel is very high this year.} « Le prix du carburant est très élevé cette année. »
\item \textbf{Maîtriser les collocations clés.} \emph{to raise/lower prices} (augmenter/baisser les prix), \emph{to charge a fee} (demander des frais), \emph{to afford something} (avoir les moyens de), \emph{value for money} (rapport qualité-prix), \emph{to be worth it} (en valoir la peine).
\item \textbf{Employer les structures de comparaison.} \eng{X is cheaper than Y.} « X est moins cher que Y. » / \eng{X is the cheapest of the three.} « X est le moins cher des trois. » / \eng{X is twice as expensive as Y.} « X coûte deux fois plus cher que Y. »
\item \textbf{Reformuler pour vérifier.} Avant d'écrire, demande-toi : qui dépense ? quoi coûte ? Cela évite les confusions \emph{cost/pay/spend}.
\end{enumerate}
\end{methode}

\begin{popfigure}
\begin{tikzpicture}[
  mindmap,
  concept color=popblue,
  level 1 concept/.append style={level distance=3.5cm, sibling angle=90, concept color=popblue!70},
  level 2 concept/.append style={level distance=2.8cm, sibling angle=45, concept color=popgreen!70},
  every node/.style={font=\small, text=white},
  edge from parent/.style={draw=popline, thick}
]
\node[concept, align=center]{Consumer\\Economy}
  child[clockwise from=90] {node[concept] {Prices \& Money}
    child {node[concept] {cost / pay / spend}}
    child {node[concept] {high / low price}}
    child {node[concept] {afford / charge}}
    child {node[concept] {value for money}}
  }
  child[clockwise from=0] {node[concept] {Quality \& Comparison}
    child {node[concept] {cheaper / best buy}}
    child {node[concept] {durability / guarantee}}
    child {node[concept] {bargain / deal}}
    child {node[concept] {as ... as / twice as}}
  }
  child[clockwise from=-90] {node[concept] {Shopping \& Bargaining}
    child {node[concept] {discount / reduce}}
    child {node[concept] {last price / meet halfway}}
    child {node[concept] {change / receipt}}
    child {node[concept] {second-hand / brand-new}}
  }
  child[clockwise from=180] {node[concept] {Community Resources}
    child {node[concept] {health centre / school}}
    child {node[concept] {water supply / electricity}}
    child {node[concept] {taxes / council / funds}}
    child {node[concept] {public transport / roads}}
  };
\end{tikzpicture}
\legende{Mindmap : champs lexicaux de la leçon 15 (Économie de consommation et ressources communautaires).}
\end{popfigure}

\begin{exoresolu}[Vocabulary in context — prices and money]
\textbf{Énoncé.} Complete the following sentences with the correct word from the box: \emph{afford, costs, charge, worth, raise, spent, low, value}.

\eng{1. Many families in Douala cannot \dots\ school fees this year.\\
2. A taxi ride from the market to the town centre \dots\ about 500 francs.\\
3. Some banks \dots\ high fees for money transfers.\\
4. Traders decided to \dots\ the price of tomatoes because of the fuel shortage.\\
5. This phone is good \dots\ for money: it is cheap and durable.\\
6. She \dots\ all her savings on school books.\\
7. The price of garri is very \dots\ during the harvest season.\\
8. The more expensive uniform is \dots\ it because it lasts three years.}

\tcblower
\textbf{Solution détaillée.}

\textbf{1. afford.} Structure : \emph{can/cannot afford + nom} = avoir (ou non) les moyens de. \eng{Many families in Douala cannot afford school fees this year.} « Beaucoup de familles à Douala ne peuvent pas payer les frais de scolarité cette année. »

\textbf{2. costs.} Le sujet est une chose (\emph{a taxi ride}), donc le verbe est \emph{cost} ; à la 3e personne : \emph{costs}. \eng{A taxi ride from the market to the town centre costs about 500 francs.} « Une course en taxi du marché au centre-ville coûte environ 500 francs. »

\textbf{3. charge.} Collocation : \emph{to charge fees} = demander des frais. \eng{Some banks charge high fees for money transfers.} « Certaines banques facturent des frais élevés pour les transferts d'argent. »

\textbf{4. raise.} \emph{To raise prices} = augmenter les prix (verbe transitif régulier ; ne pas confondre avec \emph{rise}, intransitif : \eng{Prices are rising.} « Les prix augmentent. »).

\textbf{5. value.} Expression figée : \emph{good value for money} = bon rapport qualité-prix. \eng{This phone is good value for money.} « Ce téléphone a un bon rapport qualité-prix. »

\textbf{6. spent.} Structure : \emph{spend + argent + on + nom} ; au prétérit, \emph{spend} devient \emph{spent} (verbe irrégulier). \eng{She spent all her savings on school books.} « Elle a dépensé toutes ses économies en livres scolaires. »

\textbf{7. low.} Un prix est \emph{high} ou \emph{low}, jamais « cheap ». \eng{The price of garri is very low during the harvest season.} « Le prix du gari est très bas pendant la saison des récoltes. »

\textbf{8. worth.} Structure : \emph{to be worth it} = en valoir la peine. \eng{The more expensive uniform is worth it because it lasts three years.} « L'uniforme plus cher en vaut la peine car il dure trois ans. »

\textbf{Conclusion.} Retiens les trois verbes pivots : la chose \emph{costs}, la personne \emph{pays (for)} ou \emph{spends (on)} ; et les deux expressions figées \emph{value for money} et \emph{to be worth it}.
\end{exoresolu}

\courssub{Savoir-faire 2 : Comparer la qualité et déterminer le meilleur achat}

\begin{methode}[Comparer qualité et prix pour trouver le « best buy »]
\begin{enumerate}
\item \textbf{Identifier les critères.} Avant de comparer, liste les critères : \emph{price, quality, durability, size, guarantee, after-sales service}.
\item \textbf{Employer les comparatifs et superlatifs correctement.} Adjectifs courts : \emph{cheap $\rightarrow$ cheaper $\rightarrow$ the cheapest} ; adjectifs longs : \emph{expensive $\rightarrow$ more expensive $\rightarrow$ the most expensive} ; irréguliers : \emph{good $\rightarrow$ better $\rightarrow$ the best}, \emph{bad $\rightarrow$ worse $\rightarrow$ the worst}.
\item \textbf{Utiliser les structures d'équivalence et de proportion.} \eng{X is as reliable as Y.} « X est aussi fiable que Y. » / \eng{X is not as durable as Y.} « X n'est pas aussi durable que Y. » / \eng{X is twice as expensive as Y.} « X coûte deux fois plus cher que Y. »
\item \textbf{Conclure avec le vocabulaire du choix.} \emph{the best buy} (le meilleur achat), \emph{a bargain} (une bonne affaire), \emph{a good deal} (une bonne affaire), \emph{the best option} : \eng{All things considered, the second-hand generator is the best buy.} « Tout compte fait, le groupe électrogène d'occasion est le meilleur achat. »
\item \textbf{Justifier ton choix.} Utilise \emph{because, since, as, that is why} : \eng{I would choose shop B because it offers a one-year guarantee.} « Je choisirais le magasin B car il offre une garantie d'un an. »
\end{enumerate}
\end{methode}

\begin{exoresolu}[Guided writing — choosing the best buy]
\textbf{Énoncé.} Read the advertisement, then answer the questions in full sentences.

\begin{textebox}[Advertisement — Back-to-school offers in Yaoundé]
\eng{MOKOLO ELECTRONICS: Brand-new solar lamp, 12,000 francs. Two-year guarantee. Very bright, charges phones.\\
CITY SHOP: Same model, second-hand, 7,000 francs. No guarantee, slight scratches, works perfectly.\\
MAMA GRACE STORE: New lamp, smaller model, 9,500 francs. Six-month guarantee, less bright, does not charge phones.}
\end{textebox}

\eng{1. Which lamp is the cheapest?\\
2. Which lamp is the most expensive?\\
3. Is the City Shop lamp as bright as the Mokolo lamp?\\
4. Which one is the best buy in your opinion? Justify your answer in two or three sentences.}

\tcblower
\textbf{Solution détaillée.}

\textbf{Question 1.} On compare trois prix : 12,000 / 7,000 / 9,500 francs. Superlatif de \emph{cheap} : \emph{the cheapest}. Réponse : \eng{The City Shop lamp is the cheapest: it costs 7,000 francs.} « La lampe de City Shop est la moins chère : elle coûte 7 000 francs. »

\textbf{Question 2.} Superlatif d'un adjectif long : \emph{the most expensive}. \eng{The Mokolo Electronics lamp is the most expensive, at 12,000 francs.} « La lampe de Mokolo Electronics est la plus chère, à 12 000 francs. »

\textbf{Question 3.} Structure d'égalité : \emph{as + adjectif + as}. Le texte dit que City Shop vend le même modèle qui « works perfectly », donc oui. \eng{Yes, it is as bright as the Mokolo lamp because it is the same model and it works perfectly.} « Oui, elle est aussi lumineuse que la lampe de Mokolo car c'est le même modèle et elle fonctionne parfaitement. »

\textbf{Question 4.} Réponse rédigée avec justification (\emph{because}, \emph{although}) : \eng{In my opinion, the Mokolo Electronics lamp is the best buy. Although it is the most expensive, it has a two-year guarantee and it can charge phones, so it offers the best value for money. A student who buys it will not need to replace it soon.} « À mon avis, la lampe de Mokolo Electronics est le meilleur achat. Bien qu'elle soit la plus chère, elle a une garantie de deux ans et peut recharger les téléphones, donc elle offre le meilleur rapport qualité-prix. Un élève qui l'achète n'aura pas besoin de la remplacer de sitôt. »

\textbf{Conclusion.} Une bonne réponse de « best buy » contient toujours : un choix clair (\emph{the best buy is...}), une concession éventuelle (\emph{although...}) et une justification liée au rapport qualité-prix (\emph{value for money}, \emph{guarantee}, \emph{durability}).
\end{exoresolu}

\courssub{Savoir-faire 3 : Réemployer le lexique à l'oral — marchander et dialoguer au marché}

\begin{methode}[Conduire un dialogue d'achat et de négociation]
\begin{enumerate}
\item \textbf{Ouvrir l'échange poliment.} Vendeur : \eng{Can I help you?} « Puis-je vous aider ? » / \eng{What are you looking for?} « Que cherchez-vous ? » Client : \eng{How much is this?} « Combien coûte ceci ? » / \eng{How much do you sell the plantains?} « À combien vendez-vous les plantains ? »
\item \textbf{Demander et annoncer le prix.} \eng{It costs 2,000 francs.} « Ça coûte 2 000 francs. » / \eng{I am selling it for 2,000 francs.} « Je le vends 2 000 francs. »
\item \textbf{Négocier avec des formules figées.} \eng{That is too expensive!} « C'est trop cher ! » / \eng{Can you reduce the price?} « Pouvez-vous baisser le prix ? » / \eng{Can you give me a discount?} « Pouvez-vous me faire une remise ? » / \eng{My last price is 1,800 francs.} « Mon dernier prix est 1 800 francs. » / \eng{Let us meet halfway: 1,900 francs.} « Coupons la poire en deux : 1 900 francs. »
\item \textbf{Conclure la transaction.} \eng{It is a deal!} « C'est entendu ! / Affaire conclue ! » / \eng{I will take it.} « Je le prends. » / \eng{Here is your change.} « Voici votre monnaie. » / \eng{Thank you, come again.} « Merci, revenez. »
\item \textbf{Soigner la prononciation des nombres.} Distingue \emph{thirteen} (« seur-tine », accent sur la 2e syllabe) et \emph{thirty} (« seur-ti », accent sur la 1re) ; \emph{fifteen} / \emph{fifty}.
\end{enumerate}
\end{methode}

\begin{exoresolu}[Dialogue completion — at Mokolo market]
\textbf{Énoncé.} Complete the dialogue between a customer and a trader with suitable expressions.

\eng{Trader: Good morning, madam. (1) \dots\ ?\\
Customer: Yes, how much (2) \dots\ this bag of rice?\\
Trader: It costs 15,000 francs.\\
Customer: 15,000! That is (3) \dots\ expensive! Can you (4) \dots\ the price?\\
Trader: For you, madam, my (5) \dots\ price is 14,000 francs.\\
Customer: Let us (6) \dots\ halfway: 13,500 francs.\\
Trader: All right, it is a (7) \dots\ ! Here is your bag.\\
Customer: Thank you. Here is the money.\\
Trader: And here is your (8) \dots\ : 500 francs.}

\tcblower
\textbf{Solution détaillée.}

\textbf{(1) Can I help you.} Formule d'accueil standard du vendeur. \eng{Good morning, madam. Can I help you?} « Bonjour madame. Puis-je vous aider ? »

\textbf{(2) is.} \eng{How much is this bag of rice?} — avec \emph{how much}, le verbe \emph{be} s'accorde avec le nom singulier (\emph{this bag}).

\textbf{(3) too.} \emph{Too + adjectif} exprime l'excès : \eng{That is too expensive!} « C'est trop cher ! » (ne pas confondre avec \emph{very} = très).

\textbf{(4) reduce / lower.} Collocation : \emph{to reduce/lower the price}. \eng{Can you reduce the price?} « Pouvez-vous baisser le prix ? »

\textbf{(5) last.} Expression idiomatique du marchandage : \eng{My last price is 14,000 francs.} « Mon dernier prix est de 14 000 francs. »

\textbf{(6) meet.} Expression figée : \emph{to meet halfway} = couper la poire en deux. \eng{Let us meet halfway.} « Arrêtons-nous à mi-chemin. »

\textbf{(7) deal.} \eng{It is a deal!} « C'est d'accord ! / Affaire conclue ! »

\textbf{(8) change.} \emph{Change} (invariable ici) = la monnaie rendue. \eng{Here is your change: 500 francs.} « Voici votre monnaie : 500 francs. »

\textbf{Conclusion.} Un dialogue de marchandage suit toujours la même architecture : accueil $\rightarrow$ demande de prix $\rightarrow$ négociation (\emph{too expensive, discount, last price, meet halfway}) $\rightarrow$ conclusion (\emph{it is a deal, change}). Mémorise ces formules comme des blocs.
\end{exoresolu}

\courssub{Savoir-faire 4 : Employer le vocabulaire des ressources communautaires et des services publics}

\begin{methode}[Parler des services publics et des ressources de la communauté]
\begin{enumerate}
\item \textbf{Mémoriser le champ lexical par catégories.} Santé : \emph{health centre, hospital, pharmacy, vaccination campaign} ; éducation : \emph{public school, library, scholarship} ; eau et énergie : \emph{water supply, tap water, electricity, power cut} ; transport : \emph{public transport, bus fare, road maintenance} ; administration : \emph{town council, community hall, public service}.
\item \textbf{Utiliser les collocations verbales.} \emph{to provide a service} (fournir un service), \emph{to run a health centre} (gérer un centre de santé), \emph{to maintain roads} (entretenir les routes), \emph{to pay taxes} (payer des impôts), \emph{to benefit from a service} (bénéficier d'un service).
\item \textbf{Comparer les services comme on compare les prix.} \eng{The district hospital is better equipped than the health centre.} « L'hôpital de district est mieux équipé que le centre de santé. » / \eng{Public transport is cheaper than taxis but slower.} « Le transport public est moins cher que les taxis mais plus lent. »
\item \textbf{Relier services et économie.} \eng{Taxes pay for community services.} « Les impôts financent les services collectifs. » \eng{The council should lower water bills.} « La mairie devrait baisser les factures d'eau. »
\end{enumerate}
\end{methode}

\begin{savaistu}[Le saviez-vous ? — Services publics au Cameroun et au Royaume-Uni]
Au Cameroun, les \emph{councils} (mairies) gèrent les marchés, l'état civil, les écoles primaires et l'entretien des routes rurales. Les impôts locaux (\emph{council taxes}) financent ces services. Au Royaume-Uni, le \emph{National Health Service (NHS)} offre des soins gratuits financés par l'impôt national ; les \emph{local councils} gèrent les ordures, les bibliothèques et le logement social. Aux États-Unis, pas de système de santé national gratuit : l'assurance privée domine, mais il existe \emph{Medicaid} pour les plus démunis. Le vocabulaire \emph{public service, community hall, water supply, power cut} est transversal à tous ces contextes.
\end{savaistu}

\begin{exoresolu}[Reading comprehension — community services]
\textbf{Énoncé.} Read the text and answer the questions in your own words.

\begin{textebox}[A community meeting in Bamenda]
\eng{Last Saturday, the residents of Nkwen held a meeting at the community hall to discuss local services. Many complained that the water supply is irregular and that power cuts are too frequent. The councillor explained that taxes pay for road maintenance and the health centre, but the budget is limited. A trader suggested that the council should lower bus fares for students. Finally, the community decided to raise funds to buy a generator for the health centre, because it is the most urgent need.}
\end{textebox}

\eng{1. Where did the meeting take place?\\
2. Mention two problems the residents complained about.\\
3. What do taxes pay for, according to the councillor?\\
4. What did the community finally decide to do, and why?}

\tcblower
\textbf{Solution détaillée.}

\textbf{Question 1.} Le texte dit : \emph{at the community hall}. Réponse en phrase complète : \eng{The meeting took place at the community hall in Nkwen.} « La réunion a eu lieu à la salle communautaire de Nkwen. »

\textbf{Question 2.} Deux problèmes sont cités : \emph{irregular water supply} et \emph{frequent power cuts}. \eng{The residents complained that the water supply is irregular and that power cuts are too frequent.} « Les habitants se sont plaints que l'approvisionnement en eau est irrégulier et que les coupures de courant sont trop fréquentes. »

\textbf{Question 3.} Le texte : \emph{taxes pay for road maintenance and the health centre}. \eng{According to the councillor, taxes pay for road maintenance and the health centre.} Attention à la préposition : \emph{to pay for something}.

\textbf{Question 4.} Décision finale : \emph{raise funds to buy a generator} ; justification : \emph{the most urgent need}. \eng{The community decided to raise funds to buy a generator for the health centre because it is the most urgent need.} « La communauté a décidé de collecter des fonds pour acheter un groupe électrogène pour le centre de santé, car c'est le besoin le plus urgent. »

\textbf{Conclusion.} En compréhension écrite, repère les mots du champ lexical (\emph{council, taxes, services, funds}) : ils signalent presque toujours les réponses. Reformule avec tes propres mots sans recopier de longs passages.
\end{exoresolu}

\courssub{Savoir-faire 5 : Éviter les faux amis et les calques du français}

\begin{methode}[Déjouer les faux amis du domaine économique]
\begin{enumerate}
\item \textbf{Repérer les faux amis classiques.} \emph{Actually} = en fait (pas « actuellement » $\rightarrow$ \emph{currently}) ; \emph{economic} = économique (relatif à l'économie), \emph{economical} = économe ; \emph{a library} = une bibliothèque (une librairie = \emph{a bookshop}) ; \emph{sensible} = raisonnable (sensible = \emph{sensitive}).
\item \textbf{Se méfier des calques de structure.} « J'ai acheté ça à bon marché » ne se traduit pas mot à mot : dis \eng{I bought it cheap.} « Je l'ai acheté pas cher. » « Combien ça coûte ? » = \eng{How much does it cost?} (jamais « How much it costs? » en question directe).
\item \textbf{Vérifier la préposition.} \emph{to pay for} (payer quelque chose), \emph{to spend on}, \emph{to depend on}, \emph{to be interested in}, \emph{to consist of}.
\item \textbf{Contrôler le nom dénombrable/indénombrable.} \emph{money, information, advice, furniture, news} sont indénombrables : \eng{The money is on the table.} (jamais « moneys »).
\item \textbf{Relire en traduisant mentalement à l'envers.} Si la phrase anglaise, traduite mot à mot, donne du français bizarre, c'est probablement un calque.
\end{enumerate}
\end{methode}

\begin{exoresolu}[Error correction — false friends and calques]
\textbf{Énoncé.} Each sentence contains one mistake (false friend, wrong preposition or calque from French). Correct it.

\eng{1. Actually, the prices in the market are very high.\\
2. I am interested to buy a second-hand bicycle.\\
3. She spent 3,000 francs for her school uniform.\\
4. The advices the trader gave me were useful.\\
5. How much costs this bag of onions?\\
6. This shop is very economic for students.}

\tcblower
\textbf{Solution détaillée.}

\textbf{1.} Faux ami : \emph{actually} signifie « en fait ». Pour « actuellement », il faut \emph{currently} ou \emph{at present}. Correction : \eng{Currently, the prices in the market are very high.} « Actuellement, les prix au marché sont très élevés. »

\textbf{2.} Mauvaise préposition : \emph{to be interested in + verbe-ing}. Correction : \eng{I am interested in buying a second-hand bicycle.} « Je suis intéressé par l'achat d'un vélo d'occasion. »

\textbf{3.} Calque du français « dépenser pour » : la structure anglaise est \emph{spend + montant + on + nom}. Correction : \eng{She spent 3,000 francs on her school uniform.} « Elle a dépensé 3 000 francs pour son uniforme scolaire. »

\textbf{4.} \emph{Advice} est indénombrable : pas de pluriel, verbe au singulier. Correction : \eng{The advice the trader gave me was useful.} « Les conseils que le commerçant m'a donnés étaient utiles. »

\textbf{5.} Question directe : inversion avec l'auxiliaire \emph{does}. Correction : \eng{How much does this bag of onions cost?} « Combien coûte ce sac d'oignons ? »

\textbf{6.} Faux ami : \emph{economic} = relatif à l'économie d'un pays ; pour « avantageux, pas cher », on dit \emph{cheap}, \emph{affordable} ou \emph{economical}. Correction : \eng{This shop is very affordable for students.} « Ce magasin est très abordable pour les élèves. »

\textbf{Conclusion.} Trois réflexes sauvent des points : vérifier la préposition après les verbes courants, traiter \emph{money/advice/information} comme indénombrables, et remplacer tout mot « trop français » (\emph{actually, economic, assist}) par son vrai équivalent anglais.
\end{exoresolu}

\begin{attention}[Les 5 erreurs qui coûtent des points au Bac]
\begin{enumerate}
\item \textbf{Confondre \emph{cost}, \emph{pay} et \emph{spend}.} On écrit \eng{The book costs 2,000 francs.} « Le livre coûte 2 000 francs. » / \eng{I paid 2,000 francs for the book.} « J'ai payé le livre 2 000 francs. » / \eng{I spent 2,000 francs on the book.} « J'ai dépensé 2 000 francs pour le livre. » Jamais « The book pays... » ni « I cost 2,000 francs ».
\item \textbf{Dire « the price is expensive ».} Un prix est \emph{high} ou \emph{low} ; seul l'objet est \emph{expensive} ou \emph{cheap}. Écris \eng{The price is too high.} « Le prix est trop élevé. » ou \eng{The lamp is too expensive.} « La lampe est trop chère. »
\item \textbf{Calquer les questions directes.} « Combien ça coûte ? » se dit \eng{How much does it cost?} avec l'auxiliaire \emph{does}. La forme « How much it costs? » n'existe qu'en interrogation indirecte : \eng{Can you tell me how much it costs?} « Peux-tu me dire combien ça coûte ? »
\item \textbf{Pluraliser les indénombrables.} \emph{Money, advice, information, furniture, news} n'ont pas de pluriel : \eng{The information is useful.} « L'information est utile. » Pour compter, utilise \emph{a piece of advice}, \emph{a sum of money}.
\item \textbf{Tomber dans les faux amis.} \emph{Actually} = en fait ; \emph{economic} = relatif à l'économie ; \emph{a library} = une bibliothèque ; \emph{sensible} = raisonnable. Une seule de ces confusions dans une composition coûte cher en exactitude lexicale.
\end{enumerate}
\end{attention}

\courssub{Exercices d'entraînement — Vocabulary consolidation}

\begin{exercice}[Practice exercises — Lesson 15]
\textbf{Exercise A: Multiple Choice.} Choose the correct option (a, b or c) to complete each sentence.
\begin{enumerate}
\item The price of cement has \dots\ sharply this month. \quad a) raised \quad b) risen \quad c) rose
\item I cannot \dots\ to buy a new laptop right now. \quad a) afford \quad b) pay \quad c) cost
\item This second-hand car is a real \dots\ ; it runs perfectly and was cheap. \quad a) bargain \quad b) value \quad c) fee
\item The \dots\ decided to repair the road leading to the market. \quad a) council \quad b) counsel \quad c) console
\item \dots\ water supply is a major problem in this neighbourhood. \quad a) Irregular \quad b) Irregularity \quad c) Irregularly
\end{enumerate}

\textbf{Exercise B: Word Formation.} Complete the table.
\begin{center}
\begin{tabular}{|l|l|l|l|}\hline
\rowcolor{popblueL} \textbf{Verb} & \textbf{Noun (thing/action)} & \textbf{Noun (person)} & \textbf{Adjective} \\ \hline
to consume & consumption & consumer & \dots \\ \hline
to provide & \dots & provider & \dots \\ \hline
to maintain & maintenance & \dots & \dots \\ \hline
\dots & comparison & \dots & comparable \\ \hline
to bargain & \dots & bargainer & \dots \\ \hline
\end{tabular}
\end{center}

\textbf{Exercise C: Writing.} You want to buy a smartphone. You have seen two offers:
\begin{itemize}
\item \textbf{Shop A:} Brand new, 150,000 FCFA, 1-year warranty, 128 GB.
\item \textbf{Shop B:} Refurbished, 95,000 FCFA, 3-month warranty, 64 GB.
\end{itemize}
Write a short paragraph (60–80 words) explaining which one you would choose as the \emph{best buy}. Use at least five words from the lesson (\emph{value for money, guarantee, affordable, durable, worth it, best buy, bargain, quality}).
\end{exercice}

\begin{corrige}[Exercice A]
\textbf{1. b) risen} (\emph{rise} = augmenter [intransitif] ; \emph{raise} = augmenter qqch [transitif]). \eng{The price of cement has risen sharply.}
\textbf{2. a) afford} (\emph{can afford} = avoir les moyens). \eng{I cannot afford to buy a new laptop right now.}
\textbf{3. a) bargain} (une bonne affaire). \eng{This second-hand car is a real bargain.}
\textbf{4. a) council} (conseil municipal). \emph{Counsel} = conseil/avocat ; \emph{console} = consoler.
\textbf{5. a) Irregular} (adjectif qualificatif). \eng{Irregular water supply is a major problem.}
\end{corrige}

\begin{corrige}[Exercice B]
\begin{center}
\begin{tabular}{|l|l|l|l|}\hline
\rowcolor{popblueL} \textbf{Verb} & \textbf{Noun (thing/action)} & \textbf{Noun (person)} & \textbf{Adjective} \\ \hline
to consume & consumption & consumer & \textbf{consumer / consumable} \\ \hline
to provide & \textbf{provision} & provider & \textbf{providential / provided} \\ \hline
to maintain & maintenance & \textbf{maintainer} & \textbf{maintainable} \\ \hline
\textbf{to compare} & comparison & \textbf{comparator} & comparable \\ \hline
to bargain & \textbf{bargain / bargaining} & bargainer & \textbf{bargain (adj.) / bargaining} \\ \hline
\end{tabular}
\end{center}
\end{corrige}

\begin{corrige}[Exercice C — Modèle de rédaction]
\eng{In my opinion, the smartphone from Shop A is the best buy. Although Shop B is more affordable at 95,000 FCFA, the refurbished model only has a three-month guarantee and half the storage. The brand-new phone from Shop A offers better value for money because it is more durable with a one-year warranty and 128 GB. It is worth it to pay more now to avoid repair costs later. Quality and guarantee make it the best option for a student.}
\end{corrige}

\newpage

\coursec{Exercices}

\coursec{Lesson 15 — Vocabulary: Comparing Prices, Quality \& Community Resources}

\courssub{Talking about prices and money}

\begin{aretenir}[Vocabulaire de base : l'argent et les prix]
\begin{center}
\begin{tabularx}{\linewidth}{|X|X|X|}
\hline
\rowcolor{popblueL} \textbf{English} & \textbf{Français} & \textbf{Collocations / Exemples} \\
\hline
\eng{price} & prix & \eng{high/low price}, \eng{price tag}, \eng{price increase} \\
\eng{cost} & coût, coûter & \eng{cost of living}, \eng{at no cost}, \eng{cover the cost} \\
\eng{fare} & tarif (transport) & \eng{bus fare}, \eng{train fare}, \eng{air fare} \\
\eng{fee} & frais, honoraires & \eng{entrance fee}, \eng{tuition fees}, \eng{service fee} \\
\eng{rate} & taux, tarif & \eng{exchange rate}, \eng{interest rate}, \eng{hourly rate} \\
\eng{charge} & frais (supplément) & \eng{delivery charge}, \eng{bank charges}, \eng{free of charge} \\
\eng{receipt} & reçu, ticket de caisse & \eng{keep the receipt}, \eng{ask for a receipt} \\
\eng{invoice} & facture & \eng{pay an invoice}, \eng{invoice number} \\
\eng{budget} & budget & \eng{tight budget}, \eng{stick to a budget}, \eng{annual budget} \\
\eng{afford} & avoir les moyens de & \eng{can't afford it}, \eng{afford a holiday} \\
\hline
\end{tabularx}
\end{center}
\end{aretenir}

\begin{attention}[Faux amis et pièges fréquents]
\begin{itemize}
    \item \eng{Receipt} $\neq$ \emph{recette} (cuisine) = \eng{recipe}. \eng{Receipt} = ticket de caisse / reçu. Prononciation : « ri-sit » (le \eng{p} est muet).
    \item \eng{Prescription} = ordonnance (médecin) ; \eng{recipe} = recette (cuisine).
    \item \eng{Fare} (transport) vs \eng{Fee} (service, entrée, scolarité) vs \eng{Price} (objet, bien).
    \item \eng{Cost} (verbe irrégulier : cost, cost, cost) vs \eng{Cost} (nom). \eng{It costs 500 FCFA.}
    \item \eng{Afford} est toujours précédé d'un modal (\eng{can, could, be able to}) et suivi d'un infinitif ou d'un nom, JAMAIS de \eng{to} seul. \eng{I can't afford to buy it.} / \eng{I can't afford it.}
\end{itemize}
\end{attention}

\begin{exemplebox}[Dialogue : Au marché Mokolo (Yaoundé)]
\textbf{Vendor:} Good morning, Madam. This wax print is top quality. \eng{5,000 FCFA} per yard.
\textbf{Customer:} \eng{5,000?} That's \eng{overpriced}! The vendor next door sells it for \eng{3,500}.
\textbf{Vendor:} Ah, but hers is \eng{counterfeit}. This one is \eng{genuine}. I can give you a \eng{discount} of \eng{10\%} if you take \eng{6 yards}.
\textbf{Customer:} Okay, \eng{4,000} final price and you give me a \eng{receipt}.
\textbf{Vendor:} Deal. \eng{Here is your change and the receipt.}
\end{exemplebox}

\courssub{Comparing quality and finding the best buy}

\begin{aretenir}[Qualité, valeur et comparaison]
\begin{center}
\begin{tabularx}{\linewidth}{|X|X|X|}
\hline
\rowcolor{popblueL} \textbf{English} & \textbf{Français} & \textbf{Collocations / Exemples} \\
\hline
\eng{quality} & qualité & \eng{high/poor quality}, \eng{quality control}, \eng{top quality} \\
\eng{value for money} & rapport qualité-prix & \eng{good/excellent value for money}, \eng{best value for money} \\
\eng{bargain} & bonne affaire & \eng{real bargain}, \eng{pick up a bargain}, \eng{bargain hunter} \\
\eng{overpriced} & trop cher, surévalué & \eng{slightly overpriced}, \eng{grossly overpriced} \\
\eng{affordable} & abordable, accessible & \eng{affordable housing}, \eng{affordable prices}, \eng{affordable for students} \\
\eng{cheap} & bon marché (souvent péjoratif : médiocre) & \eng{cheap and nasty}, \eng{dirt cheap} \\
\eng{expensive} & cher & \eng{too expensive}, \eng{expensive taste} \\
\eng{worth} & valoir la peine / valoir & \eng{worth buying}, \eng{worth the money}, \eng{worth \$100} \\
\eng{genuine} & authentique, vrai & \eng{genuine leather}, \eng{genuine product}, \eng{genuine parts} \\
\eng{counterfeit} / \eng{fake} & contrefait, faux & \eng{counterfeit goods}, \eng{fake brand}, \eng{spot a fake} \\
\hline
\end{tabularx}
\end{center}
\end{aretenir}

\begin{propriete}[Structures grammaticales clés pour comparer]
\begin{itemize}
    \item \textbf{Comparatif de supériorité/infériorité} : \eng{cheaper than}, \eng{more expensive than}, \eng{better quality than}, \eng{less reliable than}.
    \item \textbf{Superlatif} : \eng{the cheapest}, \eng{the most affordable}, \eng{the best value for money}.
    \item \textbf{\eng{Worth} + V-ing} (valoir la peine de) : \eng{This phone is worth buying.} (Pas \eng{worth to buy}).
    \item \textbf{\eng{Value for money}} (nom invariable) : \eng{It offers good value for money.}
    \item \textbf{\eng{Prefer} + V-ing + to + V-ing} : \eng{I prefer shopping at the market to shopping at the supermarket.}
    \item \textbf{\eng{Used to} + V-ing} (habitude) vs \textbf{\eng{Used to} + V} (passé) : \eng{I am used to comparing prices.} / \eng{I used to compare prices.}
\end{itemize}
\end{propriete}

\begin{exemplebox}[Exemples traduits]
\eng{This local rice is better value for money than the imported one.} « Ce riz local offre un meilleur rapport qualité-prix que celui importé. »
\eng{The warranty makes it worth the extra cost.} « La garantie vaut le coût supplémentaire. »
\eng{Don't be fooled by cheap prices; check the quality first.} « Ne vous laissez pas tromper par les prix bas ; vérifiez d'abord la qualité. »
\end{exemplebox}

\courssub{Shopping, bargaining and consumer behaviour}

\begin{textebox}[Script d'écoute : Négocier au Marché Central (Douala)]
\textbf{Narrator:} At the Central Market, bargaining (\eng{haggling}) is part of the culture.
\textbf{Buyer:} Excuse me, \eng{how much is this bag?}
\textbf{Seller:} \eng{15,000 FCFA}. It is \eng{genuine leather}.
\textbf{Buyer:} \eng{15,000?} That's \eng{a bit steep}. I saw a similar one for \eng{10,000}.
\textbf{Seller:} That one was \eng{fake}. For you, \eng{13,000}. Last price.
\textbf{Buyer:} I'll give you \eng{11,000}. Cash, right now.
\textbf{Seller:} \eng{12,000} and I throw in a \eng{wallet}.
\textbf{Buyer:} Done. \eng{Here is the money.} Do I get a \eng{receipt} and a \eng{warranty}?
\textbf{Seller:} \eng{No warranty on market goods}, Madam. But I guarantee the leather.
\textbf{Narrator:} The buyer made an \eng{impulse buy}. He didn't \eng{compare} prices elsewhere. He should \eng{check the stitching} and \eng{ask for a receipt} in case he needs a \eng{refund} later.
\end{textebox}

\begin{methode}[Comment marchander (Haggling) en anglais]
\begin{enumerate}
    \item \textbf{Demander le prix} : \eng{How much is this?} / \eng{What's the price?}
    \item \textbf{Exprimer la surprise / refuser} : \eng{That's too expensive.} / \eng{That's a bit steep.} / \eng{I can't afford that.}
    \item \textbf{Faire une contre-offre} : \eng{I'll give you [prix].} / \eng{Would you take [prix]?} / \eng{Can you do it for [prix]?}
    \item \textbf{Argumenter} : \eng{It's cheaper elsewhere.} / \eng{It has a small defect.} / \eng{I'm buying two.}
    \item \textbf{Conclure} : \eng{Deal.} / \eng{You've got a deal.} / \eng{Okay, sold.} / \eng{No deal, thanks anyway.}
\end{enumerate}
\end{methode}

\begin{propriete}[Verbes + Prépositions / Gérondif / Infinitif (Contexte achat)]
\begin{center}
\begin{tabular}{|l|l|l|}
\hline
\rowcolor{popblueL} \textbf{Verbe} & \textbf{Structure} & \textbf{Exemple} \\
\hline
\eng{insist on} & V-ing & \eng{She insisted on paying cash.} \\
\eng{refuse} & to V & \eng{He refused to give a refund.} \\
\eng{avoid} & V-ing & \eng{Avoid buying counterfeits.} \\
\eng{worth} & V-ing & \eng{It's worth checking.} \\
\eng{aim} & to V & \eng{We aim to protect consumers.} \\
\eng{think of} & V-ing & \eng{We're thinking of switching providers.} \\
\eng{look forward to} & V-ing & \eng{I look forward to seeing the new market.} \\
\eng{used to} & V-ing & \eng{I'm used to haggling.} \\
\hline
\end{tabular}
\end{center}
\end{propriete}

\courssub{Community resources and public services}

\begin{aretenir}[Ressources communautaires et services publics]
\begin{center}
\begin{tabularx}{\linewidth}{|X|X|X|}
\hline
\rowcolor{popblueL} \textbf{English} & \textbf{Français} & \textbf{Collocations / Contexte camerounais} \\
\hline
\eng{amenities} & équipements collectifs, commodités & \eng{public amenities}, \eng{local amenities} (eau, électricité, bancs publics) \\
\eng{infrastructure} & infrastructures & \eng{road infrastructure}, \eng{improve infrastructure} \\
\eng{public transport} & transports en commun & \eng{reliable public transport}, \eng{bus fare}, \eng{taxi fare} \\
\eng{community centre} & centre communautaire & \eng{community hall}, \eng{youth centre} \\
\eng{health centre} / \eng{clinic} & centre de santé, dispensaire & \eng{public clinic}, \eng{consultation fee} \\
\eng{water supply} & alimentation en eau & \eng{water cuts}, \eng{public tap}, \eng{borehole} \\
\eng{consumer rights} & droits des consommateurs & \eng{consumer protection}, \eng{consumer association} \\
\eng{tax} / \eng{levy} & impôt, taxe & \eng{local taxes}, \eng{council tax}, \eng{VAT} (\eng{Value Added Tax}) \\
\eng{subsidy} & subvention & \eng{government subsidy}, \eng{fuel subsidy} \\
\hline
\end{tabularx}
\end{center}
\end{aretenir}

\begin{savaistu}[Le consommateur au Cameroun]
Au Cameroun, l'\eng{Association pour la Défense des Intérêts des Consommateurs (ADIC)} et le \eng{Ministère du Commerce} luttent contre la vie chère (\eng{high cost of living}), les produits frelatés (\eng{counterfeit/substandard goods}) et les poids falsifiés. La \eng{Loi n°2011/012} du 6 mai 2011 organise la protection des consommateurs. Elle impose l'affichage des prix (\eng{price display}), la délivrance du ticket de caisse (\eng{receipt}) et interdit la vente au-dessus des prix homologués (\eng{price gouging/overcharging}). Les marchés comme \eng{Mokolo} (Yaoundé), \eng{Mfoundi} ou \eng{Central Market} (Douala) sont des lieux où le \eng{haggling} (marchandage) est la norme, contrairement aux supermarchés (\eng{supermarkets}) où les prix sont fixes (\eng{fixed prices}).
\end{savaistu}

\courssub{Word families, false friends and pronunciation}

\begin{propriete}[Familles de mots (Word Formation) — Indispensable pour le Bac]
\begin{center}
\begin{tabularx}{\linewidth}{|X|X|X|X|}
\hline
\rowcolor{popblueL} \textbf{Verbe} & \textbf{Nom (chose/abstract)} & \textbf{Nom (agent)} & \textbf{Adjectif} \\
\hline
\eng{consume} & \eng{consumption} & \eng{consumer} & \eng{consumer} (adj) \\
\eng{economize} / \eng{economise} & \eng{economy} / \eng{economics} & \eng{economist} & \eng{economic} / \eng{economical} \\
\eng{afford} & \eng{affordability} & — & \eng{affordable} \\
\eng{protect} & \eng{protection} & \eng{protector} & \eng{protective} \\
\eng{compete} & \eng{competition} & \eng{competitor} & \eng{competitive} \\
\eng{guarantee} & \eng{guarantee} / \eng{warranty} & \eng{guarantor} & \eng{guaranteed} \\
\eng{refund} & \eng{refund} & — & \eng{refundable} \\
\eng{discount} & \eng{discount} & — & \eng{discounted} \\
\hline
\end{tabularx}
\end{center}
\emph{Note :} \eng{Economic} = relatif à l'économie (macro) ; \eng{Economical} = économe, bon marché. \eng{Warranty} (garantie écrite, légale) vs \eng{Guarantee} (promesse, engagement moral ou commercial).
\end{propriete}

\begin{attention}[Prononciation : Pièges sonores]
\begin{itemize}
    \item \eng{Receipt} /rɪˈsiːt/ → « ri-sit » (p muet).
    \item \eng{Warranty} /ˈwɒrənti/ → « wo-ren-ti » (pas « wa-ran-ti »).
    \item \eng{Guarantee} /ˌɡærənˈtiː/ → « ga-ran-ti ».
    \item \eng{Voucher} /ˈvaʊtʃə(r)/ → « vou-tcheu » (bon d'achat).
    \item \eng{Budget} /ˈbʌdʒɪt/ → « ba-djit ».
    \item \eng{Quality} /ˈkwɒləti/ → « quo-lo-ti ».
    \item \eng{Counterfeit} /ˈkaʊntəfɪt/ → « coun-teu-fit ».
\end{itemize}
\end{attention}

\courssub{Practice \& Bac Preparation}

\courssub{Je vérifie mes connaissances}

\begin{exercice}[QCM : Vocabulaire des prix et de la qualité]
Chaque phrase ne comporte qu'une seule réponse correcte. Recopiez la phrase en soulignant le bon choix.
\begin{enumerate}
    \item The \eng{receipt / recipe / prescription} proves you have paid for the goods.
    \item This shirt is a real \eng{bargain / debt / loan}; it costs only 2\,000 FCFA instead of 5\,000 FCFA.
    \item The \eng{consumer / customer / seller} is the person who buys goods or services for personal use.
    \item If a product is \eng{overpriced / affordable / cheap}, it costs more than it should.
    \item The shop offers a 10\% \eng{discount / increase / tax} for students with a valid ID card.
    \item You should always check the \eng{warranty / warning / guarantee} period before buying electronic devices.
    \item The \eng{fare / fee / price} for the bus from Yaoundé to Douala has increased this year.
    \item A \eng{refund / receipt / discount} is money returned to you when you take goods back to the shop.
\end{enumerate}
\end{exercice}


\begin{exercice}[Vrai ou Faux justifié : Comportement du consommateur]
Lisez les affirmations suivantes. Écrivez \textbf{True} ou \textbf{False} et justifiez votre réponse en anglais en une phrase, en vous basant sur le vocabulaire de la leçon.
\begin{enumerate}
    \item To \eng{haggle} means to accept the first price the seller asks for.
    \item \eng{Value for money} means the product is the cheapest on the market.
    \item A \eng{counterfeit} product is an authentic, original item.
    \item \eng{Window shopping} means buying things for your windows.
    \item An \eng{impulse buy} is a purchase planned long in advance.
    \item \eng{Amenities} are facilities like parks, libraries, or public transport provided for the community.
\end{enumerate}
\end{exercice}


\begin{exercice}[Vocabulaire : Champs lexicaux et collocations]
Complétez les phrases suivantes avec le mot ou l'expression approprié(e) parmi la liste ci-dessous. Chaque mot s'utilise une seule fois.
\begin{center}
\begin{tabular}{|l|l|l|l|}
\hline
\eng{affordable} & \eng{quality} & \eng{receipt} & \eng{bargain} \\
\hline
\eng{consumer rights} & \eng{overpriced} & \eng{warranty} & \eng{discount} \\
\hline
\end{tabular}
\end{center}
\begin{enumerate}
    \item Always keep your \eng{\_\_\_\_\_\_\_\_} in case you need to return the item.
    \item This phone is too \eng{\_\_\_\_\_\_\_\_}; I can find the same model cheaper elsewhere.
    \item The association defends \eng{\_\_\_\_\_\_\_\_} against dishonest traders.
    \item We got a 15\% \eng{\_\_\_\_\_\_\_\_} because we bought in bulk.
    \item Is this product \eng{\_\_\_\_\_\_\_\_} for a student budget?
    \item The \eng{\_\_\_\_\_\_\_\_} of this fabric is excellent; it will last years.
    \item The \eng{\_\_\_\_\_\_\_\_} covers repairs for two years.
    \item At the Marché Central, you can often find a good \eng{\_\_\_\_\_\_\_\_} if you know how to negotiate.
\end{enumerate}
\end{exercice}


\begin{exercice}[Grammaire en contexte : Formes verbales et prépositions]
Mettez le verbe entre parenthèses à la forme correcte (infinitif, gérondif, participe, temps conjugué) et ajoutez la préposition nécessaire le cas échéant.
\begin{enumerate}
    \item I am used to \eng{(compare)} prices before \eng{(buy)} expensive items.
    \item The seller refused \eng{(give)} me a refund without the receipt.
    \item She insists on \eng{(pay)} cash; she doesn't trust credit cards.
    \item This car is worth \eng{(buy)}; it is in perfect condition.
    \item They are thinking of \eng{(switch)} to a cheaper internet provider.
    \item The government aims \eng{(protect)} consumers from unfair practices.
    \item He avoided \eng{(shop)} at that store after his bad experience.
    \item We look forward to \eng{(see)} the new community centre open next month.
\end{enumerate}
\end{exercice}


\courssub{Je m'entraîne}

\begin{exercice}[Traduction : Du français vers l'anglais]
Traduisez les phrases suivantes en anglais en utilisant le vocabulaire précis de la leçon (prix, qualité, achat, services publics).
\begin{enumerate}
    \item Ce téléphone est trop cher ; je ne peux pas me le permettre.
    \item Il a négocié le prix pendant dix minutes et a obtenu une réduction.
    \item Ce produit est une contrefaçon ; la qualité est très mauvaise.
    \item Les services publics comme l'eau et l'électricité sont des ressources communautaires essentielles.
    \item Nous comparons toujours les prix pour trouver le meilleur rapport qualité-prix.
    \item Avez-vous une garantie pour cet ordinateur portable ?
    \item Le ticket de caisse est indispensable pour un remboursement.
    \item Ce supermarché propose des produits abordables pour les familles à faible revenu.
\end{enumerate}
\end{exercice}


\begin{exercice}[Transformation : Reformulation et registres]
Réécrivez chaque phrase en commençant par les mots indiqués, sans en changer le sens.
\begin{enumerate}
    \item The price of fuel has gone up again.
    \begin{quote} There has been \eng{\_\_\_\_\_\_\_\_}
    \end{quote}
    \item "How much does this cost?" she asked the shopkeeper.
    \begin{quote} She asked the shopkeeper \eng{\_\_\_\_\_\_\_\_}
    \end{quote}
    \item It is not worth repairing that old fridge.
    \begin{quote} That old fridge is not \eng{\_\_\_\_\_\_\_\_}
    \end{quote}
    \item The market vendors sell fresh produce at low prices.
    \begin{quote} Fresh produce \eng{\_\_\_\_\_\_\_\_}
    \end{quote}
    \item You should check the expiry date before buying.
    \begin{quote} Don't forget \eng{\_\_\_\_\_\_\_\_}
    \end{quote}
    \item The community centre offers free Wi-Fi to residents.
    \begin{quote} Free Wi-Fi \eng{\_\_\_\_\_\_\_\_}
    \end{quote}
    \item "I will give you a discount if you buy two," said the seller.
    \begin{quote} The seller promised \eng{\_\_\_\_\_\_\_\_}
    \end{quote}
    \item Public transport in Buea is more reliable than in Douala.
    \begin{quote} Public transport in Douala is not \eng{\_\_\_\_\_\_\_\_}
    \end{quote}
\end{enumerate}
\end{exercice}


\begin{exercice}[Texte à trous : La vie chère à Yaoundé]
Complétez le texte suivant avec les mots de la liste. Attention : il y a \textbf{trois} mots en trop. Mettez le verbe à la forme appropriée si nécessaire.
\begin{center}
\begin{tabular}{|l|l|l|l|l|}
\hline
\eng{afford} & \eng{bargain} & \eng{consumer} & \eng{discount} & \eng{fare} \\
\hline
\eng{fee} & \eng{overpriced} & \eng{quality} & \eng{receipt} & \eng{refund} \\
\hline
\eng{warranty} & \eng{value for money} & \eng{haggle} & \eng{amenities} & \eng{counterfeit} \\
\hline
\end{tabular}
\end{center}

\begin{textebox}[Smart Shopping in Yaoundé]
With the rising cost of living in Yaoundé, being a smart \eng{\_\_\_\_\_\_\_\_} (1) is essential. Many people struggle to \eng{\_\_\_\_\_\_\_\_} (2) basic necessities. At the Marché Mokolo, customers often \eng{\_\_\_\_\_\_\_\_} (3) to get a better price. A good strategy is to check the \eng{\_\_\_\_\_\_\_\_} (4) of the product before buying; cheap items that break quickly are not good \eng{\_\_\_\_\_\_\_\_} (5). Be careful of \eng{\_\_\_\_\_\_\_\_} (6) goods, especially electronics, which have no \eng{\_\_\_\_\_\_\_\_} (7). Always ask for a \eng{\_\_\_\_\_\_\_\_} (8) to get a \eng{\_\_\_\_\_\_\_\_} (9) if the item is faulty. Supermarkets sometimes offer a \eng{\_\_\_\_\_\_\_\_} (10) for loyalty card holders. Finally, using public transport like the \eng{bus fare} (11) is cheaper than taking a taxi. Access to good community \eng{\_\_\_\_\_\_\_\_} (12) like markets and transport helps the local economy.
\end{textebox}
\end{exercice}


\begin{exercice}[Appariement : Mots, définitions et contexte]
Associez chaque mot anglais (1--10) à sa définition française (A--J) et donnez un exemple de collocation (mot qui va naturellement avec) en anglais.
\begin{center}
\begin{tabular}{|c|l|c|l|}
\hline
\textbf{Mot} & \textbf{Anglais} & \textbf{Def.} & \textbf{Français} \\
\hline
1 & \eng{Receipt} & A & Marchander, négocier le prix \\
2 & \eng{Haggle} & B & Contrefaçon \\
3 & \eng{Counterfeit} & C & Ticket de caisse \\
4 & \eng{Warranty} & D & Prix du transport (bus, train) \\
5 & \eng{Fare} & E & Garantie (constructeur) \\
6 & \eng{Amenities} & F & Accessible, à prix raisonnable \\
7 & \eng{Affordable} & G & Équipements collectifs, commodités \\
8 & \eng{Overpriced} & H & Remboursement \\
9 & \eng{Refund} & I & Rapport qualité-prix \\
10 & \eng{Value for money} & J & Trop cher, surévalué \\
\hline
\end{tabular}
\end{center}
\textit{Présentez vos réponses sous forme de tableau : Numéro -- Lettre -- Collocation anglaise (ex: \eng{keep a receipt}).}
\end{exercice}


\courssub{Je me prépare au Bac}

\begin{exercice}[Compréhension et expression écrites : Texte original type Bac]
Lisez attentivement le texte ci-dessous et répondez aux questions qui suivent en anglais.

\begin{textebox}[Community Resources vs. Private Services: The Case of Buea Municipality]
In Buea, the capital of the Southwest Region, residents face a daily dilemma: should they rely on community resources or opt for private services? The municipal council provides essential \eng{amenities} such as public water points, community health centres, and the main market. These services are funded by local \eng{taxes} and are supposed to be \eng{affordable} for everyone. However, frequent water cuts and long queues at the public clinic push many families towards private alternatives.

Private boreholes sell water at 100 FCFA per bucket, while the public tap is free but often dry. Private clinics offer faster service but at a high \eng{fee}. A consultation can cost 5\,000 FCFA, excluding medication. "It is a matter of \eng{value for money}," says Mr. Nfor, a teacher. "If I lose a day's work waiting at the public clinic, I lose more money than the private consultation fee."

The same comparison applies to transport. The public \eng{fare} for the "Car Rapid" is fixed by the government, but drivers often \eng{overcharge} passengers during rush hours. Private "Okada" (motorcycle taxis) are faster but riskier and more expensive. Consumers must constantly \eng{compare} costs, \eng{quality}, and reliability to make the \eng{best buy} for their time and money. The municipality has promised to improve \eng{infrastructure} and enforce \eng{consumer protection} laws, but progress is slow. Until then, the "smart consumer" in Buea is the one who knows where to find the right balance between price and service.
\end{textebox}

\textbf{Partie A : Comprehension Questions (10 marks)}
Answer the following questions in your own words as far as possible. Use complete sentences.
\begin{enumerate}
    \item What are three examples of community resources provided by the Buea municipal council mentioned in the text?
    \item Why do some residents prefer private clinics despite the high fees?
    \item Explain the expression "value for money" as used by Mr. Nfor in the context of the text.
    \item What problem do passengers face with public "Car Rapid" transport during rush hours?
    \item According to the last paragraph, what two things has the municipality promised to do?
\end{enumerate}

\textbf{Partie B : Language in Use (10 marks)}
\begin{enumerate}
    \item Find in the text words or expressions that mean the same as:
    \begin{itemize}
        \item facilities / useful services (paragraph 1)
        \item reasonable price / not expensive (paragraph 1)
        \item charge too much (paragraph 3)
        \item apply / make sure laws are respected (paragraph 4)
    \end{itemize}
    \item Rewrite the following sentences using the prompts given:
    \begin{enumerate}
        \item "The public tap is free but often dry," the journalist reported.
        \begin{quote} The journalist reported that \eng{\_\_\_\_\_\_\_\_}
        \end{quote}
        \item Private clinics are more expensive than public ones.
        \begin{quote} Public clinics are not \eng{\_\_\_\_\_\_\_\_}
        \end{quote}
        \item Residents must compare prices. It is necessary.
        \begin{quote} Residents \eng{\_\_\_\_\_\_\_\_}
        \end{quote}
        \item The municipality will improve infrastructure. This will help consumers.
        \begin{quote} If the municipality \eng{\_\_\_\_\_\_\_\_}
        \end{quote}
    \end{enumerate}
    \item Complete the sentences with the correct form of the words in brackets:
    \begin{enumerate}
        \item The \eng{\_\_\_\_\_\_\_\_} (consume) association held a meeting yesterday.
        \item We need an \eng{\_\_\_\_\_\_\_\_} (economy) solution to the water crisis.
        \item This product offers great \eng{\_\_\_\_\_\_\_\_} (value) for money.
        \item Is this housing \eng{\_\_\_\_\_\_\_\_} (afford) for young couples?
    \end{enumerate}
\end{enumerate}
\end{exercice}


\begin{exercice}[Traduction : Thème et version (10 marks)]
\textbf{A. Version (English to French) :} Traduisez en français le passage suivant extrait du texte ci-dessus (depuis "Private boreholes..." jusqu'à "...private consultation fee.").
\begin{quote}
"Private boreholes sell water at 100 FCFA per bucket, while the public tap is free but often dry. Private clinics offer faster service but at a high fee. A consultation can cost 5,000 FCFA, excluding medication. 'It is a matter of value for money,' says Mr. Nfor, a teacher. 'If I lose a day's work waiting at the public clinic, I lose more money than the private consultation fee.'"
\end{quote}

\textbf{B. Thème (French to English) :} Traduisez en anglais.
\begin{quote}
"Au marché de Mfoundi, les vendeurs fixent souvent des prix élevés pour les touristes. Un acheteur averti doit savoir marchander pour obtenir une réduction. La semaine dernière, j'ai acheté un tissu de bonne qualité à un prix abordable parce que j'ai comparé les offres de trois boutiques différentes. Cependant, il faut faire attention aux contrefaçons qui n'ont aucune garantie. Si un produit est défectueux, le consommateur a droit à un remboursement sur présentation du ticket de caisse."
\end{quote}
\end{exercice}


\begin{exercice}[Production écrite : Composition guidée (10 marks)]
\textbf{Sujet :} \eng{Compare two products or services available in your town (e.g., two phone networks, two transport options, two food shops, a public vs. private clinic) and explain which one is the "best buy" for you.}

\textbf{Consignes :}
\begin{itemize}
    \item Write between \textbf{80 and 100 words}.
    \item Structure your text: Introduction (present the two options), Comparison (price, quality, reliability, accessibility), Conclusion (your choice with reasons).
    \item Use at least \textbf{8 vocabulary items} from this lesson (e.g., \eng{affordable, overpriced, bargain, discount, warranty, value for money, consumer, amenities, fare, fee, refund, quality, haggle, counterfeit}).
    \item Underline the vocabulary items used.
\end{itemize}

\textbf{Grille d'évaluation indicative :}
\begin{center}
\begin{tabular}{|l|c|}
\hline
\textbf{Critères} & \textbf{Points} \\
\hline
Respect du nombre de mots et structure (Intro/Comparaison/Conclusion) & 2 \\
Richesse et exactitude du vocabulaire thématique (8 mots soulignés minimum) & 4 \\
Exactitude grammaticale (temps, prépositions, ordre des mots) & 4 \\
Cohérence, liaison logique et pertinence de l'argumentation & 2 \\
\hline
\textbf{Total} & \textbf{10} \\
\hline
\end{tabular}
\end{center}
\end{exercice}

\begin{exemplebox}[Modèle de composition (Exemple de réponse attendue)]
In Douala, commuters often choose between \eng{public transport} (Car Rapid) and \eng{Okada} (motorcycle taxis). The \eng{bus fare} is fixed and \eng{affordable} (150 FCFA), but buses are overcrowded, unreliable, and drivers sometimes \eng{overcharge}. \eng{Okada} are faster and available everywhere, but the \eng{fare} is higher (300--500 FCFA), riskier, and offers no \eng{warranty} for safety. As a \eng{consumer}, I \eng{compare} \eng{quality}, cost, and time. For daily commuting, the bus offers better \eng{value for money} despite delays. However, for urgent trips, an \eng{Okada} is the best \eng{bargain} for my time. \textbf{(92 words)}
\end{exemplebox}

\newpage

\coursec{Corrigés des exercices}

\begin{corrige}[Exercice 1 — QCM : Vocabulaire des prix et de la qualité]
\begin{enumerate}
    \item \eng{receipt} (ticket de caisse, reçu). \eng{Recipe} = recette de cuisine ; \eng{prescription} = ordonnance médicale. Faux ami classique à éliminer définitivement : un \eng{receipt} n'est ni une recette ni une ordonnance.
    \item \eng{bargain} (bonne affaire). \eng{Debt} = dette ; \eng{loan} = prêt (bancaire).
    \item \eng{consumer} (consommateur final, celui qui utilise le bien ou le service). \eng{Customer} = client (relation commerciale avec un vendeur) ; \eng{seller} = vendeur.
    \item \eng{overpriced} (trop cher par rapport à sa valeur réelle). \eng{Affordable} = abordable ; \eng{cheap} = bon marché (parfois péjoratif : de mauvaise qualité).
    \item \eng{discount} (réduction, remise). \eng{Increase} = augmentation ; \eng{tax} = taxe, impôt.
    \item \eng{warranty} (garantie constructeur, écrite, pour les appareils et produits manufacturés). \eng{Warning} = avertissement. \eng{Guarantee} serait compréhensible, mais \eng{warranty} est le terme technique exact pour un appareil.
    \item \eng{fare} (tarif de transport : bus, taxi, avion). \eng{Fee} = frais de service, honoraires, droits (scolaires, d'inscription) ; \eng{price} = prix d'un bien.
    \item \eng{refund} (remboursement). \eng{Receipt} = reçu ; \eng{discount} = réduction.
\end{enumerate}
\textbf{Points de vigilance :} relire la boîte \emph{Faux amis} sur le trio \eng{receipt / recipe / prescription} et le trio \eng{fare / fee / price}, très fréquents au Bac. Retenir aussi la prononciation : \eng{receipt} se prononce « ri-site » (le \eng{p} est muet).
\end{corrige}

\begin{corrige}[Exercice 2 — Vrai ou Faux justifié]
\begin{enumerate}
    \item \textbf{False.} To \eng{haggle} means to negotiate the price with the seller, not to accept the first offer.
    \item \textbf{False.} \eng{Value for money} means good quality for the price paid, not necessarily the lowest price.
    \item \textbf{False.} A \eng{counterfeit} product is a fake imitation of a genuine brand, not an authentic item.
    \item \textbf{False.} \eng{Window shopping} means looking at goods in shop windows without buying anything.
    \item \textbf{False.} An \eng{impulse buy} is an unplanned, spontaneous purchase, not a planned one.
    \item \textbf{True.} \eng{Amenities} are public facilities such as parks, libraries or transport provided for the community.
\end{enumerate}
\textbf{Points de vigilance :} la justification doit être une phrase complète en anglais, avec sujet et verbe ; éviter les réponses d'un seul mot au Bac, qui ne rapportent pas le point de langue.
\end{corrige}

\begin{corrige}[Exercice 3 — Champs lexicaux et collocations]
\begin{enumerate}
    \item \eng{receipt} — collocation : \eng{keep your receipt}.
    \item \eng{overpriced} — adjectif précédé de \eng{too} : \eng{too overpriced}.
    \item \eng{consumer rights} — expression fixe, toujours au pluriel.
    \item \eng{discount} — \eng{a 15\% discount}.
    \item \eng{affordable} — \eng{affordable for a student budget}.
    \item \eng{quality} — \eng{the quality of this fabric}.
    \item \eng{warranty} — \eng{the warranty covers repairs}.
    \item \eng{bargain} — \eng{find a good bargain}.
\end{enumerate}
\textbf{Points de vigilance :} chaque mot ne sert qu'une fois ; vérifier l'accord article + nom (\eng{a discount}, \eng{a bargain}) et la préposition qui suit l'adjectif (\eng{affordable for}, \eng{the quality of}).
\end{corrige}

\begin{corrige}[Exercice 4 — Formes verbales et prépositions]
\begin{enumerate}
    \item \eng{comparing} / \eng{buying} — \eng{be used to} + V-ing (habitude acquise) ; \eng{before} + V-ing.
    \item \eng{to give} — \eng{refuse} + infinitif en \eng{to}.
    \item \eng{paying} — \eng{insist on} + V-ing.
    \item \eng{buying} — \eng{be worth} + V-ing (jamais \eng{worth to buy}).
    \item \eng{switching} — \eng{think of} + V-ing.
    \item \eng{to protect} — \eng{aim} + infinitif en \eng{to}.
    \item \eng{shopping} — \eng{avoid} + V-ing.
    \item \eng{seeing} — \eng{look forward to} + V-ing (ici, \eng{to} est une préposition, pas une marque d'infinitif).
\end{enumerate}
\textbf{Points de vigilance :} les deux pièges majeurs sont \eng{look forward to seeing} et \eng{worth buying} ; les francophones écrivent souvent l'infinitif par calque du français (« j'ai hâte de voir », « ça vaut la peine d'acheter »). Après toute préposition (\eng{before, of, on, without, by}), le verbe anglais prend la forme en \eng{-ing}.
\end{corrige}

\begin{corrige}[Exercice 5 — Traduction : du français vers l'anglais]
\begin{enumerate}
    \item \eng{This phone is too expensive; I can't afford it.} (Variante : \eng{It is overpriced; I can't afford it.}) — \eng{afford} s'emploie presque toujours avec un modal (\eng{can / can't afford}).
    \item \eng{He haggled over the price for ten minutes and got a discount.}
    \item \eng{This product is counterfeit; the quality is very poor.} (Variante : \eng{This product is a fake.})
    \item \eng{Public utilities like water and electricity are essential community resources.}
    \item \eng{We always compare prices to find the best value for money.}
    \item \eng{Do you have a warranty for this laptop?}
    \item \eng{The receipt is essential for a refund.} (Variante : \eng{You need the receipt to get a refund.})
    \item \eng{This supermarket offers affordable products for low-income families.}
\end{enumerate}
\textbf{Points de vigilance :} ne pas traduire « se permettre (financièrement) » par \eng{permit} — c'est \eng{afford} ; « rapport qualité-prix » = \eng{value for money} (expression invariable, sans article) ; « contrefaçon » = \eng{counterfeit} ou \eng{fake}, jamais \eng{falsification}.
\end{corrige}

\begin{corrige}[Exercice 6 — Transformation : reformulation]
\begin{enumerate}
    \item \eng{There has been an increase in the price of fuel.} (Variante : \eng{There has been a rise in fuel prices.})
    \item \eng{She asked the shopkeeper how much it cost.} — discours rapporté : recul du temps (\eng{does it cost} $\rightarrow$ \eng{it cost}), suppression de l'inversion et de l'auxiliaire \eng{do}.
    \item \eng{That old fridge is not worth repairing.} — \eng{worth} + V-ing.
    \item \eng{Fresh produce is sold at low prices by market vendors.} — voix passive ; \eng{produce} (produits agricoles) est ici singulier indénombrable.
    \item \eng{Don't forget to check the expiry date before buying.} — \eng{forget to} + infinitif (action future à accomplir).
    \item \eng{Free Wi-Fi is offered to residents by the community centre.} — passif ; structure \eng{offer something to somebody}.
    \item \eng{The seller promised to give me a discount if I bought two.} — \eng{promise} + infinitif ; recul du temps dans la subordonnée (\eng{buy} $\rightarrow$ \eng{bought}).
    \item \eng{Public transport in Douala is not as reliable as in Buea.} — comparatif d'égalité à la forme négative (\eng{not as ... as}).
\end{enumerate}
\textbf{Points de vigilance :} en discours rapporté, toujours vérifier le backshift et l'ordre sujet--verbe dans la question indirecte (\eng{how much it cost}, jamais \eng{how much did it cost}). Attention aussi à \eng{forget to do} (ne pas oublier de faire) contre \eng{forget doing} (avoir oublié d'avoir fait).
\end{corrige}

\begin{corrige}[Exercice 7 — Texte à trous : Smart Shopping in Yaoundé]
\begin{enumerate}
    \item \eng{consumer}
    \item \eng{afford}
    \item \eng{haggle}
    \item \eng{quality}
    \item \eng{value for money}
    \item \eng{counterfeit}
    \item \eng{warranty}
    \item \eng{receipt}
    \item \eng{refund}
    \item \eng{discount}
    \item \eng{fare} (\eng{bus fare} : tarif du bus)
    \item \eng{amenities}
\end{enumerate}
\textbf{Mots en trop (non utilisés) :} \eng{fee}, \eng{overpriced}, \eng{bargain}.

\textbf{Points de vigilance :} au trou 5, \eng{value for money} est invariable ; au trou 6, \eng{counterfeit} est ici adjectif devant \eng{goods} ; au trou 11, le mot \eng{bus} figurant déjà dans le texte, seul \eng{fare} convient (\eng{fee} désignerait des frais de service, pas un tarif de transport).
\end{corrige}

\begin{corrige}[Exercice 8 — Appariement : mots, définitions, collocations]
\begin{center}
\begin{tabular}{|c|c|l|}
\hline
\rowcolor{popblueL} \textbf{N°} & \textbf{Lettre} & \textbf{Collocation possible} \\
\hline
1 & C & \eng{keep a receipt} / \eng{ask for a receipt} \\
2 & A & \eng{haggle over the price} / \eng{haggle with the seller} \\
3 & B & \eng{counterfeit goods} / \eng{spot a counterfeit} \\
4 & E & \eng{a two-year warranty} / \eng{under warranty} \\
5 & D & \eng{bus fare} / \eng{pay the fare} \\
6 & G & \eng{local amenities} / \eng{public amenities} \\
7 & F & \eng{affordable housing} / \eng{affordable prices} \\
8 & J & \eng{slightly overpriced} / \eng{grossly overpriced} \\
9 & H & \eng{get a refund} / \eng{a full refund} \\
10 & I & \eng{good value for money} / \eng{best value for money} \\
\hline
\end{tabular}
\end{center}
\textbf{Points de vigilance :} une collocation doit être naturelle en anglais ; \eng{make a discount} est un calque du français — on dit \eng{give a discount} ou \eng{offer a discount}. De même, on dit \eng{pay the fare}, jamais \eng{pay the price of the bus}.
\end{corrige}

\begin{corrige}[Exercice 9 — Bac Prep : compréhension et langue en contexte]
\textbf{Partie A — Comprehension Questions (10 points, 2 pts par question : 1 pt sens, 1 pt langue)}
\begin{enumerate}
    \item \eng{The council provides public water points, community health centres and the main market.}
    \item \eng{They prefer private clinics because the public ones have long queues and frequent water cuts, so waiting there makes them lose time and money.}
    \item \eng{It means getting a service whose quality justifies its price: paying the private fee costs him less than losing a day's work at the public clinic.}
    \item \eng{During rush hours, drivers often overcharge passengers instead of respecting the fixed government fare.}
    \item \eng{The municipality has promised to improve infrastructure and to enforce consumer protection laws.}
\end{enumerate}

\textbf{Partie B — Language in Use (10 points)}
\begin{enumerate}
    \item Synonymes dans le texte (1 pt) : \eng{amenities} ; \eng{affordable} ; \eng{overcharge} ; \eng{enforce}.
    \item Reformulation (2 pts) :
    \begin{enumerate}
        \item \eng{The journalist reported that the public tap was free but often dry.} (backshift : \eng{is} $\rightarrow$ \eng{was}).
        \item \eng{Public clinics are not as expensive as private ones.}
        \item \eng{Residents have to compare prices.} (ou \eng{need to} / \eng{must}).
        \item \eng{If the municipality improves infrastructure, it will help consumers.} (premier conditionnel : présent + \eng{will}).
    \end{enumerate}
    \item Formation de mots (2 pts) :
    \begin{enumerate}
        \item \eng{consumer} (\eng{consumer association}, expression fixe).
        \item \eng{economical} (économe, qui fait économiser — pas \eng{economic}).
        \item \eng{value} (\eng{value for money}, invariable).
        \item \eng{affordable} (suffixe \eng{-able} sur le verbe \eng{afford}).
    \end{enumerate}
\end{enumerate}

\textbf{Points de vigilance :} répondre par des phrases complètes sans recopier intégralement le texte ; distinguer soigneusement \eng{economic} (qui relève de l'économie, de la macro-économie : \eng{economic growth}) et \eng{economical} (qui fait économiser : \eng{an economical car}) ; dans le conditionnel, jamais \eng{will} dans la proposition introduite par \eng{if}.
\end{corrige}

\begin{corrige}[Exercice 10 — Traduction : thème et version (10 points)]
\textbf{A. Version (5 points)}

« Les forages privés vendent l'eau à 100 FCFA le seau, alors que le robinet public est gratuit mais souvent à sec. Les cliniques privées offrent un service plus rapide, mais à des tarifs élevés. Une consultation peut coûter 5 000 FCFA, médicaments non compris. "C'est une question de rapport qualité-prix", explique M. Nfor, enseignant. "Si je perds une journée de travail à attendre au centre de santé public, je perds plus d'argent que le prix de la consultation privée." »

\textbf{B. Thème (5 points)}

\eng{At Mfoundi market, vendors often set high prices for tourists. A smart buyer must know how to haggle to get a discount. Last week, I bought a good quality fabric at an affordable price because I compared the offers of three different shops. However, one must beware of counterfeit goods, which have no warranty. If a product is faulty, the consumer is entitled to a refund on presentation of the receipt.}

\textbf{Points de vigilance :} « médicaments non compris » = \eng{excluding medication} ; « avoir droit à » = \eng{be entitled to} ; « se méfier de » = \eng{beware of} (sans \eng{to} après) ; « acheteur averti » = \eng{smart buyer} ou \eng{savvy shopper} ; accorder le verbe du relatif avec son antécédent : \eng{goods, which have}.
\end{corrige}

\begin{corrige}[Exercice 11 — Production écrite : composition guidée (10 points)]
\textbf{Barème indicatif :}
\begin{center}
\begin{tabular}{|l|c|}
\hline
\rowcolor{popblueL} \textbf{Critères} & \textbf{Points} \\
\hline
Respect du nombre de mots (80--100) et structure en trois parties & 2 \\
Vocabulaire thématique : au moins 8 items soulignés, employés correctement & 4 \\
Exactitude grammaticale (temps, prépositions, comparatifs) & 2 \\
Cohérence, connecteurs logiques, pertinence de l'argumentation & 2 \\
\hline
\textbf{Total} & \textbf{10} \\
\hline
\end{tabular}
\end{center}

\textbf{Proposition de rédaction modèle :}

\eng{In Bamenda, students can choose between two internet providers: Camtel and MTN. Camtel offers an affordable monthly fee and unlimited data, which is excellent value for money, but the connection is slow and often cut. MTN is faster and more reliable, but its bundles are slightly overpriced for a student budget, and there is no refund when the network fails. As a careful consumer, I always compare quality, price and reliability before paying. For my research work, MTN is the best buy because speed matters more than savings; but for downloading films, Camtel is a real bargain. The best choice depends on your needs.} \textbf{(97 mots)}

\textbf{Points de vigilance :}
\begin{itemize}
    \item Souligner les 8 items lexicaux exigés (\eng{affordable, fee, value for money, overpriced, refund, consumer, compare, bargain}).
    \item Employer au moins un comparatif (\eng{faster than}, \eng{more reliable than}) et un connecteur de contraste (\eng{but, however, whereas}).
    \item Éviter les calques : « le meilleur achat » = \eng{the best buy} ; « être coupé » (réseau, connexion) = \eng{be cut} ou \eng{go down}, pas un calque du français.
    \item Compter les mots : toute copie hors de la fourchette 80--100 perd le point de respect des consignes.
\end{itemize}
\end{corrige}

\newpage

\coursec{Fiche bilan}

\begin{popfigure}
\begin{center}
\begin{tikzpicture}[mindmap, grow cyclic,
  every node/.style={concept, text=white, align=center, font=\small},
  root concept/.append style={concept color=popblue, font=\bfseries\small},
  level 1/.append style={level distance=3.1cm, sibling angle=60},
  level 2/.append style={level distance=2.1cm, sibling angle=45, font=\scriptsize}]
\node[root concept, align=center]{Best Buys\\ Consumer\\ Economy}
  child[concept color=poporange]{ node {Prices \& Money}
    child { node[align=center]{cheap, costly\\ discount, receipt} }
    child { node[align=center]{afford, charge\\ save, spend} }
  }
  child[concept color=poppurple]{ node {Comparing Quality}
    child { node[align=center]{good value\\ for money} }
    child { node[align=center]{durable, reliable\\ worth buying} }
  }
  child[concept color=popgreen]{ node {Shopping \& Bargaining}
    child { node[align=center]{bargain, haggle\\ reduce the price} }
    child { node[align=center]{shop around\\ compare prices} }
  }
  child[concept color=poppink]{ node {Community Services}
    child { node[align=center]{clinic, library\\ water supply} }
    child { node[align=center]{public transport\\ town hall} }
  }
  child[concept color=popgold]{ node {Word Families}
    child { node[align=center]{economy\\ economic\\ economical} }
    child { node[align=center]{value\\ valuable\\ invaluable} }
  }
  child[concept color=popdark]{ node {False Friends}
    child { node[align=center]{actually\\ = en fait} }
    child { node[align=center]{sensible\\ = raisonnable} }
  };
\end{tikzpicture}
\end{center}
\legende{Carte mentale : le vocabulaire des meilleurs achats et des services communautaires}
\end{popfigure}

\begin{bilan}
\textbf{1. Lexique clé : parler des prix et de l'argent}
\begin{itemize}
  \item \eng{price} (n.m.) : le prix ; \eng{cheap / expensive} : bon marché / cher ; \eng{costly} : coûteux.
  \item \eng{to cost -- cost -- cost} : coûter ; \eng{to charge} : facturer, demander (un prix) ; \eng{to afford} : avoir les moyens de.
  \item \eng{a discount} : une remise ; \eng{a receipt} : un reçu ; \eng{change} (n.) : la monnaie (à rendre).
  \item \eng{to save money} : économiser ; \eng{to spend money on} : dépenser pour ; \eng{to waste money} : gaspiller.
\end{itemize}

\textbf{2. Lexique clé : comparer la qualité et trouver la meilleure affaire}
\begin{itemize}
  \item \eng{good value for money} : d'un bon rapport qualité-prix ; \eng{worth buying} : qui vaut la peine d'être acheté.
  \item \eng{durable / hard-wearing} : durable / résistant ; \eng{reliable} : fiable ; \eng{low-quality / second-rate} : de mauvaise qualité.
  \item \eng{the best buy} : le meilleur achat ; \eng{a bargain} : une bonne affaire ; \eng{a rip-off} : une arnaque.
  \item \eng{to shop around} : comparer les prix avant d'acheter ; \eng{to compare prices} : comparer les prix.
\end{itemize}

\textbf{3. Lexique clé : marchander et services communautaires}
\begin{itemize}
  \item \eng{to bargain / to haggle} : marchander ; \eng{to reduce the price} : baisser le prix ; \eng{a fixed price} : un prix fixe.
  \item \eng{a consumer} : un consommateur ; \eng{consumer rights} : les droits du consommateur.
  \item \eng{community resources} : les ressources communautaires ; \eng{public services} : les services publics.
  \item \eng{a health centre} : un centre de santé ; \eng{water supply} : l'adduction d'eau ; \eng{public transport} : les transports en commun.
\end{itemize}

\textbf{4. Grammaire et structures à connaître par c\oe ur}
\begin{center}
\begin{tabularx}{\linewidth}{|X|X|}
\hline
\rowcolor{popblueL} \textbf{Structure} & \textbf{Exemple} \\
\hline
Comparatif : \eng{-er} / \eng{more ... than} & \eng{This phone is cheaper than that one.} \\
\hline
Superlatif : \eng{the -est} / \eng{the most} & \eng{It is the best buy in the market.} \\
\hline
Irréguliers : \eng{good / better / the best} ; \eng{bad / worse / the worst} & \eng{The second shop offers better quality.} \\
\hline
\eng{too + adjectif} / \eng{adjectif + enough} & \eng{It is too expensive. / It is cheap enough.} \\
\hline
\eng{to be worth + V-ing} & \eng{This bag is worth buying.} \\
\hline
\eng{to afford + nom / to + V} & \eng{We cannot afford to waste money.} \\
\hline
\end{tabularx}
\end{center}

\textbf{5. Faux amis à éviter}
\begin{itemize}
  \item \eng{actually} = en fait (et non « actuellement » $\rightarrow$ \eng{currently}).
  \item \eng{sensible} = raisonnable (et non « sensible » $\rightarrow$ \eng{sensitive}).
  \item \eng{economic} = économique (relatif à l'économie) ; \eng{economical} = économe, économique (qui coûte peu).
  \item \eng{to demand} = exiger (et non « demander » $\rightarrow$ \eng{to ask for}).
\end{itemize}

\textbf{6. Méthodes express}
\begin{itemize}
  \item Pour comparer deux produits : prix (\eng{cheaper than}), qualité (\eng{more durable}), conclusion (\eng{the best buy is ...}).
  \item Au marché : saluer, demander le prix (\eng{How much is it?}), marchander poliment (\eng{Can you reduce the price?}), conclure (\eng{I'll take it.}).
  \item Familles de mots : \eng{economy $\rightarrow$ economic $\rightarrow$ economical $\rightarrow$ economically} ; \eng{value $\rightarrow$ valuable $\rightarrow$ invaluable}.
\end{itemize}
\end{bilan}

\begin{aretenir}[Checklist avant le Bac]
\begin{itemize}
  \item $\square$ Je sais nommer les prix, les remises et les moyens de paiement en anglais.
  \item $\square$ Je maîtrise les comparatifs et superlatifs, y compris \eng{good / better / the best}.
  \item $\square$ J'emploie correctement \eng{good value for money} et \eng{the best buy}.
  \item $\square$ Je sais utiliser \eng{to be worth + V-ing} et \eng{too ... / ... enough}.
  \item $\square$ Je connais le vocabulaire du marchandage : \eng{to bargain, to haggle, to reduce the price}.
  \item $\square$ Je sais parler des services communautaires : \eng{health centre, water supply, public transport}.
  \item $\square$ Je distingue \eng{economic} et \eng{economical}.
  \item $\square$ J'évite les faux amis : \eng{actually}, \eng{sensible}, \eng{to demand}.
  \item $\square$ Je connais les familles de mots : \eng{economy, value, consume, compare}.
  \item $\square$ Je peux tenir un dialogue d'achat au marché et rédiger un paragraphe comparant deux produits.
  \item $\square$ Je prononce correctement \eng{cheap} (« tchipe »), \eng{receipt} (« ri-site », p muet) et \eng{bargain} (« bar-guine »).
  \item $\square$ Je sais conclure une comparaison avec une recommandation claire : \eng{In my opinion, the best buy is ...}.
\end{itemize}
\end{aretenir}

\findecours

\end{document}
